\documentclass[11pt]{ctexart}
\usepackage[a4paper,margin=1in]{geometry}
\usepackage{graphicx}
\usepackage{xcolor}
\usepackage{hyperref}
\definecolor{refgray}{gray}{0.45}
\title{\textbf{市场最新观点汇总}\\[0.4em]{\large 地缘风险溢价退潮与AI结构性扩张并行，全球资产定价逻辑正从宏观叙事切换至微观结构再定价}}
\author{KC桌面 自动生成}
\date{260616}
\begin{document}
\maketitle
\tableofcontents
\newpage
\section{一页摘要}
\begin{itemize}
  \item 市场对美联储降息的定价过于激进，而对财政供给压力和地缘供给侧冲击的持续性则可能低估。
  \item AI基础设施投资逻辑正从'GPU短缺'转向'系统级重构'，互联层与存储环节的定价权正在转移。
  \item 美伊和平协议触发全球资产重定价，能源价格、美元走势和风险偏好三个核心变量正在被重新校准。
  \item 中国权益市场的核心叙事从'周期性复苏'切换到'结构性供给优势'，出口韧性与内需疲弱的分化加剧。
  \item 亚洲央行正从'被迫加息'转向'主动加息'，市场对此定价不足，尤其是韩国、台湾、新加坡。
  \item 上游资本开支的'超级周期'叙事可能被证伪，资本纪律比油价更顽固，供给增量有限。
\end{itemize}
\section{宏观与利率：财政供给与通胀再平衡主导定价}
\textbf{全球利率曲线的定价锚点正在从政策利率转向财政供给与期限溢价，市场对美联储降息的预期过于激进，而对通胀粘性及财政压力的持续性估计不足。}

\begin{itemize}
  \item JPM指出，美国利率定价的锚点正从'政策利率'转向'财政供给'，由于2027财年存在巨大融资缺口，财政部预计将在8月移除前瞻指引中的'至少'措辞，为多季度息票增发做准备，这将系统性抬升长端利率。
  \item JPM预测2年期和10年期美债收益率将在2026年底分别升至4.20\%和4.70\%，而市场OIS定价未来1-2年加息35bp，比美联储自身预期更鹰。
  \item GS在6月FOMC前瞻中认为，加息的可能性被市场系统性高估，当前通胀脉冲主要由关税、油价和存储芯片价格等一次性供给冲击构成，而非需求驱动的持续性通胀，加息门槛远比市场想象的高。
  \item BofA的6月全球基金经理调查显示，认为美联储将在未来12个月内加息的基金经理占比从16\%飙升至40\%，为近四年最高，利率预期的剧烈转向正成为仓位调整的核心驱动力。
  \item GS的供应链拥堵综合指数已降至'2'，基本回到疫情前水平，这并非需求崩溃驱动的被动去库存，而是港口基建、铁路运营等系统效率提升的结果，供应链新常态可能提前到来。
  \item JPM的EM USD 10s30s利差曲线报告显示，新兴市场债券的期限结构正从'风险偏好驱动'转向'信用质量驱动'，投资级债券的久期溢价被系统性压缩，而高收益债券的期限溢价几乎未动。
\end{itemize}
\begin{figure}[htbp]
\centering
\includegraphics[width=0.88\linewidth]{figures/fig_150.jpg}
\caption{Exhibit 9 - Exhibit 1: Our weekly composite index increased in the most recent week (-4\% w/w); the bottleneck scale remained at '2'; overall bottleneck levels remain well below peak congestion levels when scale was at '10' and now imply levels}
\end{figure}
\begin{figure}[htbp]
\centering
\includegraphics[width=0.88\linewidth]{figures/fig_151.jpg}
\caption{Exhibit 2 - Exhibit 2: Our average weekly bottleneck score in June is tracking at 2.0; this result is down significantly from the Dec21/Jan22 peak of congestion and close to in line with the pre-Covid baseline GS Weekly Congestion Scale, Score}
\end{figure}
\begin{figure}[htbp]
\centering
\includegraphics[width=0.88\linewidth]{figures/fig_166.jpg}
\caption{Figure 1 - Figure 1: EM USD Aggregate 10s30s}
\end{figure}
\begin{figure}[htbp]
\centering
\includegraphics[width=0.88\linewidth]{figures/fig_167.jpg}
\caption{Figure 2 - Figure 2: UST 10s30s yield curve slope}
\end{figure}
{\small\color{refgray}\textbf{References:} [R001] JPM：市场真正定价的不是衰退，而是财政与通胀的再平衡; [R011] 美国银行：市场真正低估的不是AI泡沫，而是通胀二次探底对仓位结构的重塑; [R023] GS：美国供应链的“新常态”比市场预期的更早到来; [R024] JPM：新兴市场债券曲线正在发出被低估的信用分化信号; [R034] GS：市场误读了美联储的利率路径，最不确定的不是加息，而是何时降息}

\section{AI/算力与半导体：从GPU短缺到系统级重构，互联与存储定价权转移}
\textbf{AI基础设施的投资重心正从算力芯片转向系统架构与关键组件，互联层、存储（HDD）和ABF载板等环节的定价权正在发生结构性转移，市场对此定价不足。}

\begin{itemize}
  \item MS指出，AI基础设施正经历从'GPU短缺'到'系统级重构'的转折点，下一代平台（Vera Rubin、Kyber架构）带来电源架构、散热、PCB等设计升级，价值分配正从芯片端向系统集成与关键组件端迁移。
  \item MS半导体团队调研发现，互联层正从'必要组件'升级为'系统架构的定义者'，Astera Labs的PCIe交换机被客户称为'开放生态的NVSwitch'，Marvell的互联业务年增70\%且无放缓迹象。
  \item MS大幅上调HDD（硬盘驱动器）盈利预测，认为供需缺口至少延续到2028年，需求以40-50\%的速度增长而供给仅增30-35\%，定价权正从买方转向卖方，近线HDD价格目标在2-3年内从14.5美元/TB提升至25-30美元/TB。
  \item Bernstein通过韩国5月半导体设备进口数据（同比+51\%）验证了记忆体结构性扩张的加速，光刻设备进口创季度第二高，测试设备进口同比+103\%，但日本设备商TEL的收入动能出现边际放缓。
  \item MS关于ABF载板的报告指出，AT\&S的扩产（15-20亿欧元）由AMD等客户长期订单支持，这意味着供给短缺被提前证实而非推迟，2027年起ABF载板将进入供给短缺周期。
  \item GS路演反馈显示，投资者关注点已从AI需求转向货币化验证和资本支出可持续性，CPU因Agentic AI推理需求被重新定价，台积电的利润弹性可能被低估。
  \item Citi报告指出，AI推理层正从'成本竞赛'转向'稀缺性定价'，最前沿模型推理价格六周内翻倍而性能仅提升4个点，A100 GPU租赁价格上涨11\%，稀缺性正在被货币化。
\end{itemize}
\begin{figure}[htbp]
\centering
\includegraphics[width=0.88\linewidth]{figures/fig_001.jpg}
\caption{Exhibit 1 - EXHIBIT 1: May South Korea imports for SPE was \textbackslash{}\$3.2bn, -5\% MoM.}
\end{figure}
\begin{figure}[htbp]
\centering
\includegraphics[width=0.88\linewidth]{figures/fig_005.jpg}
\caption{Exhibit 6 - EXHIBIT 5: May tester imports from Japan and Malaysia collectively was +5\% MoM.}
\end{figure}
\begin{figure}[htbp]
\centering
\includegraphics[width=0.88\linewidth]{figures/fig_011.jpg}
\caption{Exhibit 13 - EXHIBIT 13: South Korea monthly SPE imports from Netherlands shows good directional correlation with ASML's South Korea revenue divided by 3. May Litho Imports of €928mn increased 28\% MoM and \textbackslash{}\textasciitilde{}150\% YoY.}
\end{figure}
\begin{figure}[htbp]
\centering
\includegraphics[width=0.88\linewidth]{figures/fig_269.jpg}
\caption{Exhibit 1 - Exhibit 1: Our new base case STX forecast implies nearline price/EB of \textbackslash{}\$19 in CY27 and \textbackslash{}\$22 in CY28 vs. \textbackslash{}\$22/\textbackslash{}\$29 in our new bull case forecasts}
\end{figure}
\begin{figure}[htbp]
\centering
\includegraphics[width=0.88\linewidth]{figures/fig_270.jpg}
\caption{Exhibit 2 - Exhibit 2: Our new base case WDC forecast implies nearline price/EB near \textbackslash{}\$19.50 in CY27 and \textbackslash{}\$22 in CY28 vs. \textbackslash{}\$22.50/\textbackslash{}\$29 in our new bull case forecasts}
\end{figure}
{\small\color{refgray}\textbf{References:} [R003] 摩根斯坦利：存储器价格的结构性重置，远比市场理解的更深远; [R004] Bernstein：韩国设备进口数据揭示的不是周期复苏，而是记忆体结构性扩张的加速; [R010] 摩根斯坦利：ABF载板供需拐点已至，真正稀缺的不是产能，是客户锁定的产能; [R014] GS：AI ASIC仍是首选主题，但市场低估了CPU和下一个供应链瓶颈的定价权; [R020] 摩根斯坦利：AI基础设施正在经历从“GPU短缺”到“系统级重构”的转折点; [R021] 摩根斯坦利：AI互联正在从“规模扩张”转向“架构重构”，市场低估了互联层定价权的转移; [R026] 摩根斯坦利：市场低估的不是AI需求本身，而是需求在亚洲供应链中产生的结构性再定价; [R035] Citi：AI推理层正在从“成本竞赛”转向“稀缺性定价”，市场低估了中间层的结构性机会; [R037] 摩根斯坦利：市场正在严重低估HDD周期的长度与定价弹性}

\section{能源与大宗：供需错位结构性固化，上游资本开支超级周期存疑}
\textbf{即便中东供应恢复，原油市场仍面临供需错位格局，上游资本开支的'超级周期'叙事可能被证伪；中国钛白粉出口韧性正在重塑全球定价权。}

\begin{itemize}
  \item MS在美伊备忘录签署后下调油价预测，但强调物理原油市场的疲软早在政治突破前就已形成，驱动因素是美国出口高企（同比+400万桶/天）与中国进口低迷（同比-600万桶/天），短期内看不到逆转迹象。
  \item MS指出，即便中东日均约1100万桶的产量处于停产，现货市场已出现反常疲软，布伦特现货价格、迪拜升贴水等指标均显示未售出货盘高于正常水平。
  \item Bernstein认为2026年全球上游资本开支不仅不会进入超级周期，反而可能出现约400亿美元的收缩，现金流的改善并未转化为投资意愿，资本纪律比油价更顽固。
  \item Bernstein指出，已发现资源中约三分之一尚未开发，但深海油转化为产量难度大，页岩油产量进入平台期，边际增长桶的控制权正在从美国独立公司转向国家石油公司。
  \item GS报告显示，中国钛白粉出口韧性远超预期，2026年前四个月净出口量同比增加8.8万吨（+14\%），即便面临硫磺供应紧张和印度反倾销税不确定性，出口供给弹性依然强劲。
  \item GS认为中国钛白粉的高出口量将压制2025年下半年西方生产商预期的涨价幅度，全球钛白粉定价权正从供给侧发生结构性转移。
\end{itemize}
\begin{figure}[htbp]
\centering
\includegraphics[width=0.88\linewidth]{figures/fig_027.jpg}
\caption{Exhibit 6 - Exhibit 6: US seaborne net exports are nearly 4 mb/d higher than last year...}
\end{figure}
\begin{figure}[htbp]
\centering
\includegraphics[width=0.88\linewidth]{figures/fig_028.jpg}
\caption{Exhibit 7 - Exhibit 7: ...whilst China's seaborne net imports are down \textbackslash{}\textasciitilde{} 6mb/d}
\end{figure}
\begin{figure}[htbp]
\centering
\includegraphics[width=0.88\linewidth]{figures/fig_225.jpg}
\caption{EXHIBIT 1 - EXHIBIT 1: We forecast \textbackslash{}\textasciitilde{}\textbackslash{}\$560 bln in upstream capital expenditures in 2026 including \textbackslash{}\textasciitilde{}\textbackslash{}\$56 bln toward exploration}
\end{figure}
\begin{figure}[htbp]
\centering
\includegraphics[width=0.88\linewidth]{figures/fig_226.jpg}
\caption{EXHIBIT 2 - EXHIBIT 2: The world has discovered abundant resources, largely in shale and shelf}
\end{figure}
\begin{figure}[htbp]
\centering
\includegraphics[width=0.88\linewidth]{figures/fig_302.jpg}
\caption{Exhibit 1 - Exhibit 1: Chinese Net Exports}
\end{figure}
\begin{figure}[htbp]
\centering
\includegraphics[width=0.88\linewidth]{figures/fig_303.jpg}
\caption{Exhibit 2 - Exhibit 2: Chinese Net TiO2 Exports MoM, in kt}
\end{figure}
{\small\color{refgray}\textbf{References:} [R013] 摩根斯坦利：油价真正的风险不是地缘缓和，而是供需错位的结构性固化; [R029] Bernstein：上游资本开支的“超级周期”叙事，市场可能信错了; [R041] GS：中国钛白粉出口韧性正在重塑全球定价权，而非简单的产能外溢}

\section{地缘风险与汇率：和平红利下的货币重估与亚洲资产重新定价}
\textbf{美伊和平协议触发全球资产定价逻辑从'地缘风险溢价'向'宏观基本面与相对价值'切换，油价下行对亚洲的影响远不止于通胀改善，而是央行政策路径与AI资本开支叠加效应的结构性重置。}

\begin{itemize}
  \item NOM认为市场低估了中东协议对全球资产定价的连锁效应，交易重心应从方向性押注转向基于各国本地因素的多空配对交易，短期看弱美元。
  \item DB筛选出四支被低估的货币——瑞典克朗、南非兰特、智利比索和印度卢比，认为市场对'和平红利'的定价仍停留在情绪层面，忽略了供给端成本结构变化的长期含义。
  \item NOM指出，油价每跌10\%，亚洲经常账户平均改善约0.3\%的GDP，泰国、印度、菲律宾和韩国是最大赢家，但受益程度和方式存在显著分化。
  \item Bernstein认为美伊和平对印度市场是'反弹而非反转'，真正的催化剂是伊朗原油制裁豁免和恰巴哈尔港的战略价值，这两点当前定价最不充分。
  \item NOM的人民币中间价模型显示，模型预测值（6.7575）较实际中间价系统性偏低，逆周期因子调节力度正在减弱，中间价形成逻辑可能从'被动跟随'转向'主动引导'。
  \item GS的5月中国外汇流入数据显示净流入约500亿美元，结构上资本金融项目出现边际改善，外资通过债券通净买入130亿人民币债券，为2025年4月以来首次。
\end{itemize}
\begin{figure}[htbp]
\centering
\includegraphics[width=0.88\linewidth]{figures/fig_021.jpg}
\caption{Figure 1 - Figure 1: SEK has underperformed since March G10 FX performance: since start March vs last 3d}
\end{figure}
\begin{figure}[htbp]
\centering
\includegraphics[width=0.88\linewidth]{figures/fig_024.jpg}
\caption{Figure 4 - Figure 4: INR positioning remains short}
\end{figure}
\begin{figure}[htbp]
\centering
\includegraphics[width=0.88\linewidth]{figures/fig_014.jpg}
\caption{Exhibit 1 - Exhibit 1: Our preferred gauge suggests increased FX inflows in May}
\end{figure}
\begin{figure}[htbp]
\centering
\includegraphics[width=0.88\linewidth]{figures/fig_016.jpg}
\caption{Exhibit 3 - Exhibit 3: Foreign investors started to buy CNY bonds in May China monthly bond flow in the interbank market}
\end{figure}
{\small\color{refgray}\textbf{References:} [R002] NOM：市场低估了中东协议对全球资产定价的连锁效应，而非仅关注地缘风险本身; [R005] Bernstein：美伊和平对印度而言是反弹，不是反转; [R012] DB：地缘风险溢价退潮，最被低估的不是避险资产，而是这些被错杀的货币; [R018] NOM：油价暴跌不是短期交易，而是亚洲资产重新定价的起点; [R028] NOM：人民币中间价模型正在发出一个被市场忽视的信号; [R006] GS：5月中国外汇流入的真实信号被市场低估了}

\section{中国权益与地产：结构性供给优势被低估，城市更新资金落地成关键变量}
\textbf{中国市场的核心叙事正从'周期性复苏'切换到'结构性供给优势'，出口韧性与内需疲弱的分化加剧，城市更新资金落地可能成为打破地产市场均衡点的外力。}

\begin{itemize}
  \item MS认为中国在全球AI和能源转型资本开支超级周期中的供给侧地位被低估，中国正在成为全球资本品和电子供应链中不可替代的'供给锚点'，出口市场份额有望在2030年达到16.5\%。
  \item MS指出，MSCI中国年初至今-10\%的表现掩盖了结构性分化，CSI 300同期上涨5\%，夏普比率在2025年达到1.44，高于MSCI新兴市场的1.70和标普500的0.66。
  \item MS报告揭示了中国经济的深层裂痕：5月出口在AI周期推动下全面回暖（电子集成电路同比+110\%），但私人信贷需求全面恶化，居民短期和中长期贷款均为净偿还。
  \item GS认为市场低估了'城市更新'作为房地产下行缓冲器的真实权重，中央层面对城市更新的资金支持正从框架走向明细，257亿超长期特别国债专项用于地下管网升级。
  \item GS数据显示，新房成交量进入'平台期'（环比+1\%，同比持平），市场正处在一个需要外力打破的均衡点，而城市更新资金落地正是这个外力。
  \item GS指出中国券商的结构性机会不在A股交易量，而在香港的资产负债表扩张，头部券商香港子公司杠杆（11倍）和ROE（16\%）均显著高于集团整体水平。
\end{itemize}
\begin{figure}[htbp]
\centering
\includegraphics[width=0.88\linewidth]{figures/fig_127.jpg}
\caption{Exhibit 1 - Exhibit 1: Primary GFA sold last week was \$+1\textbackslash{}\%\$ wow and flat yoy in c.75 cities}
\end{figure}
\begin{figure}[htbp]
\centering
\includegraphics[width=0.88\linewidth]{figures/fig_128.jpg}
\caption{Exhibit 2 - Exhibit 2: Primary GFA sold YTD on average was -13\% yoy in c.75 cities, and -11\%/-43\% vs. 2024/2023 level}
\end{figure}
\begin{figure}[htbp]
\centering
\includegraphics[width=0.88\linewidth]{figures/fig_119.jpg}
\caption{Exhibit 1 - Exhibit 1: Brokers are demonstrating strong performance, +7\%/+5\% vs. +1\%/+3\% of banks for A/H share since June 9. As of June 15 Price performance since June 9}
\end{figure}
\begin{figure}[htbp]
\centering
\includegraphics[width=0.88\linewidth]{figures/fig_120.jpg}
\caption{Exhibit 2 - Exhibit 2: The average leverage for the offshore subsidiaries of the three brokers in our coverage is 11x, compared to the group average leverage of 6x As of 2025}
\end{figure}
\begin{figure}[htbp]
\centering
\includegraphics[width=0.88\linewidth]{figures/fig_121.jpg}
\caption{Exhibit 3 - Exhibit 3: The average ROE for the offshore subsidiaries of the three brokers in our coverage is 16\%, compared to the group average leverage of 9\% As of 2025}
\end{figure}
{\small\color{refgray}\textbf{References:} [R008] 摩根斯坦利：中国权益市场正在被低估的并非需求，而是供给侧的结构性定价能力; [R009] 摩根斯坦利：贸易韧性掩盖了国内需求的深层裂痕; [R016] GS：中国券商真正的结构性机会不在A股交易量，而在香港的资产负债表扩张; [R017] GS：市场低估了“城市更新”作为中国房地产下行缓冲器的真实权重}

\section{亚洲央行与区域市场：从被动加息转向主动加息，市场定价不足}
\textbf{亚洲央行的货币政策叙事正从'弱势下的被动加息'进入'强势增长驱动的主动加息'阶段，韩国、台湾、新加坡正在酝酿的加息由科技繁荣和增长韧性驱动，市场对此定价不充分。}

\begin{itemize}
  \item JPM明确区分了两类加息：印尼和菲律宾属于'被迫加息'（汇率压力+通胀失控），而韩国、台湾、新加坡属于'主动加息'（科技出口强劲、增长超预期、通胀风险上升）。
  \item JPM预计菲律宾央行将加息50个基点以锚定通胀预期，印尼紧缩周期尚未结束，市场可能低估了其持续性的风险。
  \item 韩国1Q25 GDP上修至7.5\%，全年增速上调至3.7\%，央行转鹰，预计未来一年加息100bp。
  \item 台湾5月CPI 2.2\%、PPI 14.1\%，下周央行会议很关键，市场关注是否暗示年底前加息。
  \item BofA的6月亚太基金经理调查显示，能源安全焦虑骤降（从91\%降至18\%），亚太（除日本）企业盈利预期进一步走强，净看多比例达50\%。
  \item BofA调查显示，科技硬件在仓位中的权重首次超过半导体，印尼取代印度成为最不受欢迎的市场，北亚依然是区域配置的焦点。
\end{itemize}
\begin{figure}[htbp]
\centering
\includegraphics[width=0.88\linewidth]{figures/fig_097.jpg}
\caption{Exhibit 2 - Exhibit 2: APAC ex-Japan growth expectations turned positive in June, marking a continued recovery from April's sharp downturn Net \% expecting a stronger Global / APAC ex-Japan economy}
\end{figure}
\begin{figure}[htbp]
\centering
\includegraphics[width=0.88\linewidth]{figures/fig_101.jpg}
\caption{Exhibit 6 - Exhibit 6: Energy security concerns have normalized: high/extreme concern dropped 73pts since Apr, while moderate concern became the majority view How concerned are you about energy security risks for APAC region in current geopoli}
\end{figure}
\begin{figure}[htbp]
\centering
\includegraphics[width=0.88\linewidth]{figures/fig_116.jpg}
\caption{Exhibit 21 - Exhibit 21: FMS investors are now more overweight in Tech Hardware than Semis, with Financial Services emerging as a new favored sector Asia Pacific ex-Japan sector sentiment: Net \% overweight APAC ex-Japan sector sentiment: Net \% F}
\end{figure}
\begin{figure}[htbp]
\centering
\includegraphics[width=0.88\linewidth]{figures/fig_217.jpg}
\caption{Figure 1 - Figure 1: Indonesia - BI-FXPI versus change in policy rate}
\end{figure}
\begin{figure}[htbp]
\centering
\includegraphics[width=0.88\linewidth]{figures/fig_219.jpg}
\caption{Figure 3 - Figure 3: Philippines headline CPI vs. diffusion index \$\textasciicircum{}\{1\}\$}
\end{figure}
{\small\color{refgray}\textbf{References:} [R025] JPM：亚洲正在从“被迫加息”转向“主动加息”，市场对此定价不足; [R015] 美国银行：市场真正低估的不是需求，而是AI定价权从算力向硬件的转移}

\section{能源转型与电力基础设施：AI基建的真正瓶颈在电力供给}
\textbf{AI基础设施的真正瓶颈不在资本而在电力，电力变压器交货周期拉长至128-144周，数据中心选址逻辑正从'靠近用户'转向'靠近电力'，先进核能商业化路径正从政府补贴转向科技巨头需求驱动。}

\begin{itemize}
  \item MS指出，电力变压器交货周期已从此前的12-16周拉长至平均128周，美国电网互联积压项目规模已超过现有装机容量的两倍，电力供给正从'配套条件'升级为'独立瓶颈'。
  \item MS认为，数据中心选址逻辑正在重构，开发者必须首先确认电力可用性才能推进项目，拥有富余电力容量的区域将获得显著竞争优势。
  \item JPM与TerraPower的炉边对话揭示，Meta选择Natrium反应堆并提供里程碑式开发资金，标志着先进核能商业化路径从'政府补贴驱动'转向'科技巨头需求驱动'。
  \item JPM指出，先进核能的真正拐点不是技术，而是客户愿意为第一座电站付费，Meta的参与解决了'先有鸡还是先有蛋'的FOAK（首堆）风险问题。
  \item MS关于天然气轮市场的报告显示，市场对需求侧的叙事已高度拥挤，但供给侧的结构性变化——非传统发电解决方案的产能扩张——正在重塑2030年的竞争格局。
  \item Bernstein关于液冷数据中心的报告指出，冷却液分配单元（CDU）是比冷板更被低估的环节，兼具技术壁垒、服务粘性和长期定价权，市场规模到2030年将达到中高十亿美元级别。
\end{itemize}
\begin{figure}[htbp]
\centering
\includegraphics[width=0.88\linewidth]{figures/fig_137.jpg}
\caption{EXHIBIT 3 - EXHIBIT 3: Overview of Liquid Cooling and Key Components ```mermaid graph TD}
\end{figure}
\begin{figure}[htbp]
\centering
\includegraphics[width=0.88\linewidth]{figures/fig_138.jpg}
\caption{EXHIBIT 4 - EXHIBIT 4: Still Early Days for Liquid Cooling Evolution of frontier GPU rack computing power Rack power (kW) for frontier chips}
\end{figure}
\begin{figure}[htbp]
\centering
\includegraphics[width=0.88\linewidth]{figures/fig_253.jpg}
\caption{Exhibit 1 - Exhibit 1: Change to 2028 consensus EBITA following results and the 2025 CMD (14th Nov), and change in share price in the next 30 days. More recently, as consensus has moved higher, the rate of consensus upgrades has been decreasin}
\end{figure}
\begin{figure}[htbp]
\centering
\includegraphics[width=0.88\linewidth]{figures/fig_254.jpg}
\caption{Exhibit 2 - Exhibit 2: MS Gas turbine / power solutions supply model. We have updated our supply model for recent announcements, and the potential 2030 supply rises to 135GW by 2030 (we forecast 116GW supply in our Jan 2026 report: The capacit}
\end{figure}
{\small\color{refgray}\textbf{References:} [R027] 摩根斯坦利：AI基建的真正瓶颈不在资本，而在电力与结构性供给约束; [R033] 摩根斯坦利：天然气轮市场的真正拐点不在2026，而在供给侧的隐性再定价; [R036] JPM：先进核能真正的拐点不是技术，而是客户愿意为第一座电站付费; [R019] Bernstein：液冷数据中心最被低估的环节不是冷板，而是冷却液分配单元}

\section{生物科技与医疗：并购结构性加速，IRA谈判规则影响被低估}
\textbf{生物科技并购正从偶发事件变为系统性力量，大型药企为填补未来5年产品线缺口而加速并购；CMS拟议规则将皮下注射制剂纳入价格谈判范围，对默沙东、强生等公司的影响被市场低估。}

\begin{itemize}
  \item MS指出，2026年生物科技并购活动正在加速，GSK以106亿美元收购Nuvalent、Incyte以20亿美元收购Vega Therapeutics等交易密集程度超出简单的'行业整合'叙事，并购是结构性的而非偶发性的。
  \item MS认为，并购的驱动力是大药企面临的专利悬崖压力，2026年至今IPO平均规模已达3.3亿美元，远超2025年全年水平，SMID Cap生物科技正进入'复兴期'。
  \item GS报告指出，CMS拟议规则将皮下注射制剂纳入2029年药品价格谈判候选范围，Keytruda Qlex、Opdivo Qvantig和Darzalex Faspro面临额外定价压力。
  \item GS量化分析显示，在最保守假设下，强生Darzalex Faspro年销售额影响约19亿美元（占EBIT约4\%），默沙东Keytruda Qlex影响约9亿美元（占EBIT约3\%）。
  \item 市场对皮下制剂纳入谈判的反应迅速（消息公布当日默沙东-2.8\%，强生-2.2\%），但GS认为影响远不止于表面数字，这是对行业定价策略的结构性打击。
\end{itemize}
{\small\color{refgray}\textbf{References:} [R031] GS：市场低估了IRA谈判规则对皮下注射制剂的连锁影响; [R032] 摩根斯坦利：生物科技并购不是脉冲，而是结构性的重新定价}

\section{消费与科技硬件：AI眼镜渠道扩容与印度混动拐点}
\textbf{AI眼镜的购买决策正从'电子消费品'向'视力健康产品'迁移，能完成处方闭环的零售商才是真正赢家；印度强混动市场的拐点即将到来，车型供给翻倍将驱动渗透率从2.3\%提升至7-8\%。}

\begin{itemize}
  \item GS评估Best Buy与Meta合作对眼镜零售商的影响，认为对于National Vision和Warby Parker并非威胁，真正构成竞争壁垒的是能否在同一个场景里完成验光、处方、镜片定制和框架适配的闭环。
  \item GS指出，超过50\%的Best Buy消费者希望在购买AI眼镜前亲眼看到实物，消费者决策链条的核心是确认'这副眼镜戴在我脸上是否合适、是否可以配上我的处方'。
  \item Bernstein认为印度强混动市场的规模几乎完全由车型供给决定，2025年仅有8款在售车型却实现83\%的销量增长，供给侧的拐点即将到来。
  \item Bernstein预测，印度强混动渗透率将从当前的2.3\%在2030年前提升至7-8\%，驱动因素不是政策变化，而是车型供给的集中释放（从8款增至约27款）。
  \item MS关于欧洲软件股的报告显示，自2021年高点以来欧洲应用软件板块已连续近1700天未创新高，累计回撤超35\%，但AI技术浪潮正在重塑行业，数据基础设施和网络安全表现相对坚挺。
\end{itemize}
\begin{figure}[htbp]
\centering
\includegraphics[width=0.88\linewidth]{figures/fig_301.jpg}
\caption{Exhibit 3 - Exhibit 3: Search intensity for Best Buy AI glasses has largely led peers, though we see an uptick in Warby Parker trends in recent weeks Trailing 4 week search intensity for: ai glasses best buy (US); america's best meta (US); ai}
\end{figure}
\begin{figure}[htbp]
\centering
\includegraphics[width=0.88\linewidth]{figures/fig_293.jpg}
\caption{EXHIBIT 2 - \#\# INDIA BUYS AS MANY HYBRIDS AS IT IS OFFERED Strong Hybrid Retail Volumes}
\end{figure}
\begin{figure}[htbp]
\centering
\includegraphics[width=0.88\linewidth]{figures/fig_296.jpg}
\caption{EXHIBIT 6 - EXHIBIT 6: India will go from 8 strong-hybrid nameplates in 2025 to roughly 27 by 2030; expansion led by Hyundai's confirmed 8 by 2030 plan, Maruti's in-house series-HEV programme, and a wave of Korean, Japanese and Indian launches}
\end{figure}
{\small\color{refgray}\textbf{References:} [R040] GS：AI眼镜渠道扩容的真正赢家不是Best Buy，而是能完成处方闭环的零售商; [R039] Bernstein：印度市场真正低估的不是电动化，而是混动供给侧的再定价; [R038] 摩根斯坦利：欧洲软件股正经历近二十年最深的估值回撤，但市场忽视了结构性分化}

\section{交通运输：油运定价权结构性重置，VLCC日租金弹性飙升}
\textbf{霍尔木兹海峡若重开，最被低估的不是油价，而是油运定价权的结构性重置，VLCC市场集中度提升使运价弹性翻了3倍以上，日租金可能飙升至25-35万美元/天。}

\begin{itemize}
  \item GS报告指出，2021至2025年间VLCC市场每1个百分点的净需求变化仅带来约6000美元/天的运价波动，而2026年前两个月这个弹性系数已跃升至2万美元/天，翻了3倍以上。
  \item GS认为，前十大合规船东的市场份额从2025年的47\%飙升至68\%，集中度在一年内完成质的飞跃，头部船东有能力通过控制运力投放节奏将运价维持在高位更长时间。
  \item GS情景推演显示，若霍尔木兹海峡在7月底前恢复出口，VLCC日租金可能冲到25万美元/天；若叠加伊朗制裁解除，甚至可能到35万美元/天，对应航运公司盈利57\%和114\%的上行空间。
  \item GS指出，全球原油库存比2月低了7.05亿桶，补库存需求将直接拉动运价，而供给侧的议价权已从分散走向集中。
\end{itemize}
\begin{figure}[htbp]
\centering
\includegraphics[width=0.88\linewidth]{figures/fig_141.jpg}
\caption{Exhibit 3 - Exhibit 2: We estimate that restocking could add 4.8\% of global seaborne crude oil trade demand if Gulf exports normalize by end of Jul'26 Global crude inventory change vs. Feb'26}
\end{figure}
\begin{figure}[htbp]
\centering
\includegraphics[width=0.88\linewidth]{figures/fig_142.jpg}
\caption{Exhibit 3 - Exhibit 3: VLCC TCE elasticity to crude tanker net demand / (supply) VLCC TCE (China)'s elasticity (US\textbackslash{}\$'000/day) to crude tanker S-D change}
\end{figure}
{\small\color{refgray}\textbf{References:} [R022] GS：霍尔木兹海峡若重开，最被低估的不是油价，而是油运定价权的结构性重置}

\section{结语}
综合来看，当前市场正处于多重结构性力量交织的节点：地缘风险溢价退潮、AI基础设施从粗放扩张进入精细化重构、全球利率定价锚点从政策利率转向财政供给、以及亚洲央行货币政策立场的主动切换。投资者需要超越短期情绪波动，聚焦于供给侧的结构性变化和定价权的转移，这些才是决定未来资产回报的核心变量。
\newpage
\section{报告覆盖清单}
\begin{itemize}
  \item [R001] 宏观与利率：财政供给与通胀再平衡主导定价 - JPM：市场真正定价的不是衰退，而是财政与通胀的再平衡
  \item [R002] 地缘风险与汇率：和平红利下的货币重估与亚洲资产重新定价 - NOM：市场低估了中东协议对全球资产定价的连锁效应，而非仅关注地缘风险本身
  \item [R003] AI/算力与半导体：从GPU短缺到系统级重构，互联与存储定价权转移 - 摩根斯坦利：存储器价格的结构性重置，远比市场理解的更深远
  \item [R004] AI/算力与半导体：从GPU短缺到系统级重构，互联与存储定价权转移 - Bernstein：韩国设备进口数据揭示的不是周期复苏，而是记忆体结构性扩张的加速
  \item [R005] 地缘风险与汇率：和平红利下的货币重估与亚洲资产重新定价 - Bernstein：美伊和平对印度而言是反弹，不是反转
  \item [R006] 地缘风险与汇率：和平红利下的货币重估与亚洲资产重新定价 - GS：5月中国外汇流入的真实信号被市场低估了
  \item [R007] 未被主题摘要引用 - JPM：市场真正低估的不是增长，而是信用脉冲对风格切换的定价权
  \item [R008] 中国权益与地产：结构性供给优势被低估，城市更新资金落地成关键变量 - 摩根斯坦利：中国权益市场正在被低估的并非需求，而是供给侧的结构性定价能力
  \item [R009] 中国权益与地产：结构性供给优势被低估，城市更新资金落地成关键变量 - 摩根斯坦利：贸易韧性掩盖了国内需求的深层裂痕
  \item [R010] AI/算力与半导体：从GPU短缺到系统级重构，互联与存储定价权转移 - 摩根斯坦利：ABF载板供需拐点已至，真正稀缺的不是产能，是客户锁定的产能
  \item [R011] 宏观与利率：财政供给与通胀再平衡主导定价 - 美国银行：市场真正低估的不是AI泡沫，而是通胀二次探底对仓位结构的重塑
  \item [R012] 地缘风险与汇率：和平红利下的货币重估与亚洲资产重新定价 - DB：地缘风险溢价退潮，最被低估的不是避险资产，而是这些被错杀的货币
  \item [R013] 能源与大宗：供需错位结构性固化，上游资本开支超级周期存疑 - 摩根斯坦利：油价真正的风险不是地缘缓和，而是供需错位的结构性固化
  \item [R014] AI/算力与半导体：从GPU短缺到系统级重构，互联与存储定价权转移 - GS：AI ASIC仍是首选主题，但市场低估了CPU和下一个供应链瓶颈的定价权
  \item [R015] 亚洲央行与区域市场：从被动加息转向主动加息，市场定价不足 - 美国银行：市场真正低估的不是需求，而是AI定价权从算力向硬件的转移
  \item [R016] 中国权益与地产：结构性供给优势被低估，城市更新资金落地成关键变量 - GS：中国券商真正的结构性机会不在A股交易量，而在香港的资产负债表扩张
  \item [R017] 中国权益与地产：结构性供给优势被低估，城市更新资金落地成关键变量 - GS：市场低估了“城市更新”作为中国房地产下行缓冲器的真实权重
  \item [R018] 地缘风险与汇率：和平红利下的货币重估与亚洲资产重新定价 - NOM：油价暴跌不是短期交易，而是亚洲资产重新定价的起点
  \item [R019] 能源转型与电力基础设施：AI基建的真正瓶颈在电力供给 - Bernstein：液冷数据中心最被低估的环节不是冷板，而是冷却液分配单元
  \item [R020] AI/算力与半导体：从GPU短缺到系统级重构，互联与存储定价权转移 - 摩根斯坦利：AI基础设施正在经历从“GPU短缺”到“系统级重构”的转折点
  \item [R021] AI/算力与半导体：从GPU短缺到系统级重构，互联与存储定价权转移 - 摩根斯坦利：AI互联正在从“规模扩张”转向“架构重构”，市场低估了互联层定价权的转移
  \item [R022] 交通运输：油运定价权结构性重置，VLCC日租金弹性飙升 - GS：霍尔木兹海峡若重开，最被低估的不是油价，而是油运定价权的结构性重置
  \item [R023] 宏观与利率：财政供给与通胀再平衡主导定价 - GS：美国供应链的“新常态”比市场预期的更早到来
  \item [R024] 宏观与利率：财政供给与通胀再平衡主导定价 - JPM：新兴市场债券曲线正在发出被低估的信用分化信号
  \item [R025] 亚洲央行与区域市场：从被动加息转向主动加息，市场定价不足 - JPM：亚洲正在从“被迫加息”转向“主动加息”，市场对此定价不足
  \item [R026] AI/算力与半导体：从GPU短缺到系统级重构，互联与存储定价权转移 - 摩根斯坦利：市场低估的不是AI需求本身，而是需求在亚洲供应链中产生的结构性再定价
  \item [R027] 能源转型与电力基础设施：AI基建的真正瓶颈在电力供给 - 摩根斯坦利：AI基建的真正瓶颈不在资本，而在电力与结构性供给约束
  \item [R028] 地缘风险与汇率：和平红利下的货币重估与亚洲资产重新定价 - NOM：人民币中间价模型正在发出一个被市场忽视的信号
  \item [R029] 能源与大宗：供需错位结构性固化，上游资本开支超级周期存疑 - Bernstein：上游资本开支的“超级周期”叙事，市场可能信错了
  \item [R030] 未被主题摘要引用 - Bernstein：福克斯收购Roku，战略逻辑成立但价格昂贵，市场低估了广告协同的潜力
  \item [R031] 生物科技与医疗：并购结构性加速，IRA谈判规则影响被低估 - GS：市场低估了IRA谈判规则对皮下注射制剂的连锁影响
  \item [R032] 生物科技与医疗：并购结构性加速，IRA谈判规则影响被低估 - 摩根斯坦利：生物科技并购不是脉冲，而是结构性的重新定价
  \item [R033] 能源转型与电力基础设施：AI基建的真正瓶颈在电力供给 - 摩根斯坦利：天然气轮市场的真正拐点不在2026，而在供给侧的隐性再定价
  \item [R034] 宏观与利率：财政供给与通胀再平衡主导定价 - GS：市场误读了美联储的利率路径，最不确定的不是加息，而是何时降息
  \item [R035] AI/算力与半导体：从GPU短缺到系统级重构，互联与存储定价权转移 - Citi：AI推理层正在从“成本竞赛”转向“稀缺性定价”，市场低估了中间层的结构性机会
  \item [R036] 能源转型与电力基础设施：AI基建的真正瓶颈在电力供给 - JPM：先进核能真正的拐点不是技术，而是客户愿意为第一座电站付费
  \item [R037] AI/算力与半导体：从GPU短缺到系统级重构，互联与存储定价权转移 - 摩根斯坦利：市场正在严重低估HDD周期的长度与定价弹性
  \item [R038] 消费与科技硬件：AI眼镜渠道扩容与印度混动拐点 - 摩根斯坦利：欧洲软件股正经历近二十年最深的估值回撤，但市场忽视了结构性分化
  \item [R039] 消费与科技硬件：AI眼镜渠道扩容与印度混动拐点 - Bernstein：印度市场真正低估的不是电动化，而是混动供给侧的再定价
  \item [R040] 消费与科技硬件：AI眼镜渠道扩容与印度混动拐点 - GS：AI眼镜渠道扩容的真正赢家不是Best Buy，而是能完成处方闭环的零售商
  \item [R041] 能源与大宗：供需错位结构性固化，上游资本开支超级周期存疑 - GS：中国钛白粉出口韧性正在重塑全球定价权，而非简单的产能外溢
\end{itemize}
\section{逐篇报告摘录}
\subsection{[R001] JPM：市场真正定价的不是衰退，而是财政与通胀的再平衡}
{\small\color{refgray}归类：宏观与利率：财政供给与通胀再平衡主导定价}

这份JPM最新宏观策略报告的核心判断，值得每一位资产配置者和产业决策者仔细推敲：当前金融市场正在定价的，并非一次传统意义上的经济衰退或政策宽松周期，而是一个由财政主导、通胀粘性与供给冲击相互缠绕的“再平衡”过程。全球利率曲线、汇率走势和大宗商品定价，都正在被这一深层逻辑重新校准。

报告明确传递了一个信号：市场对美联储降息的定价过于激进，而对财政供给压力和地缘供给侧冲击的持续性则可能低估。这不是一份简单的“看多美债”或“看空美股”的报告，而是一幅关于全球资产定价锚点正在发生结构性位移的路线图。

为什么现在重要？因为过去几个月，市场在“软着陆”叙事与“衰退交易”之间反复摇摆，但真正的结构性力量——美国财政赤字的长期化、全球央行紧缩的惯性、以及霍尔木兹海峡问题对能源和金属供应链的持续扰动——正在被短期情绪掩盖。JPM这份报告的价值，在于将这些力量系统性地纳入了一个可检验的分析框架。

报告最值得关注的一个判断，是美国利率定价的锚点正在发生迁移。传统上，10年期美债收益率主要由对未来联邦基金利率的预期主导，但JPM的模型显示，当前曲线定价中，财政供给因素和期限溢价的权重正在显著上升。

报告明确指出，由于2027财年及之后存在巨大的融资缺口，美国财政部预计将在今年8月移除其前瞻指引中的“至少”措辞，为从2027年2月开始的多季度系列息票增发做准备。这意味着，即使美联储维持利率不变，长端利率也可能因为供给压力而系统性抬升。JPM预测，2年期和10年期美债收益率将在2026年底分别升至4.20\%和4.70\%。

这一判断的含义是：投资者不能再用“美联储何时降息”来推断长端利率走势。财政供给正在成为独立的定价变量。对于配置型机构而言，这意味着需要重新评估“持有长债”的风险回报，尤其是在期限溢价处于历史低位、而供给即将放量的背景下。

2. 美联储的“新范式”：缩表与降息可以同时进行，但市场可能低估了难度

一个容易被忽视的细节是，报告专门讨论了新任主席沃什的货币政策框架。沃什主张同时推进缩表与降息，认为这样可以减少美联储在市场中的足迹，同时对金融条件产生更中性的影响。但报告通过其公平价值模型测算了一个关键数据：即使美联储资产负债表规模缩减接近1万亿美元，也只能换来约25个基点的降息空间。

这个数字揭示了“新范式”的内在矛盾。市场目前定价的降息幅度远超模型所暗示的空间。如果沃什试图推动降息，其前提条件——大幅缩表——本身就会对金融条件产生收紧效应。换句话说，市场期待的“鸽派转向”可能以另一种方式实现，但最终对资产价格的影响远非线性。

对于投资者而言，这意味着需要区分“降息交易”和“缩表交易”。前者利好风险资产，后者则可能通过抽离流动性对风险偏好形成压制。当前市场可能过于关注降息本身，而低估了实现这一目标所需付出的流动性代价。

3. 霍尔木兹海峡的“薛定谔状态”：供给侧的半恢复比完全中断更具破坏性

报告在大宗商品部分使用了一个精妙的标题：“薛定谔的霍尔木兹海峡”。这并非故弄玄虚，而是精准描述了当前能源和金属市场的核心矛盾——海峡流量正在恢复，但远未回到战前水平。

数据显示，6月经霍尔木兹海峡的石油流量已从3月的220万桶/日回升至510万桶/日，但仍仅为战前水平的约25\%。西北欧的天然气储存处于历史同期低位。更值得关注的是，铝市场出现的严重缺口，正在缓慢地通过供应链传导，这个过程在海峡完全重新开放后仍将持续。

这种“半恢复”状态可能比完全中断更具破坏性。因为完全中断会触发价格急涨和需求破坏，从而快速出清市场；而半恢复则让市场持续处于“短缺但未恐慌”的状态，导致价格在高位缓慢爬升，供应链调整被延迟，最终积累更大的调整压力。

对于产业决策者而言，这意味着不能以“地缘风险消退”为由放松对供应链风险的管控。铝、天然气以

[... middle omitted ...]

指出，如果美联储对更强美国周期的确认被视为鹰派，美元将进一步走强，因为美元相对于利率公平价值仍然折价。历史经验表明，在加息周期的首次行动前6个月，美元通常会升值约5\%。

这意味着，新兴市场货币面临的压力可能持续。尽管报告对部分高收益新兴市场货币转为超配，但前提是这些国家的央行准备好加息。对于中国而言，人民币汇率面临的外部约束并未解除，美元强势周期可能比市场预期的更长。

5. 报告尚未完全回答的关键问题：通胀粘性的来源正在从“需求”转向“供给”

这份报告在多个领域都展现了深刻的洞察，但它留下了一个尚未完全回答的关键问题：通胀粘性的来源是否正在发生结构性转变？

报告预测美国核心PCE在2026年四季度仍将维持在3.4\%的高位，同时失业率将缓慢下降至4.1\%。这意味着“菲利普斯曲线”似乎再次变得陡峭——低失业率正在推动通胀粘性。但报告同时也指出，财政扩张、能源供给冲击和贸易政策（关税）正在从供给侧推高价格。这两种力量叠加在一起，使得通胀的“弹性”远高于传统模型预测。

\subsection{[R002] NOM：市场低估了中东协议对全球资产定价的连锁效应，而非仅关注地缘风险本身}
{\small\color{refgray}归类：地缘风险与汇率：和平红利下的货币重估与亚洲资产重新定价}

NOM：市场低估了中东协议对全球资产定价的连锁效应，而非仅关注地缘风险本身

这份NOM研报的核心判断，并非简单地“看多”或“看空”某个市场，而是提出一个更具结构性的观点：市场正在低估美伊临时协议签署后，全球资产定价逻辑从“地缘风险溢价”向“宏观基本面与相对价值”切换的深度和速度。NOM认为，真正的交易机会不在于赌局势升级或缓和，而在于捕捉这种定价逻辑切换带来的、由本地因素驱动的相对价值交易。

报告发布的时间点——美伊即将于6月19日签署谅解备忘录前夕——本身就极具信号意义。自5月22日以来，NOM已经基于“局势缓和”的假设调整了策略仓位。如今，随着协议签署的确认，NOM进一步强化了这一立场，并做出了数项关键调整。这并非一次简单的战术性调整，而是对全球宏观交易主题的一次系统性重估。

报告最值得关注的信号，不是某个具体的汇率目标价，而是其背后清晰的投资框架：在美元整体偏弱的背景下，交易的重心应从“做多或做空美元”这种笼统的方向性押注，转向“基于各国本地因素（经济增长、货币政策、资本流动、政治风险）进行多空配对交易”。这意味着，未来一段时间的赢家，将是那些能够精准识别并利用各国基本面差异的投资者。

研报明确指出，美伊临时协议的签署，其意义远超出中东地缘政治本身。它直接触发了全球资产定价框架的三个核心变量的变化：能源价格、美元走势和风险偏好。

首先，协议的核心内容之一——重新开放霍尔木兹海峡——直接指向能源供给侧。一旦海峡通航恢复，全球能源价格面临的下行压力将显著增加。这对于高度依赖能源进口的经济体（如新西兰、印度、韩国）是重大利好，而对于能源出口国（如澳大利亚）则构成拖累。NOM正是基于这一逻辑，维持了对澳元兑纽元的空头头寸。

其次，能源价格的回落将逆转此前因能源危机而强化的“贸易条件”效应，即能源出口国货币的强势和进口国货币的弱势。这种逆转，是NOM判断美元整体将走弱的核心论据之一。一个更弱的美元环境，为新兴市场货币，尤其是那些基本面扎实、出口竞争力强的货币（如人民币、新台币），提供了升值的空间。

最后，局势缓和直接提升了全球风险偏好。这并非简单的“避险情绪消退”，而是意味着资金将从“避险资产”（如美元、日元）和“防御性头寸”中流出，重新配置到那些此前因风险溢价而被低估的资产上。NOM报告中提到的“低beta于美元和石油”的交易（如做多瑞郎/日元、欧元/英镑），正是在这种风险偏好改善环境下，寻求相对价值的表现。

因此，这份研报的核心洞察是：不要仅仅把美伊协议看作一个“利好”或“利空”消息。它是一个触发器，迫使市场重新审视所有资产的定价基础——能源成本、美元流动性和风险情绪。谁能更早、更准确地理解这种定价框架的切换，谁就能在新的市场格局中占据先机。

2. 美元整体偏弱，但真正的机会在于“本地因素”驱动的相对价值交易

NOM对美元的整体判断是“选择性做空”。这并非一个无条件的看空美元宣言，而是一个有条件的策略：在美元整体走弱的背景下，寻找那些由各自经济体本地因素驱动、能够跑赢美元指数的货币对。

做空美元/人民币（目标6.60）：这是NOM高确信度的头寸。理由并非简单看空美元，而是基于中国的强劲出口、持续的贸易顺差、人民币被低估的估值（模型显示平均低估9.7\%）以及央行引导汇率走低的信号。这是典型的“本地因素”驱动——中国自身的经济基本面和政策意图，才是人民币走强的根本动力。

做空美元/新台币（目标30.5）：逻辑与人民币类似，但叠加了台湾地区在人工智能、半导体等科技领域的强劲增长和出口优势。台湾的贸易顺差扩大、GDP高增长预期，以及科技股带来的资本流入，都是支持新台币走强的本地因素。

做多欧元/印度卢比（目标113.8）：这是一个非常精妙的相对价值交易。NOM认为，在能源价格下跌的环境中，欧元

[... middle omitted ...]

：从单边押注转向相对价值，韩国与印度呈现不同逻辑

NOM在亚洲利率市场的策略同样体现了“从方向到相对价值”的转变。报告明确指出，在近期利率大幅下行后，市场焦点应转向相对价值交易。

韩国利率：NOM维持“接收韩国2年远期5年利率”的头寸，但将确信度从4/5下调至3/5。这表明，尽管韩国利率在和平协议下仍有下行空间（作为高beta市场，对地缘风险缓和最为敏感），但前期涨幅已部分消化了利好。报告隐含的逻辑是：韩国更强的财政状况意味着其利率进一步下行的空间可能受限，因此需要降低单边押注的力度，等待更好的入场时机。

印度利率：NOM推荐“做平印度2年与5年NDOIS利差”的交易（Sep-2s5s NDOIS flattener）。这是一个典型的曲线交易，而非方向性押注。其背后逻辑是：印度央行（RBI）可能通过买入长期国债来压低长端利率，同时市场对短期利率的预期可能因通胀和资本外流而保持坚挺。这种曲线平坦化交易，正是对印度本地货币政策、财政状况和资本流动复杂博弈的综合反映。

\subsection{[R003] 摩根斯坦利：存储器价格的结构性重置，远比市场理解的更深远}
{\small\color{refgray}归类：AI/算力与半导体：从GPU短缺到系统级重构，互联与存储定价权转移}

2026年，全球存储器市场规模将从2025年的约2200亿美元跃升至约8900亿美元。增量部分——约6000亿美元——已经超过了智能手机、PC或服务器任何一个单一品类的完整市场空间。

这不是一个周期性波动。这是存储器作为数字世界基础输入品，其价格行为正在经历一次结构性重置。

摩根斯坦利在最新发布的《Chipflation – Implications of a Memory Crisis》全球宏观论坛报告中，用一个核心概念概括了这一变化：Chipflation。芯片通胀。报告由首席固定收益策略师Vishwanath Tirupattur领衔，联合半导体、IT硬件、商品策略、美国公共政策和美国经济研究团队共同撰写。这不是一份单一的行业报告，而是一次跨资产、跨领域的系统性推演。

市场目前对这场存储器危机的定价，仍停留在“供应紧张”的周期叙事中。但报告揭示的真相是：AI正在从根本上改变存储器的需求结构和定价逻辑，而供给端的三家寡头格局和地缘政治约束，使得传统的“扩产-降价”循环已经失效。

以下是我们从这份报告中提炼出的五个核心洞察，以及一个报告尚未完全回答的关键问题。

过去一年，存储器价格上涨超过六倍。对于习惯了半导体价格长期下行的产业界和投资者来说，这个数字本身就是一个信号。

摩根斯坦利的分析框架清晰地指出，这不是一次普通的库存周期反转。报告用了一个关键区分：这是一个供给问题，还是一个需求问题？答案是两者兼具，但驱动因素完全不同。

从需求端看，AI创造了突然且价格弹性极低的需求跃升。HBM（高带宽存储器）的使用量，在单个AI芯片层面增长了7倍，在AI系统层面增长了65倍，在AI集群层面增长了惊人的1800倍。到2028年，HBM将占据领先制程存储器晶圆的约34\%，而2023年这一比例仅为约6\%。

从供给端看，全球DRAM市场约90\%的产出由三家厂商控制——三星、SK海力士和美光。这种高度集中的供给结构，在面临价格弹性极低的AI需求时，天然具有定价权。更关键的是，中国在全球存储器产能中的份额正从18\%上升至2028年的超过23\%，占全球净晶圆增量的约30\%。地缘政治因素使得这一产能扩张路径充满不确定性。

这两个力量合在一起，意味着存储器价格可能不会像过去那样在周期高点后迅速回落。报告的核心判断是：价格弹性较低的需求遇到了缺乏弹性的供给，这才是结构性重置的本质。

对于产业决策者而言，这意味着过去十年“存储器越来越便宜”的假设已经失效。任何依赖存储器作为核心成本项的终端产品——从服务器到PC到智能手机——都需要重新审视其成本结构和定价策略。

摩根斯坦利的分析揭示了一个被市场低估的连锁效应：当AI需求成为存储器厂商的优先分配对象时，消费级市场将出现结构性缺口。

报告给出的数字非常具体。在2027年的供应分配推演中，总DRAM+HBM供应为5760亿Gb，但扣除HBM分配后，剩余的非服务器供应仅为1730亿Gb。再扣除其他非服务器需求，最终分配给PC和智能手机的只有1300亿Gb，而这一市场的需求为1450亿Gb。这意味着约190亿Gb的缺口，相当于需求的13\%。

这个缺口的具体含义是什么？报告估算，这将导致2027年约15\%的PC DRAM短缺和约12\%的智能手机DRAM短缺，对应约5800万台PC和1.34亿部智能手机。

这不是一个“可能发生”的情景。这是基于当前产能规划、AI需求增长轨迹和厂商优先分配逻辑的基准预测。存储器厂商没有动力去优先满足利润率较低的消费级市场，尤其是在HBM和AI服务器DRAM的利润率远超传统产品的情况下。

这个判断的产业含义是：消费电子市场的竞争逻辑正在被改写。过去，终端厂商可以通过与存储器供应商的长期合作获得稳定供应。但在AI优先的分配逻辑下，消费电子品牌将面临

[... middle omitted ...]

26.9\%，上游中间投入品价格涨幅更为显著。但报告估算，存储器成本对2026年整体CPI的拉动仅为约10个基点。其中，PC和智能手机两个品类合计贡献约8个基点，游戏主机的贡献虽然单品类高达125个基点，但由于权重极低，对整体CPI的影响微乎其微。

这一差异的关键在于：CPI指数对电子产品进行了质量调整，而ASP（平均售价）系列数据没有。报告通过散点图展示了PC领域ASP与CPI之间的关系：ASP上涨30\%时，CPI仅上涨约10\%。在智能手机领域，两者的关联性甚至更弱。

这意味着什么？终端制造商正在承担大部分成本压力。他们要么通过产品升级（如增加内存、提升性能）来合理化涨价，要么直接压缩自身利润空间。报告暗示，消费端生产商正在面临艰难的利润决策。

对于投资者而言，这意味着需要关注两个层面的传导：第一，哪些终端品牌有能力将成本转嫁给消费者？第二，那些无法转嫁成本的品牌，其利润率将在未来几个季度面临显著压力。芯片通胀对CPI的温和影响，恰恰反映了企业利润被挤压的现实。

\subsection{[R004] Bernstein：韩国设备进口数据揭示的不是周期复苏，而是记忆体结构性扩张的加速}
{\small\color{refgray}归类：AI/算力与半导体：从GPU短缺到系统级重构，互联与存储定价权转移}

Bernstein：韩国设备进口数据揭示的不是周期复苏，而是记忆体结构性扩张的加速

半导体设备进口数据，通常被市场视为周期风向标。当韩国5月半导体设备进口同比飙升51\%时，许多投资者的第一反应是“周期回来了”。但这份Bernstein研报揭示的图景远比“周期复苏”复杂——它指向的是记忆体领域一次由技术节点跃迁驱动的结构性资本支出扩张，而非简单的需求回暖。理解这一区别，对于判断设备供应链上各家公司的相对受益程度至关重要。

市场之所以容易误读，是因为传统框架下进口激增往往与下游需求旺盛挂钩。然而，Bernstein团队通过拆解进口来源国别和设备品类，发现本轮增长的核心驱动力来自韩国两大记忆体巨头——三星和SK海力士——对1c纳米DRAM节点的激进导入，以及由此带来的光刻强度（litho intensity）显著提升。这不是一个“量”的故事，而是一个“质”的故事。

这份报告最值得关注的判断是：韩国设备进口的强劲势头，正在从总量层面验证记忆体资本开支的结构性上调，但不同设备供应商之间将出现显著分化。光刻设备供应商ASML有望持续受益于高价值量系统的密集交付，测试设备供应商Advantest可能迎来超预期的季度收入，而日本设备商东京电子（TEL）的收入动能则显示出边际放缓的迹象。这种分化，恰恰是投资者需要重新校准预期的地方。

1. 光刻设备进口创下季度内第二高纪录，ASML在韩收入动能持续强化

韩国从荷兰进口的半导体设备——几乎可以视为ASML光刻系统的代理变量——在5月达到9.28亿欧元，这是有记录以来季度内第二高的单月数字，环比增长28\%，同比增幅约150\%。Bernstein的回归模型据此估算，ASML第二季度在韩国的系统销售额约为23.1亿欧元，虽然环比下降18\%，但同比翻倍有余。

这意味着什么？韩国将在ASML第二季度系统销售总额中占据约37\%的份额。对于一个季度前市场还在担忧记忆体资本开支放缓的行业而言，这一比例本身就说明了很多问题。更关键的是，Bernstein指出，这一强劲势头很可能得到了1c节点DRAM产能扩张和快速采用的支撑——1c节点相比前代产品需要显著更高的光刻强度。

从投资含义来看，ASML在记忆体领域的收入集中度正在上升，这既是机遇也是风险。机遇在于，记忆体客户的资本开支计划通常更为集中和可预测，一旦进入交付周期，收入能见度较高；风险在于，如果下游记忆体价格出现超预期下跌，客户可能临时调整交付节奏。但就当前数据而言，ASML在韩国的收入动能尚未见顶。

2. 测试设备进口同比翻倍，Advantest季度收入可能显著超出市场预期

测试设备是这份报告中另一个值得深挖的信号。韩国从日本和马来西亚进口的测试设备在5月环比增长5\%，同比增幅高达103\%。Bernstein指出，Advantest的记忆体测试收入与韩国测试设备进口数据之间存在良好的相关性。

基于回归模型，Bernstein估算Advantest第二季度在韩国的测试收入可能环比增长84\%。作为对比，市场共识预期Advantest整体公司收入仅环比增长3\%。即便Bernstein也谨慎地指出这一回归仅基于两个月的数据且近期相关性有所减弱，但如此巨大的差距仍然值得投资者高度关注。

这里的核心逻辑在于：记忆体客户在导入新节点时，测试环节的需求往往滞后于光刻和刻蚀等前道设备。如果1c节点的产能爬坡正在加速，那么测试设备的需求释放可能才刚刚开始。Advantest的韩国收入如果确实如模型所示出现井喷，将意味着市场对其全年业绩的预期存在系统性低估。

与ASML和Advantest形成对比的是，东京电子相关品类（CVD、干法刻蚀、清洗、涂胶显影、RTP等）的韩国进口数据在5月环比下降27\%，尽管同比仍增长52\%。Bernstein的

[... middle omitted ...]

着东京电子的长期竞争力受损。但如果市场将韩国设备进口总量增长简单线性外推到所有设备供应商，东京电子的季度业绩就可能成为预期差来源。这一点，对于持有或关注日本半导体设备板块的投资者尤为重要。

4. 韩国记忆体资本开支与进口数据出现短期背离，但下半年有望加速修复

一个值得注意的细节是：尽管设备进口数据持续攀升，三星和SK海力士第一季度的合并资本开支却环比下降。Bernstein将其归因于季节性因素以及部分基础设施投资在2025年第四季度前置。

这种短期背离并不罕见，但关键在于后续走势。两家公司均已明确指引2026年资本开支将大幅增长，重点投向基础设施和战略设备。Bernstein据此判断，未来资本开支将恢复增长，甚至可能很快超过当前的进口数据。

这意味着什么？当前设备进口的强劲势头并非昙花一现，而是更大规模资本开支周期的前奏。如果三星和SK海力士下半年启动新一轮设备采购，进口数据有望在现有高位基础上进一步上移。这对整个半导体设备板块而言，是一个中期层面的积极信号。

\subsection{[R005] Bernstein：美伊和平对印度而言是反弹，不是反转}
{\small\color{refgray}归类：地缘风险与汇率：和平红利下的货币重估与亚洲资产重新定价}

印度市场正在庆祝一个期待已久的信号：美伊之间可能达成真正的和平协议，霍尔木兹海峡重新开放。油价应声回落至83美元，航空、石油营销公司（OMC）、医疗保健乃至IT板块的情绪迅速升温。但这份Bernstein研报的核心判断是——市场应当区分“乌云散去”和“阳光普照”是两回事。反弹是确定的，但反转的条件远未成熟。真正值得深挖的，不是短期的情绪修复，而是和平协议可能带来的两个结构性变量：伊朗原油制裁豁免，以及恰巴哈尔港的战略价值。这两点，才是决定印度能否从“反弹”走向“重估”的关键。

这份研报发布于一个微妙的时点。此前三个月，中东的和平声明屡屡被证明是短暂的，市场已对“消息面反弹”产生疲劳。但Bernstein认为这次不同——美伊双方官方渠道同步确认，且得到调解方背书，协议的可靠性高于以往。然而，研报同时警告，魔鬼在细节中。双方出于国内政治需要，可能在短期内各自叙事，真正的条款细节需要数周才能明朗。这意味着，当前市场的定价更多基于情绪，而非事实。

对于投资者而言，核心问题不是“要不要参与反弹”，而是“反弹的持续性和深度由什么决定”。Bernstein的答案很明确：真正的催化剂不是海峡重开，而是伊朗原油回归和地缘战略通道的打通。而这两点，恰恰是当前市场讨论最少、定价最不充分的变量。

1. 霍尔木兹重开对油价的缓解可能低于预期，因为需求侧会同步反弹

油价从高位回落至83美元，市场的第一反应是成本端压力缓解。但Bernstein指出，霍尔木兹海峡重开带来的供给增量，可能很快被需求端的同步复苏所抵消。全球经济活动在战争阴云消散后有望加速，尤其是航空和工业用油需求将回升。因此，油价更可能在一个较窄的区间内波动，而非趋势性下行。报告明确判断，油价重回三位数的概率已经大幅下降，但同样难以跌至60-70美元的低位区间。这一判断意味着，依赖低油价的板块（如航空、石化）的利好是有限的、阶段性的，而非结构性的。

对于印度而言，油价的“不涨”比“大跌”更重要。印度作为原油净进口国，油价稳定在80-90美元区间，已经足以让财政赤字目标不再恶化。但如果油价只是温和回落，则无法为印度提供额外的财政空间。

2. 印度市场的反弹是“减负”而非“加力”，估值和结构性逆风并未改变

Bernstein在6月初调整行业配置时，已经将中东停火作为前提假设。因此，当前的市场反弹更多是此前策略的兑现，而非新的超额收益来源。报告重申了对Nifty年底26000点的目标价，并维持中性立场。这一目标价意味着，当前市场已经部分定价了和平利好，剩余空间有限。

更关键的是，印度市场面临的结构性逆风并未因和平协议而消失。估值仍然偏贵，印度在全球AI产业链中处于不利位置，原油均价仍将高于去年，近期燃料和LPG的涨价将推高通胀读数——而去年同期的通胀基数较低，这意味着同比数据可能超预期。此外，印度气象局已确认厄尔尼诺现象正在发展，雨季的不确定性尚未解除。这些因素叠加，使得印度央行在可预见的未来既不可能加息，也不可能降息。利率环境对成长股估值的压制将持续存在。

因此，反弹的逻辑是“风险溢价下降”，而非“盈利预期上调”。投资者需要警惕的是，当情绪修复完成后，市场将重新面对这些未解决的宏观问题。

3. 伊朗原油制裁豁免是比海峡重开更大的催化剂，但存在执行风险

报告中最具洞察力的部分，是对伊朗原油制裁豁免的分析。伊朗方面声称，和平协议包括对其原油出口的制裁豁免。如果这一条款落地，伊朗每日300-400万桶的产能可能部分回归市场。Bernstein测算，这将推动油价跌破80美元。尽管仍难以回到60-70美元区间，但“低于80美元”对于印度的财政和经常账户而言，是一个质的临界点。

报告附带的财政敏感性分析表清晰地展示了这一临界点的重要性：当油价维持在90美元时，印度FY2

[... middle omitted ...]

和巴基斯坦港口的替代通道，这对于能源安全、中亚资源获取以及对抗中国“一带一路”在中亚的影响力，都具有深远意义。对于相关的基础设施建设、港口运营和物流企业而言，这可能是一个持续数年的结构性主题，而非一次性的消息面炒作。

但报告也承认，这一前景高度依赖美伊关系的长期稳定以及伊朗国内的政治走向。目前来看，恰巴哈尔港的商业化运营仍处于早期阶段，其财务影响难以量化。

5. 报告尚未完全回答的关键问题：反弹能否转化为结构性重估？

Bernstein的分析框架清晰而严谨，但有一个问题悬而未决：如果和平协议能够持续，且伊朗原油和恰巴哈尔港的利好逐步兑现，印度市场是否应该迎来一次全面的重估？目前，报告给出的答案是否定的——它坚持中性立场和26000点的目标。但这一判断建立在“原油均价仍高于去年”和“估值偏高”的假设之上。如果伊朗原油回归超预期，油价跌破75美元，同时恰巴哈尔港带动新一轮基础设施投资和贸易增长，那么盈利预期和估值中枢都可能上移。届时，26000点的目标可能显得保守。

\subsection{[R006] GS：5月中国外汇流入的真实信号被市场低估了}
{\small\color{refgray}归类：地缘风险与汇率：和平红利下的货币重估与亚洲资产重新定价}

市场对中国资本流动的讨论，长期停留在“流出压力”的叙事框架里。但GS最新发布的5月外汇数据解读，给出了一个截然不同的信号：中国在5月录得了约500亿美元的净外汇流入，这不仅是连续第二个月净流入，而且结构上发生了值得关注的转变。

这份报告的核心价值，不在于告诉你“流入还是流出”这个结论，而在于它拆解了流入的构成——经常项目依然强劲，但真正值得关注的，是资本金融项目下出现了一些此前数月未曾见到的边际变化。这些变化能否持续，将直接决定人民币资产定价的底层逻辑是否正在被重新书写。

1. 5月净流入的核心驱动力仍是贸易顺差，但“结汇意愿”出现了微妙信号

GS偏好的外汇流量指标显示，5月净流入500亿美元，高于4月的370亿美元。其中，经常项目贡献了370亿美元的净流入，货物贸易相关的净流入从4月的470亿美元升至530亿美元。

表面上看，这延续了“出口强劲、顺差扩大”的老故事。但GS报告里一个容易被忽视的数字是：5月的结汇比率降至51\%，低于4月的56\%，更远低于一季度74\%的平均水平。

这意味着什么？企业赚了更多外汇，但选择“持有”而不是“兑换”人民币的比例在上升。从直觉上看，这似乎是一个负面信号——企业不结汇，说明对人民币汇率缺乏信心。

但这里需要做两层思考。第一，51\%的结汇比率仍然高于 年多数月份的水平，说明企业并非在系统性抛售人民币。第二，更重要的是，结汇比率下降的同时，整体外汇流入却在增加，这意味着“不结汇”的量级本身也在扩大——企业手中积累的外汇头寸正在变厚。一旦汇率预期出现拐点，这些“蓄水池”里的外汇可能集中释放，形成人民币升值的脉冲式力量。

这个结构性问题，报告没有完全展开，但它对汇率定价的隐含含义值得每一个持有人民币资产的投资者思考。

2. 债券市场正在经历一个被忽视的拐点：外资重新净买入中国国债

这是整份报告中最具“信号意义”的变化。GS数据显示，5月北向债券通录得130亿美元的净流入，而4月还是100亿美元的净流出。这是自2025年4月以来，外资首次净买入中国在岸债券。

驱动因素非常集中：中国国债（CGB）。报告明确指出，净买入主要由CGB驱动，而政策性银行债、同业存单等其他品种并未出现类似的资金回流。

这个信号的解读需要谨慎。外资买中国国债，可能出于几个不同的动机：利率套利（中美利差变化）、汇率对冲后的正收益、或者单纯的资产配置再平衡。GS报告没有给出外资买入的详细动机拆解，但我们可以从数据中做合理推断。

如果外资是在赌人民币升值，那么买入国债只是“搭便车”行为，一旦汇率预期逆转，资金也会迅速撤离。但如果外资是基于中国国债在全球利率体系中的“独立定价”价值而买入，那么这就是一个更持久的配置型力量。

目前来看，后者的可能性在上升。在全球主要央行同步降息或维持高利率的背景下，中国国债收益率曲线的“独立性”正在吸引那些追求分散化收益的长期资金。这一点，报告没有说透，但它是决定债券流入能否持续的关键假设。

3. 官方储备和商业银行资产的变化，揭示了一个更完整的资金循环

GS报告还提供了一个非常有趣的视角：将官方外汇储备和商业银行的净外部资产变化放在一起看。

5月，官方外汇储备经估值调整后增加了390亿美元，而商业银行的净外部资产增加了150亿美元。两者合计，5月中国整体对外资产增加了约540亿美元。

这个数字与GS偏好的外汇流量指标（500亿美元）基本吻合，说明数据内部一致性较好。但更值得关注的是结构：官方储备和商业银行资产都在增加，但增速不同。

官方储备的增加，意味着央行在某种程度上在“吸收”外汇流入，这有助于维持人民币汇率的稳定性。而商业银行净外部资产的增加，则意味着部分外汇流入了银行体系，而非全部被央行“截留”。这为商业银行提供了更多的外汇头寸，可能在未来的跨境

[... middle omitted ...]

债券、房地产）。5月的数据显示，经常项目顺差依然庞大，但结汇比率下降；债券市场出现了边际改善，但仅限于国债；股票市场的外资流入数据，报告并未提及。

这意味着，外汇流入的“质量”可能不如总量看起来那么乐观。如果流入的外汇只是停留在银行体系或企业账户上，而没有真正进入人民币资产市场，那么它对资产价格的支撑作用就是有限的。

但反过来看，这也意味着一旦条件成熟（如汇率预期稳定、国内政策明朗、风险偏好回升），这些“沉睡”的外汇可能成为推动资产价格上涨的增量资金。关键在于，企业——尤其是出口企业——是否愿意将手中的外汇转化为人民币资产。

GS这份报告提供了高质量的数据和清晰的框架，但对于几个关键问题，它没有给出明确的答案，或者说，这些答案需要更多数据来验证。

第一，外资买国债的可持续性。5月的130亿美元净流入，是单月脉冲，还是趋势反转的开始？这取决于全球利率环境的变化和中国国债的相对吸引力。如果美国经济数据继续超预期，美债收益率再次上行，外资买中国国债的动力可能会减弱。

\subsection{[R007] JPM：市场真正低估的不是增长，而是信用脉冲对风格切换的定价权}
{\small\color{refgray}归类：未被主题摘要引用}

JPM：市场真正低估的不是增长，而是信用脉冲对风格切换的定价权

上周五，美联储与中国人民银行几乎同时发布了两组关键数据——美国一季度资金流量表与中国5月社融数据。这不是一个巧合，而是一个信号。对于大多数投资者而言，这两组数据各自指向不同的方向：美国私人部门信用扩张加速，中国社融连续第三个月放缓。但JPM量化策略团队在其最新报告中提出了一个更值得深思的判断：合在一起看，这两组数据共同指向一个方向——Growth风格相对于Value风格仍有空间，但信心的基础正在变薄。

这不是一篇关于宏观预测的报告。这是一篇关于“如何用信用脉冲作为信号，对风格切换进行定价”的报告。JPM构建了一个简单但有效的框架：将美国与中国信用脉冲信号取平均值，作为判断Growth与Value相对表现的领先指标。自2018年以来的回测显示，这一信号的历史命中率为53\%，年化超额收益达到9.9\%。目前，该信号仍指向超配Growth，但信号强度已显著低于三个月前。

这个判断的核心在于：市场正在定价的，可能不是增长本身，而是信用创造的速度与方向。而这一框架的真正价值，不在于它给出了一个明确的交易信号，而在于它揭示了市场定价中一个被广泛忽视的结构性变量。

大多数投资者关注信用数据时，着眼的是经济增长的领先指标。JPM提供了一个不同的视角：信用脉冲是风险定价的代理变量，尤其对Growth风格的资产定价高度敏感。

这一逻辑链条值得拆解。Growth资产的定价高度依赖远期的现金流折现，而折现率中的风险溢价部分对宏观流动性环境极为敏感。信用脉冲，即私人部门信用创造的边际变化，恰恰是捕捉这种风险溢价波动的有效工具。当信用扩张加速时，风险溢价下降，Growth资产相对受益；当信用收缩时，风险溢价上升，Value资产相对抗跌。

JPM报告中的两张图表清晰地展示了这一关系。美国信用脉冲与MSCI中国指数中Growth相对于Value的回报差之间存在高度相关性。同样，中国信用脉冲自2018年A股纳入MSCI以来，也呈现出类似的同步关系。

这里的关键洞察是：信用脉冲之所以有效，不是因为它预测了经济增长，而是因为它捕捉了市场中“风险偏好”的边际变化。而风险偏好的变化，往往先于基本面数据的变化。

JPM报告中最引人注目的数据点，是中美信用脉冲的背离。美国一季度信用脉冲从+1.6个百分点上升至+2.5个百分点，继续处于扩张区间。而中国5月社融数据则标志着连续第三个月的放缓，中国信用脉冲已滑落至-0.4个百分点。

这一背离导致了一个关键结果：中美平均信用脉冲虽然仍在上升，但仅略高于其六个月移动平均线。信号的方向没有改变，但信号的力度明显减弱。

这意味着什么？从量化策略的角度看，信号强度的减弱意味着交易信心的下降。从产业决策的角度看，这意味着市场正在进入一个“分裂”的状态——美国信用扩张对Growth的支撑，正在被中国信用收缩的拖累所抵消。对于全球配置的投资者而言，这一背离意味着不能简单地押注单一方向，而需要更精细地拆解不同市场、不同板块对信用脉冲的敏感度差异。

这一背离背后，隐藏着一个更深层的问题：中美信用周期的不同步，是否正在成为一种结构性特征，而非周期性波动？如果答案是肯定的，那么基于历史回测构建的信用脉冲信号，其有效性可能需要重新校准。

JPM报告披露了其信用脉冲信号的历史表现：自2018年以来，命中率53\%，年化超额收益9.9\%。对于任何一个量化策略而言，这都不是一个可以忽略的数字。但真正值得思考的问题是：这个信号在什么时候会失效？

53\%的命中率意味着，在将近一半的时间里，信号指向了错误的方向。JPM没有在报告中详细展开这些“失败案例”的共同特征，但这恰恰是投资者需要自行推演的关键。

从逻辑上推断，信用脉冲信号可能在以下几种情况下失效：第一，当

[... middle omitted ...]

来，取决于信用脉冲的下一步方向。

4. 报告尚未完全回答的关键问题：信用脉冲的边际变化如何与估值水平交互

JPM的报告提供了一个清晰的框架，但有一个关键问题没有被充分讨论：信用脉冲信号的有效性，是否取决于市场的估值水平？

逻辑上，当Growth资产的估值处于极端高位时，即使信用脉冲信号指向Growth，其超额收益的空间也可能被估值压缩。反之，当Value资产的估值处于极端低位时，即使信用脉冲信号指向Value，Value的反弹也可能被基本面恶化所抵消。

这一问题在当前市场环境下尤为重要。Growth资产的估值在经历了过去几年的扩张后，已经处于历史较高分位。如果信用脉冲信号的有效性在估值极端时下降，那么当前“仍超配Growth”的信号就需要被更谨慎地解读。

JPM的报告没有提供这一维度的分析。这并非报告的缺陷，而是任何量化框架都无法完全覆盖的边界。对于投资者而言，这意味着不能将信用脉冲信号作为唯一的决策依据，而需要将其与估值、盈利增长、市场情绪等变量结合使用。

\subsection{[R008] 摩根斯坦利：中国权益市场正在被低估的并非需求，而是供给侧的结构性定价能力}
{\small\color{refgray}归类：中国权益与地产：结构性供给优势被低估，城市更新资金落地成关键变量}

摩根斯坦利：中国权益市场正在被低估的并非需求，而是供给侧的结构性定价能力

这份报告发布的时间点值得玩味。2026年6月，MSCI中国年初至今回报为-10\%，而MSCI新兴市场同期上涨18\%。表面看，中国资产似乎又一次跑输了。但摩根斯坦利亚洲策略团队在题为“Forging New Horizons”的投资者报告中提出了一个与市场情绪截然不同的判断：中国市场的核心叙事正在从“周期性复苏”切换到“结构性供给优势”，而这一切换尚未被全球定价体系充分反映。

报告的核心主张不是中国需求有多强，而是中国在全球AI和能源转型资本开支超级周期中的供给侧地位，正在系统性地重塑其出口竞争力和企业盈利结构。这意味着，投资者如果继续用过去三年“弱复苏、强政策依赖”的框架来看待中国权益，可能会错过一个更根本的变量——中国正在成为全球资本品和电子供应链中不可替代的“供给锚点”。

这份报告提供的新信号在于：中国在全球出口市场的份额不仅没有因供应链多元化而下滑，反而有望在2030年达到16.5\%甚至更高。而这一扩张的驱动力，并非传统的低成本制造，而是深度嵌入全球AI基础设施和能源转型的供给能力。以下是我们从报告中提炼出的五层核心洞察，以及一个尚未被完全回答的关键问题。

1. 指数表现掩盖了结构性分化：CSI 300的风险调整后收益才是真正的信号

报告开篇用一组数据直接挑战了“中国资产表现不佳”的流行叙事。从2024年9月至今，MSCI中国总回报为28\%，MSCI新兴市场为49\%，S\&P 500为30\%。单看绝对回报，中国确实落后。但摩根斯坦利引入了更精细的衡量维度——夏普比率。

2025年全年，CSI 300的夏普比率为1.44，高于MSCI EM的1.70吗？不，MSCI EM更高。但请注意，CSI 300的年度化标准差仅为15\%，远低于MSCI中国的24\%和MSCI EM的15\%——实际上CSI 300的波动率更低。这意味着，在2025年，以人民币计价的A股核心资产在风险调整后的回报表现，实际上优于海外投资者习惯使用的MSCI中国框架所呈现的画面。2026年至今，CSI 300的夏普比率为0.68，而MSCI中国为-1.34。这组数据的含义很清楚：A股和港股之间的分化，不仅仅是市场结构问题，更反映了不同投资者群体对同一套基本面事实的定价差异。

对于产业决策者而言，这意味着“中国资产”不再是单一类别。A股核心资产（以CSI 300为代表）正在展现出更优的风险收益特征，而这一特征主要来自低波动率而非高回报。这反过来指向一个判断：中国本土机构投资者对国内经济基本面的定价，可能比海外投资者更稳定、更前瞻。

2. 指数权重结构是理解“跑输”的关键：通信服务和软件板块的拖累是周期性的，而非结构性的

报告进一步拆解了MSCI中国年初至今表现不佳的行业归因。最引人注目的是两个板块：通信服务（权重18.5\%）和软件与服务（权重1\%），年初至今回报分别为-35\%和-35\%。这两个板块合计权重接近20\%，对指数形成了巨大拖累。而同期，半导体（权重2.5\%，若按proforma权重则为7\%）回报为正7.5\%，资本品（权重4\%，proforma权重9\%）回报为正2\%。

这意味着，MSCI中国指数层面的疲弱，很大程度上是由少数几个受到监管调整或行业周期影响的板块驱动的。如果剔除通信服务和软件板块，指数的实际表现会好得多。更关键的是，半导体和资本品板块的proforma权重显著高于实际权重，说明这些结构性受益板块在指数中的代表性不足，其真实增长潜力被指数构成所掩盖。

对于投资者而言，这提出了一个操作层面的问题：如果继续以MSCI中国作为中国权益的基准，可能会系统性地低配那些真正受益于全球AI和能源转型资本开支的板块。报告隐含的建议是，行业选择和因子

[... middle omitted ...]

，其中约7000亿美元来自AI相关需求。中国资本品公司正在全球AI基础设施中获取份额，尤其是在电力设备、冷却系统、精密制造等领域。第二个是能源转型。中东冲突加速了全球对可再生能源和电力设备的需求，而中国在太阳能制造的关键环节控制着超过80\%的产能。

这两个驱动力共同指向一个结论：中国出口的竞争优势正在从“成本优势”转向“供应链深度和规模不可替代性”。报告用“China has been widening its lead in global export since 2022”来描述这一趋势，并指出中国在电子集成电路和“新三样”（电动车、锂电池、太阳能电池）领域的出口增速正在加速。2026年初至今，电子集成电路出口增速已接近80\%，“新三样”出口增速约60\%，远高于其他品类。

这意味着，对于全球制造业企业而言，中国不再仅仅是“低成本制造基地”，而是全球AI和能源转型供应链中不可绕过的“供给锚点”。任何试图完全“去中国化”的供应链重组，都将面临巨大的成本和时间代价。

\subsection{[R009] 摩根斯坦利：贸易韧性掩盖了国内需求的深层裂痕}
{\small\color{refgray}归类：中国权益与地产：结构性供给优势被低估，城市更新资金落地成关键变量}

这份摩根斯坦利在2026年亚洲夏季研讨会上发布的报告，标题本身就浓缩了一个核心矛盾：“更强劲的贸易，更疲弱的国内需求”。这不是一句中性的描述，而是一个正在撕裂中国宏观叙事的张力结构。对于关注中国资产定价的读者而言，真正重要的不是出口数据有多亮眼，而是当外需的强风暂时掩盖了内需的病灶时，资产价格的风险定价是否已经足够充分。

报告提供的最新信号是：2026年5月的数据显示，出口在AI周期和全球科技产业链重构的推动下出现了罕见的广度改善，电子集成电路、计算机、手机等品类同比增速惊人。但与此同时，私人信贷需求却在全面恶化——5月新增短期居民贷款为负70亿元，新增中长期居民贷款同样为负50亿元，均远低于过去五年的同期均值。这两个方向截然相反的指标同时出现，意味着中国经济正处于一个极其特殊的分化阶段：外循环的动能正在加速，而内循环的收缩尚未触底。

这份报告的价值，不在于它重复了“出口好、内需弱”的共识，而在于它用数据揭示了这种分化的结构细节、政策应对的局限性，以及中长期再平衡所需要的改革前提。以下是我们从这份报告中提炼出的五个核心洞察。

1. 出口的改善不是简单的周期性回暖，而是科技产业链重构带来的结构性红利

报告中的一张图清晰展示了出口改善的广度：2026年4月至5月，中国出口增速在多个品类上出现跃升。电子集成电路的同比增速从4月的约100\%进一步攀升至5月的约110\%，计算机从45\%升至65\%，手机从10\%升至45\%。即便是此前表现疲软的品类，如精炼石油产品、塑料制品，也在5月转正或加速。

但真正值得关注的不是这些数字本身，而是它们背后的驱动力。报告明确指出，全球AI周期正在驱动中国贸易增长。这意味着当前的出口改善并非传统的“外需回暖”，而是全球科技投资周期与中国制造业供应链优势的深度耦合。中国在电子制造、数据中心设备、智能电网等领域的产能，恰好契合了全球AI基础设施建设的高峰需求。

这一判断的涵义在于：出口的韧性可能比市场预期的更持久，因为它建立在资本开支周期而非库存周期之上。但同时，它也让中国经济面临一个更尖锐的结构性问题——贸易部门的繁荣能否通过就业、工资和税收渠道，有效传导至内需部门？从报告后续的数据看，这个传导机制似乎并不顺畅。

报告提供的5月信贷数据是这份报告中最值得反复研读的部分。新增短期居民贷款为负70亿元，新增中长期居民贷款为负50亿元，新增中长期企业贷款加债券融资仅为145亿元，而过去五年的同期均值是535亿元。这三个数字同时为负或大幅低于均值，在中国信贷数据中并不常见。

更关键的是，报告将这一现象与近期房地产销售的“反弹”放在一起审视。报告的原话是：“持续的居民部门去杠杆表明，近期的住房销售反弹是狭窄的。”这意味着，少数城市或少数项目的成交回暖，并不能代表整个房地产市场的企稳。居民的资产负债表修复意愿仍然强烈，提前还贷、压缩消费信贷的趋势没有逆转。

对资产定价而言，这意味着：当前市场对消费复苏的定价可能过于乐观。如果居民部门仍在主动去杠杆，那么消费股的估值扩张缺乏微观基础。任何基于“消费升级”或“内需回暖”的投资叙事，都需要先回答一个问题：居民的钱从哪里来？

3. 财政政策的节奏正在放缓，但方向已从总量刺激转向“六网”基建

报告显示，2026年前五个月政府债券的发行进度约为41\%，低于2025年同期的46\%。这意味着财政扩张的节奏不仅没有加速，反而有所放缓。但报告同时指出，这并不代表政策转向紧缩，而是意味着从三季度开始，预算执行可能加速，重点将放在“六网”建设中的AI计算网络、互联网数据中心和智能电网。

“六网”是报告提出的一个政策框架，涵盖水网、物流网、新型电力网、计算网、城市地下管网和5G网络。报告估算， 年期间，仅计算网和数据中心相关的投资规模就将超过5万亿元，新型电

[... middle omitted ...]

速最快的行业（如新能源、电动车产业链），在 年投资增速明显回落；而此前投资增速较低的行业（如AI基础设施、数据中心），则开始加速。投资的热点并没有消失，只是换了赛道。

报告坦率地指出了“全国统一大市场”和“反内卷”政策面临的三大挑战：顶层协调难以执行，地方激励仍不匹配，税收体制偏向生产端而非消费端。这意味着，即便政策意图是减少无效竞争和产能过剩，但在现有的央地关系和税收体制下，地方政府仍有强烈的动机去补贴本地企业、扩张产能、争夺投资。供需失衡的纠正，可能需要更长时间，也需要更深层次的体制改革。

报告的最后一个重要板块讨论了资本流动和人民币汇率。报告指出，中国的国际收支正在从独特的“双顺差”模式，转向更常见的“镜像”模式——经常账户顺差与资本金融账户逆差并存。而且，资本流出的结构正在发生变化：对外直接投资（FDI）的占比下降，而证券投资（portfolio investment）的占比上升。这意味着，资本流出的驱动因素从企业海外建厂，转向了居民和机构的海外资产配置。

\subsection{[R010] 摩根斯坦利：ABF载板供需拐点已至，真正稀缺的不是产能，是客户锁定的产能}
{\small\color{refgray}归类：AI/算力与半导体：从GPU短缺到系统级重构，互联与存储定价权转移}

摩根斯坦利：ABF载板供需拐点已至，真正稀缺的不是产能，是客户锁定的产能

当一家欧洲载板厂商宣布将资本支出从4亿欧元大幅上调至10-12亿欧元，同时EBITDA利润率指引从25-29\%直接跳升至32-37\%，市场的第一反应往往是“供给要过剩了”。但摩根斯坦利这份关于AT\&S马来西亚扩产的研报给出了截然相反的判断：这不是供给过剩的信号，而是供给结构性短缺提前到来的确认。

这份研报最值得关注的洞察，不在于AT\&S上调了业绩指引，而在于一个关键细节——AT\&S追加的15-20亿欧元投资，是在现有厂区增加生产线，而非新建工厂。这意味着新产能将在2027年3月前落地。而摩根斯坦利此前判断ABF载板将从2027年起进入供给短缺周期。两件事放在一起看，结论清晰：短缺不是被推迟了，而是被提前证实了。

市场正在犯一个典型错误——把“产能扩张”等同于“供给过剩”，却忽略了这次扩张的底层逻辑已经改变：客户承诺前置、资本开支锁定、产能被预购。这不再是 年的产能竞赛，而是AI/HPC需求驱动下的供给再定价。

AT\&S此次扩产的资本开支高达15-20亿欧元，几乎是此前指引的三倍。但真正让这份报告值得细读的，不是金额本身，而是资金来源——AT\&S明确表示，这笔投资“fully supported and financed by long-term customer commitments”。

这意味着什么？在传统载板行业，产能扩张通常是厂商基于需求预测做出的风险决策。景气时扩产，低谷时承担闲置成本。但这次，客户——AMD和另一家大型科技公司——直接为产能买单。这不是“厂商赌需求”，而是“客户锁产能”。

这种模式的变化对行业竞争格局有深远影响。当产能被客户预购，剩余可流通的现货产能就变得稀缺。对于未锁定产能的下游厂商，未来拿货的难度和成本都将上升。而对于像欣兴电子（Unimicron）、南亚电路板（NYPCB）和臻鼎（ZDT）这样的头部厂商，它们如果也能复制类似的客户锁定模式，将获得更强的定价权和盈利稳定性。

所以，研报的核心判断不是“供给增加”，而是“可流通供给减少”。这个差异，决定了投资者应该如何看待整个ABF载板板块的估值逻辑。

2. 2027年的供给短缺不再是预测，而是已经被AT\&S的行动确认

摩根斯坦利此前对ABF载板的供需模型预测是：2027年起出现供给缺口。AT\&S此次扩产的时间线——新产线在2027年3月前上线——与这一预测高度吻合。

这里的关键推理是：如果AT\&S认为现有产能足以满足未来需求，它不需要在当前时点就启动15-20亿欧元的扩产计划。更合理的解释是，AT\&S已经看到了来自AI/HPC客户的订单强度，现有产能根本不够用。

研报中有一个隐晦但重要的判断：“This implies strong demand and does not derail our expectation of an ABF supply deficit from 2027 onwards.” 这句话的潜台词是：即使AT\&S如此大规模扩产，仍然无法扭转供给短缺的趋势。换句话说，需求端的强度超出了市场此前最乐观的估计。

对于投资者而言，这意味着ABF载板行业的景气周期可能比市场预期的更长、更陡。那些被市场贴上“周期股”标签的载板厂商，其盈利中枢正在发生结构性上移。

3. 欣兴电子、南亚电路板、臻鼎：三家公司的差异化机会在哪里

摩根斯坦利对这三家公司均给予“Overweight”评级，但各自的驱动逻辑并不相同。

欣兴电子的估值模型采用剩余收益法，假设股权成本9.2\%、中期增长率15\%、终值增长率3\%。其风险因素中，上行风险包括PC和服务器客户需求超预期、资本支出削减或扩产暂停、以及替代技术（如CoWoS）持续存在良率问题

[... middle omitted ...]

。

三家公司的共同点在于：它们都处于ABF载板供需拐点的受益方，但各自的催化剂和风险敞口不同。投资者需要根据自身对AI需求、PC复苏、iPhone周期以及替代技术演进的不同判断，做出差异化配置。

研报在风险披露中提到了一个值得警惕的变量：替代技术。具体来说，如果CoWoS等先进封装技术持续推进，部分原本需要ABF载板的芯片可能转向无载板方案，这将从根本上改变ABF的需求结构。

这份报告对此着墨不多，只是将其列为下行风险之一。但这个问题恰恰是当前市场分歧最大的地方。一方面，台积电的CoWoS产能正在快速扩张，先进封装对传统载板的替代效应正在显现。另一方面，AI/HPC芯片的复杂度仍在提升，对ABF载板的层数、尺寸和精度要求也在同步提高。两种力量相互抵消，最终的净影响取决于技术演进的相对速度。

从投资角度看，这个不确定性意味着：ABF载板板块的上行空间可能很大，但也存在一个尚未被充分定价的技术风险。那些在替代技术领域也有布局的载板厂商，可能具备更强的风险对冲能力。

\subsection{[R011] 美国银行：市场真正低估的不是AI泡沫，而是通胀二次探底对仓位结构的重塑}
{\small\color{refgray}归类：宏观与利率：财政供给与通胀再平衡主导定价}

美国银行：市场真正低估的不是AI泡沫，而是通胀二次探底对仓位结构的重塑

这份报告的核心判断，不是关于AI是否过热，而是关于一个更根本的、正在发生的结构性转变：全球基金经理正在从一个“拥抱风险、追逐增长”的仓位状态，转向一个“警惕通胀、防御性调整”的新阶段。美国银行（BofA）2026年6月的全球基金经理调查（FMS）显示，尽管市场情绪依然偏多，但最关键的信号并非AI的拥挤度创下历史新高，而是投资者对利率预期的剧烈转向，以及由此引发的仓位再平衡。这意味着，未来几个月的资产定价逻辑，可能将从“增长叙事”切换为“利率与通胀叙事”。

为什么现在重要？因为市场正处于一个微妙的“共识分歧”点。一方面，宏观情绪在改善，全球增长与利润预期均创三个月新高；另一方面，对利率的预期却急转直下，认为未来12个月美联储将加息的基金经理占比从16\%飙升至40\%，为近四年最高。这种“增长向好、利率向上”的组合，在历史上往往不是风险资产的顶部，但却是风格切换和波动加剧的温床。美国银行自己的牛熊指标已升至8.9，触发“卖出信号”，但报告同时指出，这并非“大顶”的信号——真正的顶部通常由债券市场和选民情绪共同确认。

这份调查的价值在于，它提供了一个清晰的“市场温度计”和一套可验证的观察框架。它告诉我们，现在的拥挤交易在哪里，哪些资产正在被抛弃，以及哪些角落可能孕育着反共识的机会。但报告也留下了一些关键问题，比如AI的“繁荣”阶段还能持续多久，以及通胀二次探底的风险是否被充分定价。这些问题，正是我们接下来要深入拆解的。

这份调查中最具冲击力的数据，不是AI的拥挤度，而是利率预期的变化。认为美联储将在未来12个月内至少加息一次的基金经理占比，从5月的16\%猛增至40\%。与此同时，认为至少会降息一次的比例，从50\%骤降至28\%。这种预期的大幅扭转，直接反映在资产配置上：投资者将全球股票的超配比例从50\%削减至38\%，科技股的超配比例也从33\%降至26\%。

这意味着什么？意味着市场正在为“更高更久”的利率环境重新定价。美国银行报告指出，对短期利率的预期已升至2022年9月以来的最高水平，而这正在追赶对通胀的预期——仍有45\%的投资者预计全球CPI将走高。这种“利率预期追赶通胀预期”的过程，本身就是一种风险的释放，但同时也意味着，任何通胀数据的超预期，都可能引发利率预期的进一步上修，从而对高估值、长久期资产构成压力。

对于资产定价而言，这一转变的含义是：过去两年驱动市场的“降息交易”逻辑正在退潮，取而代之的是“通胀韧性交易”。那些对利率敏感的资产，如科技成长股、房地产投资信托（REITs），将面临持续的估值压力。而受益于利率上升的资产，如银行股、价值股，则可能获得相对收益。调查数据也印证了这一点：投资者正在买入银行股，同时削减了对科技股的敞口。

2. AI拥挤度创历史新高，但“繁荣”阶段意味着更多买家而非顶部

关于AI泡沫的讨论，这份调查提供了最直接的量化证据。80\%的基金经理认为“做多全球半导体”是目前最拥挤的交易，这是美国银行FMS历史上的最高纪录。更关键的是，当被问及AI股票处于投资周期的哪个阶段时，56\%的人选择了“繁荣”阶段——即价格获得动能，害怕错过（FOMO）吸引更多参与者。只有21\%的人选择了“狂热”阶段，即价格和估值飙升至危险区域。

这个数据分布本身就包含了一个重要的洞察：市场共识是AI尚未见顶，但拥挤度已经极高。报告将“繁荣”定义为“价格获得动能，FOMO吸引更多参与者”，这与“狂热”的本质区别在于，前者仍有增量资金进入，而后者是存量博弈的顶点。因此，56\%的人选择“繁荣”，意味着大多数投资者认为还有更多买家尚未入场，这本身就是一个看涨信号。

但这里需要警惕的是，拥挤度本身并不构成反转信号。历史上，最拥挤的交易往往能

[... middle omitted ...]

至12\%，“AI泡沫”的风险则从5\%升至28\%。这种风险排序的变化，反映了市场焦点的根本转移：从外部冲击转向内部经济变量。

“通胀二次探底”这一风险表述本身就很有意思。它不是指通胀持续高企，而是指通胀在经历一段缓和后再度抬头。这背后隐含的逻辑是：全球经济可能正在经历一个“软着陆”或“不着陆”的过程，增长韧性使得需求端压力难以根本缓解。调查中，47\%的人预期“软着陆”，40\%的人预期“不着陆”，只有5\%的人预期“硬着陆”。这种宏观背景，恰恰是通胀二次探底的温床。

对于投资者而言，这意味着资产配置的逻辑需要从“通胀下行、利率下行”切换到“通胀韧性、利率高位”。那些在低利率环境下表现优异的资产，如长期债券、成长股，可能面临持续压力。而受益于通胀和利率上行的资产，如大宗商品、银行股、价值股，则可能获得相对优势。美国银行报告中也提到了一个有趣的反向交易：做多债券、欧洲、消费、REITs，做空大宗商品、半导体、银行股。这本质上是在押注通胀二次探底不会发生，或者已经被过度定价。

\subsection{[R012] DB：地缘风险溢价退潮，最被低估的不是避险资产，而是这些被错杀的货币}
{\small\color{refgray}归类：地缘风险与汇率：和平红利下的货币重估与亚洲资产重新定价}

DB：地缘风险溢价退潮，最被低估的不是避险资产，而是这些被错杀的货币

地缘政治风险的消退，往往不是均匀地提振所有资产。真正有意义的交易，不是追涨那些已经定价了“和平”的品种，而是寻找那些风险溢价尚未完全回吐、基本面却在同步改善的货币。DB最新发布的这份外汇研究报告，核心判断很直接：美国与伊朗达成和平协议后，霍尔木兹海峡的油轮通行效率将直接决定哪些货币可以迎来系统性重估。报告筛选出了四支货币——瑞典克朗、南非兰特、智利比索和印度卢比——并给出了明确的交易逻辑。但真正值得决策者关注的，不是这些货币的短期方向，而是这份报告揭示的一个更深层问题：市场对“和平红利”的定价，可能仍然停留在情绪层面，而忽略了供给端成本结构变化的长期含义。

这份报告发布于2026年6月15日，时间点本身就值得注意。它不是在地缘冲突爆发初期、市场恐慌最严重时发布的应急分析，而是在局势已经明朗、和平协议即将落地的窗口期，系统性地重新评估汇率定价。这种时间选择本身就说明，DB认为市场对这一轮地缘风险消退的定价，仍然存在较大的结构性的偏差。

DB筛选货币的逻辑，表面上看是“油价敏感度”和“经常账户改善”，但底层逻辑其实是一个更根本的经济学命题：当一个经济体长期承受的地缘风险溢价突然消失，哪些国家能最快地把成本下降转化为竞争力提升？答案不是那些最依赖石油进口的国家，而是那些进口成本下降的同时，出口竞争力也在同步改善的国家。

瑞典克朗被列为G10中的首选交易，原因就在这里。报告指出，瑞典经济目前仍具有超过2\%的增长潜力，财政刺激持续，且正在受益于全球国防开支上升。但更关键的是，克朗是自3月以来对美元表现最弱的发达市场货币。这意味着，即便和平协议落地，克朗的反弹空间仍然远大于已经部分定价了利好的欧元或英镑。DB判断，克朗的净空头仓位仍然存在，这意味着市场尚未充分定价瑞典经济在“后冲突时代”的结构性优势。

这里有一个容易被忽略的洞察：克朗的弱势不仅仅是因为地缘风险，还因为它此前被市场当作“高贝塔版本的欧元”来交易。当欧元本身就不强时，克朗承受了双重打击。而当和平协议落地，这种“双重折扣”反而变成了“双重反弹”的潜力。对于持有欧元或美元资产的投资者来说，这是一个值得重新审视的风险收益比。

2. 利率优势只是表象，真正的护城河是“通胀目标框架的信用”

南非兰特是DB在CEEMEA区域的首选。报告给出的理由包括：油价下降、高实际利率、国内基本面改善。但真正有意思的，是报告提到的一个细节：南非储备银行（SARB）有强烈的动机维持鹰派立场，因为它正在向更低的通胀目标框架转型。

这个细节的价值在于，它揭示了一个常被市场低估的机制：一个央行在和平时期的货币政策信用，往往比它在危机时期的干预能力更重要。南非兰特的实际利率优势确实存在，但真正让兰特具备持续吸引力的，是SARB愿意为了维持通胀目标框架而牺牲短期增长空间。这种政策信用的积累，在资本流动加速的时期，会转化为更稳定的外资流入。

报告还提到，兰特可能从化肥进口和旅游流量的恢复中受益。这些细节看似零散，但合并在一起指向一个判断：南非的经济结构可能比市场想象的更有韧性。当油价下降带来的经常账户改善，与国内政治环境的相对稳定叠加，兰特的估值折价可能被快速修正。

智利比索被DB列为拉丁美洲最具价值的货币。报告的核心逻辑是：智利比索是过去一段时间里，相对于其综合贸易条件表现最差的货币之一。当油价回落、铜价维持高位，智利经常账户赤字有望在年底前收窄。

这里需要展开一个二阶判断。智利的经济结构高度依赖铜出口，而铜价在这一轮地缘冲突中并未像油价那样大幅波动。这意味着，智利承受的“输入性通胀”压力主要来自油价，而它的出口收入并未因冲突而受损。这种“成本上升、收入不变”的组合，理论上应该导致比索贬值，但贬值的幅度是否

[... middle omitted ...]

，卢比空头平仓的动能仍然充足。

更关键的是，印度储备银行（RBI）已经采取了一系列强力措施来吸引短期资本流入，同时通过长期债券市场支持措施，为印度债券纳入全球指数铺平道路。这种“短期刺激+长期结构性改革”的组合，在DB看来，意味着卢比可能迎来一个“双重流入”的阶段：先是投机性空头平仓带来的资本回流，然后是长期被动资金因指数纳入而持续流入。

对于关注新兴市场资本流动的读者来说，印度卢比的故事提供了一个重要的观察框架：当一个国家同时具备“地缘风险溢价消退”和“国内政策框架改善”两个条件时，汇率的重估往往不是线性的，而是会在某个节点加速。印度目前正处于这个加速的临界点上。

这份报告最值得钦佩的地方，是它的克制。DB没有把所有逻辑都推演到极致，而是留下了一些关键问题。比如，报告推荐做多这些货币，但同时指出，这些多头头寸可以在年底前根据对美元和美联储的判断，选择用美元或欧元来融资。这个细节暗示，DB对美元的走势并没有给出明确的方向性判断，而是选择了一种“中性”的融资策略。

\subsection{[R013] 摩根斯坦利：油价真正的风险不是地缘缓和，而是供需错位的结构性固化}
{\small\color{refgray}归类：能源与大宗：供需错位结构性固化，上游资本开支超级周期存疑}

摩根斯坦利：油价真正的风险不是地缘缓和，而是供需错位的结构性固化

这份报告最值得认真对待的判断，不是霍尔木兹海峡即将重新开放，而是：即便中东供应恢复，市场依然面临一个比想象中更顽固的供需错位格局。摩根斯坦利在美伊谅解备忘录签署后迅速下调了近两个季度的布伦特油价预测，3季度从100美元调至90美元，4季度从95美元直降至80美元。但真正值得推敲的，不是这10-15美元的调整幅度，而是调整背后的逻辑链条——报告揭示了一个比“地缘政治溢价消退”更深刻的结构性问题。

当前市场最容易被误解的地方在于：人们倾向于把油价疲软归因于霍尔木兹海峡即将重新开放这一预期，但报告通过大量实物市场数据证明，物理原油市场的疲软早在政治突破之前就已经形成，且其驱动因素——美国出口高企与中国进口低迷——短期内看不到逆转迹象。这意味着即便中东供应恢复节奏符合预期，油价的上行空间也可能比多数人想象的更窄。

这份报告发布的时间点很关键。美伊备忘录签署于6月14日，比摩根斯坦利此前假设的“7月底”提前了约两周。正是这个时间差，让机构得以重新审视此前预测中隐含的一个关键假设：在供应恢复之前，美国出口下降和中国进口回升能否为市场提供缓冲。现在来看，这个缓冲机制不仅没有启动，反而在持续恶化。

摩根斯坦利将中东产量恢复的时间线整体前移了约1-2周。按照最新估算，7月中旬开始恢复，9月恢复50\%，12月恢复80\%，剩余部分在2027年初完成。这个节奏看起来不算慢，但报告真正想表达的是：供应恢复的速度本身并不是当前油价的定价核心。

关键在于，即便在日均约1100万桶的中东原油产量处于停产的背景下，实物市场已经出现了反常的疲软。布伦特现货价格、迪拜升贴水、成品油裂解价差均在走弱，亚洲石脑油裂解价差在过去四周内下跌了13美元/桶，接近一年低点。对于一个全球约33\%的海运原油供应被切断的市场而言，这种价格行为是反直觉的。

报告用一个细节数据点破了这个矛盾：未售出货轮的数量正在上升，且范围在扩大。尼日利亚、安哥拉、刚果、CPC混合油、利比亚、阿塞拜疆、甚至俄罗斯ESPO，多个地区的现货原油都存在明显的销售困难。这不是某个单一品类的局部问题，而是全球性的需求吸收能力不足。对于一个理论上面临巨大供应缺口的市场，出现如此广泛的未售库存，只能说明一个问题：需求端的问题比供应端的问题更严重。

2. 美国出口与中国进口构成了一个“10 mb/d的缓冲垫”，但它的持续性正在被质疑

报告最核心的分析框架，是美中两国在实物市场上扮演的“缓冲角色”。美国净出口从去年的约500万桶/日跃升至近期的约900万桶/日，中国净进口则从约1300万桶/日骤降至约700万桶/日。两者合计，相当于为全球市场腾出了约1000万桶/日的供需宽松空间。

这个数字意味着什么？中东停产带来的供应缺口，大部分被美中两国的反向变动所对冲。如果没有这个缓冲，布伦特价格可能早已远超100美元。但问题在于，这个缓冲机制能持续多久？

报告通过追踪已知油轮动向给出了一个令人担忧的答案：6月美国净出口仍将维持在850-900万桶/日的水平，与5月基本持平；中国净进口在6月可能比5月更低，7月的数据也尚未显示出改善迹象。更关键的是，中国炼厂和贸易商的现货原油采购活动——通常占进口量的约一半——在4月已降至异常低的约450万桶/日，5月进一步降至约410万桶/日，而6月已知的交易量甚至低于这两个月的同期路径。

这意味着，市场此前期待的“中国需求回升”这一关键变量，至少在可预见的未来几周内不会兑现。而美国出口能否持续维持高位，同样存在不确定性——高出口本身对美国国内库存和炼厂开工率会形成压力。

3. 需求侧的真实状况比数据呈现的更复杂，报告留下了一个关键悬疑

报告指出，5月全球炼厂原油加工量同比下降约580万

[... middle omitted ...]

一个根本性的盲区：被延迟的采购需求何时会回归？如果地缘风险进一步下降，这些被压抑的需求是否会集中释放，从而在某个时间点形成价格脉冲？反之，如果延迟采购持续更久，是否会进一步压制现货价格，甚至倒逼上游减产？

这个问题的答案，将直接决定3季度布伦特能否从当前水平反弹至90美元，还是连这个目标都显得过于乐观。

4. 报告尚未完全回答的关键问题：被延迟的需求是否会形成价格脉冲

摩根斯坦利虽然下调了预测，但依然认为3季度市场偏紧、布伦特可能从当前水平反弹。然而，这个判断成立的前提是：被延迟的采购需求必须在某个时间点回归，而且回归的规模要足够大。

报告对此并未给出确定性结论。它指出了延迟采购的存在，也承认数据不足以量化其规模，但未能回答一个更关键的问题：延迟采购的回归路径是怎样的？是平滑恢复，还是集中爆发？如果大量延迟的采购需求在某个时间窗口内同时涌入市场，可能会在短期内推高价格，甚至超过摩根斯坦利目前的3季度目标。但如果回归过程缓慢且分散，则价格反弹的力度可能非常有限。

\subsection{[R014] GS：AI ASIC仍是首选主题，但市场低估了CPU和下一个供应链瓶颈的定价权}
{\small\color{refgray}归类：AI/算力与半导体：从GPU短缺到系统级重构，互联与存储定价权转移}

GS：AI ASIC仍是首选主题，但市场低估了CPU和下一个供应链瓶颈的定价权

当大多数投资者还在追问AI需求能否持续时，一份来自GS的最新路演反馈显示，对话的焦点正在发生深刻的转移。在刚刚结束的美国投资者路演中，GS团队与超过35位机构投资者进行了交流。一个清晰的信号是：市场情绪依然高度建设性，AI ASIC（专用集成电路）继续是最受追捧的投资主题。但真正值得关注的，不是需求本身是否强劲，而是讨论的层次已经进化到了下一个阶段——CPU需求、供应链定价权以及超大规模资本支出的可持续性。

这份报告的洞察力不在于重复“AI很好”，而在于它捕捉到了一个关键转折：投资者正在从“要不要买AI”转向“买谁能在下一阶段继续增长”。增量资本的配置正变得高度挑剔，半导体持仓已经成为许多机构的最大仓位之一。这意味着，能够继续证明盈利增长可持续性的公司，才能获得进一步的资本青睐。

在众多讨论中，GS提炼出了四个最核心的投资者关切：AI供应链的下一个瓶颈在哪里、生态系统的定价权如何分配、哪些公司能在未来几年维持超速增长、以及市场过热的下行风险是什么。本文将基于这份报告，拆解这些问题的逻辑层次，并指出报告尚未完全回答的关键变量。

1. 投资者的关注点已从AI需求转向货币化验证和资本支出可持续性

一个反复出现的问题是：当前AI投资周期的可持续性如何？这并非空穴来风。超大规模云服务商（CSP）正在加速资本支出，而投资者开始追问这些投入何时能转化为可观的回报。GS注意到，谷歌近期的融资活动引发了显著关注。历史上，美国投资者倾向于将外部融资视为过度投资或泡沫行为的信号。但情绪正在演变——越来越多的投资者开始将谷歌的融资能力视为对长期AI货币化机会的信心标志。

这一变化意味着什么？它表明市场正在成熟：从纯粹的“讲故事”阶段，进入需要“证明自己”的阶段。投资者不再满足于“AI有巨大潜力”的叙事，他们需要看到具体的货币化路径。对于产业链上的公司而言，这意味着那些能够清晰展示其客户AI投入回报率的公司，将获得更高的估值溢价。而那些只是泛泛受益于AI浪潮的公司，可能会面临越来越挑剔的审视。

然而，报告也指出，关于“超大规模资本支出增长是否需要在2028年前进一步加速以支撑供应链扩张”的辩论仍在持续。这是一个尚未完全解答的关键问题。如果资本支出增速放缓，供应链中哪些环节将首当其冲？这个问题本身，就指向了下一个讨论焦点。

2. 代理型AI正将CPU需求推向前台，这是市场最被低估的结构性变化

在本次路演中，最引人注目的讨论转向之一，是投资者对CPU（中央处理器）兴趣的显著提升。长期以来，市场注意力几乎完全被AI加速器（如GPU、ASIC）所吸引。但GS的报告揭示了一个关键的逻辑演变：虽然AI加速器仍然是训练工作负载的主要受益者，但CPU在推理和代理型AI（Agentic AI）应用中可能扮演更重要的角色。

代理型AI，即能够自主执行复杂任务的AI系统，其工作负载对CPU的依赖远高于传统训练。这意味着，随着AI应用从训练走向大规模部署，CPU基础设施的重要性将与加速器并驾齐驱。GS团队特别强调了这一主题，认为它对于Aspeed和WinWay等公司而言，是仍未被市场充分定价的增长驱动因素。

这一判断的“所以呢”含义非常明确：如果CPU需求因代理型AI而出现结构性增长，那么整个服务器硬件生态的竞争格局和估值逻辑都将被重塑。那些在CPU相关领域有深度布局的公司，将获得一个独立于GPU/ASIC周期的增长引擎。这不仅是一个增量机会，更可能是一个新的估值锚点。

3. 台积电的盈利预期仍过于保守，定价权辩论正从“能否”转向“何时”

作为核心持仓的台积电，在本次路演中的讨论焦点也发生了微妙但重要的变化。投资者不再质疑需求，而是开始激烈辩论其定价权。一

[... middle omitted ...]

改善。对于投资者而言，这意味着台积电的盈利安全边际可能比普遍认知的要高。

4. 联发科的AI ASIC机会：市场认知仍停留在早期，仓位配置明显不足

联发科是本次路演中讨论最多的公司之一，但这恰恰说明市场对其机会的认知尚不充分。GS指出，虽然投资者普遍理解其AI ASIC业务的长期机会，但许多人仍低估了下一个项目（2028年）中潜在的定价和内容扩展。关键在于，市场对下一代芯片的ASP（平均销售价格）提升幅度缺乏清晰认知。

更值得关注的是，投资者对一系列长期风险的辩论，在GS看来为时过早。这些风险包括：未来几代产品中内容贡献是否会下降、下一代芯片的内容是否低于预期、以及美国CSP客户是否会逐步将更多芯片设计能力内部化。报告的核心判断是：联发科的第一颗AI ASIC芯片尚未进入量产（商业化预计在2026年底），这颗芯片本身就是一个关键里程碑，它验证了联发科在领先AI ASIC项目中的参与能力。更重要的是，除了第一个项目，联发科正在积极接触多个未来的AI ASIC机会。

\subsection{[R015] 美国银行：市场真正低估的不是需求，而是AI定价权从算力向硬件的转移}
{\small\color{refgray}归类：亚洲央行与区域市场：从被动加息转向主动加息，市场定价不足}

美国银行：市场真正低估的不是需求，而是AI定价权从算力向硬件的转移

这份6月发布的美国银行亚洲基金经理调查，捕捉到了一个关键信号：投资者对AI的乐观情绪并未消退，但他们的仓位结构正在发生一次静默但深刻的重组。41\%的受访者仍然没有对AI交易的下行风险做任何对冲，而与此同时，印尼已经取代印度成为最不受欢迎的市场，科技硬件在仓位中的权重首次超过半导体。这些数据点合在一起指向一个核心判断：市场正在将AI叙事从“算力军备竞赛”重新定价为“硬件落地周期”，而多数投资者尚未完全理解这一转变对区域配置和行业选择的含义。

这并不是说AI泡沫即将破裂。恰恰相反，调查显示只有9\%的投资者认为AI对股市的正面影响已经完全被定价，41\%认为“大致合理”，另有41\%认为“部分被定价”。关键分歧在于：AI的第二波受益者可能不再是第一波中的那些名字。美国银行的这份报告提供了足够的截面数据来支撑这一判断，但报告本身并未将其提炼为可操作的框架——这正是本文试图补全的部分。

1. 能源安全焦虑的骤降正在释放被压抑的增长预期，但市场对这一传导链条的定价仍不充分

调查中最引人注目的变化不是仓位，而是情绪底层的宏观变量。4月份有91\%的受访者表示对亚太地区的能源安全问题“高度或极度担忧”，到6月这一比例骤降至18\%。与此同时，“中度担忧”成为主流观点，占比68\%。这一变化的速度和幅度在调查历史上都极为罕见。

能源安全担忧的消退直接反映在增长预期上。亚太（除日本）的企业盈利预期在6月进一步走强，净看多比例达到50\%，远高于长期均值。更值得注意的是，认为“共识盈利预测过高”的比例已降至历史第14百分位——这意味着投资者不再认为当前盈利预期存在系统性高估。

这里的逻辑链条是清晰的：能源成本的不确定性曾是压制企业盈利和投资意愿的关键变量。当这一变量快速正常化，被压抑的资本开支和产能利用率回升就会转化为盈利上修。但市场对这一传导的理解似乎仍停留在“宏观改善”层面，而没有进一步拆解哪些行业和地区最受益于能源成本正常化。报告没有提供这一层面的分析，但我们可以从仓位变化中反推：如果能源安全改善是普遍性的，那么仓位变动应该广泛分布于各行业，但实际情况并非如此——资金在向少数板块集中，这暗示投资者可能正在错过一个更结构性的机会。

2. 科技硬件超越半导体成为最受青睐的板块，这不仅是仓位轮动，更是AI落地阶段的先见信号

在亚太（除日本）市场中，基金经理的行业配置正在经历一次有意义的切换。6月数据显示，科技硬件已经超越半导体成为最受偏好的超配板块，金融服务业也首次跻身最受青睐的行业之列。同时，资金从材料和非必需消费品（除零售）中流出，转向金融服务和电信。

这一轮动的含义需要放在AI周期的背景下理解。半导体作为AI算力的核心载体，在过去18个月中享受了极高的估值溢价和仓位集中度。但科技硬件——包括服务器、网络设备、终端设备等——是AI从训练走向推理、从云端走向边缘的物理载体。当投资者开始从半导体转向硬件，这意味着市场正在为AI的“落地阶段”做布局，而不仅仅是追逐算力本身。

报告中的另一组数据支持这一判断：在问及“哪个市场最受益于AI下一阶段”时，台湾以41\%的得票率继续领先，但美国的得票率从5月的10\%翻倍至18\%。台湾在AI硬件供应链中的核心地位是众所周知的，但美国得票率的上升暗示投资者开始关注AI硬件在美国本土的部署——这可能是数据中心建设、边缘计算设备以及企业级AI基础设施投资的结果。

需要谨慎的是，这一轮动是否可持续。报告显示，仍有41\%的投资者对AI交易没有任何对冲，其中27\%保持超配。这意味着科技硬件的超配可能部分来自“被动轮动”——即投资者在减仓半导体后，将资金配置到同一主题中的其他板块，而非真正完成了对AI落地节奏的独立判断。这一假设尚未被

[... middle omitted ...]

原材料周期，而资金正在从这些板块中流出。

这背后是一套统一的逻辑：市场正在用“AI落地能力”作为区域配置的筛选标准。台湾和韩国凭借硬件供应链优势持续受益，日本受益于盈利改善和货币政策正常化预期，而印度和印尼则因为缺乏与AI落地的直接关联而遭到减仓。即使是美国，其得票率的上升也更多是因为投资者认为美国企业将在AI硬件部署中受益，而非因为美国宏观经济的相对优势。

这一框架对区域配置的含义是清晰的：未来几个季度，市场的区域偏好将越来越取决于各地在AI硬件供应链中的位置，而非传统的GDP增长或估值差异。报告没有明确说出这一点，但它提供了足够的数据来推导这一判断。

4. 日本市场的叙事正在从“治理改革”转向“盈利驱动”，但BOJ政策正常化仍是最大的不确定性

日本市场在调查中呈现出有趣的叙事切换。6月数据显示，41\%的投资者认为“盈利”是日本股市近中期的关键驱动力，高于5月的38\%；而“央行政策正常化”的提及率从19\%上升至27\%，“公司治理改革”则从19\%下降至14\%。

\subsection{[R016] GS：中国券商真正的结构性机会不在A股交易量，而在香港的资产负债表扩张}
{\small\color{refgray}归类：中国权益与地产：结构性供给优势被低估，城市更新资金落地成关键变量}

GS：中国券商真正的结构性机会不在A股交易量，而在香港的资产负债表扩张

过去一个月，中国券商股A股和H股分别上涨7\%和5\%，同期银行股仅上涨1\%和3\%。市场习惯性地将这一涨幅归因于IPO回暖或交易量放大，但GS最新发布的中国券商与资管行业研报提出了一个更值得深思的判断：券商板块相对于银行的结构性吸引力，核心驱动力并非周期性收入反弹，而是一场发生在资产负债表层面的、由香港业务驱动的ROE扩张。

这个判断如果成立，意味着当前市场对券商股的定价逻辑需要重新校准。传统上，投资者习惯于把券商当作“看天吃饭”的贝塔品种，用日均交易量、IPO规模等流量指标来线性外推盈利。但GS的报告清晰地展示了一个正在发生的结构性变化：头部券商正在通过资本注入和杠杆提升，系统性地将香港子公司的ROE拉升至集团平均水平的近两倍，从而推动集团整体回报率的持续改善。

以下，我们基于这份研报的核心逻辑，拆解这一结构性机会的底层支撑、关键证据以及尚未被市场充分定价的变量。

1. 银行与券商的分化焦点已从“谁增长更快”转向“谁的资产负债表更有定价权”

GS在报告开头就点明了一个行业层面的范式转换：随着银行体系贷款增速放缓，盈利增长已不再是银行估值的首要决定因素，投资者的关注点已转向资产负债表的韧性和资本实力。这意味着，银行股的估值框架正在从P/E向P/B回归，而P/B的核心驱动力是资产质量和资本回报率。

与此同时，券商板块正处于一个截然不同的叙事中。券商并不直接承担信用风险，其资产负债表扩张的核心驱动力是杠杆——用自有资本撬动更高的回报。GS覆盖的三家头部券商（CICC、CITIC、广发证券）的集团平均杠杆约为6倍，但其香港子公司的平均杠杆高达11倍，几乎是集团水平的两倍。更重要的是，这些香港子公司的平均ROE达到16\%，而集团平均仅为9\%。

So what：当银行在去杠杆、降增速的环境中寻找稳定性的锚点，券商恰好处于一个“加杠杆、提回报”的结构性拐点。市场对这两类资产的定价逻辑正在发生方向性的背离——银行的价值取决于其能否守住现有资本回报率，券商的价值则取决于其能否通过国际业务拉高整体回报率。当前券商股相对于银行股的超额表现，并非短期情绪驱动，而是这一结构性差异在价格上的初步反映。

报告中最具说服力的数据来自对三家券商香港子公司与集团层面的对比。CICC的香港子公司以30\%的总资产占比贡献了48\%的收入和47\%的净利润，ROE达到15\%，是集团水平（8\%）的近两倍。CITIC的香港子公司杠杆最高，达到16.6倍，ROE高达22\%，远高于集团平均的10\%。广发证券的香港业务规模相对较小，但ROE与集团持平。

这些数字背后是一个清晰的战略意图：头部券商正在通过香港子公司进行资产负债表扩张，而这一扩张的燃料来自资本市场融资。报告特别提到，CITIC在完成160亿元人民币的再融资后，其国际业务杠杆将从16.6倍降至10.6倍，这不仅没有削弱其回报潜力，反而为未来进一步的杠杆提升创造了空间。

CICC在管理层交流中明确将中期ROE目标上调至13-15\%，并明确指出资本向香港部署是主要驱动力。这意味着，香港业务不再是券商的一个“海外分部”，而是正在成为集团整体ROE提升的核心引擎。

So what：对于投资者而言，评估一家券商未来ROE走势的关键指标，不应再是A股交易量或经纪业务市场份额，而应是其香港子公司的杠杆水平和资本注入节奏。那些能够持续为香港子公司补充资本、同时保持资产质量的公司，将享有更确定的ROE上行路径。这一判断也意味着，券商股估值已经从“周期性折价”向“结构性溢价”的切换过程中。

3. IPO回暖是催化剂，但真正支撑估值的是IPO背后的“结构性杠杆”

报告用较大篇幅描述了香港IPO市场的复苏。GS策略团队指出，2025年

[... middle omitted ...]

H股券商估值从12倍PE跌至7倍PE。

So what：市场真正应该关注的，不是IPO带来的短期投行收入增量，而是IPO活动如何为券商的香港子公司提供资产负债表扩张的“原材料”——即通过承销、配售、做市等业务，香港子公司可以更高效地运用杠杆、扩大资产规模、提升ROE。IPO不是终点，而是杠杆扩张的通道。这也是为什么GS在推荐标的时，更看重CICC和CITIC在香港市场的领先地位和清晰的资本部署计划，而非广发证券更低的估值。

尽管报告对券商通过香港业务提升ROE的逻辑阐述得相当充分，但有一个关键问题并未完全展开：杠杆扩张的边界在哪里？或者说，香港子公司的ROE提升是否有天花板？

报告显示，CITIC香港子公司的杠杆已高达16.6倍（再融资后降至10.6倍），而CICC和广发证券的香港子公司杠杆分别为7.6倍和9.9倍。从国际投行的经验来看，高杠杆确实能放大ROE，但也意味着对市场波动的敏感度更高。2022年港股市场的剧烈调整，就对部分券商的香港业务产生了显著冲击。

\subsection{[R017] GS：市场低估了“城市更新”作为中国房地产下行缓冲器的真实权重}
{\small\color{refgray}归类：中国权益与地产：结构性供给优势被低估，城市更新资金落地成关键变量}

GS：市场低估了“城市更新”作为中国房地产下行缓冲器的真实权重

当市场目光仍聚焦在周度成交数据的“平台期”时，一份来自GS的最新周报揭示了一个更值得关注的信号：中央层面对城市更新的资金支持正在从框架走向明细。这或许才是当前房地产叙事中最被低估的结构性变量。

这份报告的核心判断并非关于6月销售能否反弹，而是关于政策供给侧的“再定价”。在成交量与情绪双双进入平台期的背景下，GS分析师王逸团队用一周的篇幅，详细拆解了发改委、财政部在“十五五”城市更新框架下的资金分配逻辑。这不仅仅是政策例行通报，它可能正在重塑开发商资产价值的底层评估逻辑。

对于高净值读者和产业决策者而言，当前需要追问的不是“市场何时见底”，而是“当政策资金从万亿级走向落地，哪些企业真正具备吸收这些红利的资产负债表能力”。

GS周报的数据本身并不令人兴奋：第24周新房销售面积环比微增1\%，同比持平；二手房成交面积环比持平，同比微增1\%。这是一个典型的“平台期”信号——市场既没有进一步恶化，也没有出现趋势性好转。

但这份数据背后的含义，比表面数字更重要。GS跟踪的约75个城市数据显示，6月至今新房销售面积中位数环比下降17\%，同比下降25\%。这意味着，如果没有外部变量的介入，市场自发修复的动力正在衰减。二手房市场虽然相对活跃（YTD同比+1\%，较2024年水平+22\%），但中介价格预期指数（CSI）环比下降0.9个百分点，卖方挂牌价格预期（CAI）虽然稳定，但中介情绪已经先行走弱。

数据传递的清晰信号是：市场正处在一个需要外力打破的均衡点。而这个外力，GS认为正是城市更新的资金落地。

GS周报中花费最大篇幅讨论的，并非市场数据，而是中央部委对城市更新资金的具体分配。这本身就说明，该机构认为政策细节比短期成交数据更具前瞻意义。

发改委此次明确了两笔关键资金：970亿元中央预算内投资，专项用于老旧小区和危旧房改造，预计覆盖约800万户家庭；1600亿元超长期特别国债（同比增加250亿元），专项用于地下管网升级。财政部则确认了重点城市补贴机制，2026年新增15个城市，累计覆盖约50个重点城市。

第一，资金规模在增长，但更重要的是资金结构在优化。970亿元的老旧小区改造资金，对应的是约800万户家庭，这意味着每户家庭的平均财政支持力度在提升。这不再是象征性的“刷墙式”改造，而是具备实质拉动力的资本支出。

第二，财政部强调的“多管齐下融资工具箱”——直接财政补贴、地方政府专项债、税收激励——表明，城市更新正在从一个部门政策升级为跨部门协同的资本运作框架。这与过去单纯依赖开发商预售资金回流的模式，有本质区别。

对于房地产企业而言，这意味着资产价值的评估逻辑正在发生位移。过去，一个项目的价值取决于销售去化速度；未来，位于城市更新重点区域的项目，其价值可能更多取决于能否获得财政资金配套和改造补贴。

GS通过其自有的浮法玻璃供需模型（GSPC）推算出，5月竣工面积同比降幅可能仍为高双位数（high-teens），而全年竣工面积预计同比仅下降1\%。这一数据看似矛盾，实则揭示了行业供给端的结构性分化。

高双位数的月度竣工下降，反映的是过去两年新开工面积大幅萎缩的滞后效应。而全年降幅收窄至1\%，则暗示下半年可能存在竣工脉冲——这很可能来自城市更新和保交付项目的集中推进。

更值得关注的是新开工数据。GS根据300城土地成交趋势和全国水泥发货率推算，5月新开工面积同比降幅可能在低二十位数（low-twenties）。这意味着，即使有政策托底，开发商主动投资的意愿仍然极低。土地市场的冷清，正在转化为未来1-2年可售资源的持续收缩。

这种“竣工韧性、新开工疲弱”的分化格局，对投资者的含义是：短期内，库存去化周期可能因竣工交付而有所改善（当前库存月数为27

[... middle omitted ...]

这些周期底部。

但这里需要区分“低估值”和“价值陷阱”。历史低谷时期的市净率虽然更高（例如2008年离岸0.7倍、在岸1.6倍），但彼时的资产质量和盈利前景与当下完全不同。

当前估值折价的合理性，取决于两个关键假设：第一，NAV本身的减值是否已经充分反映；第二，城市更新政策能否真正转化为开发商的盈利改善。GS的周报并没有给出明确的结论，但提供了观察框架——持续跟踪城市更新资金的落地速度和覆盖范围，将是判断估值修复拐点的核心变量。

5. 报告尚未完全回答的关键问题：资金撬动效应与开发商资产负债表之间的鸿沟

尽管GS详细拆解了城市更新的资金来源，但报告本身并未完全回答一个关键问题：这些资金如何从中央政府传导至开发商的资产负债表？

970亿元的老旧小区改造资金，更多是直接惠及居民和施工企业，而非开发商。1600亿元的管网升级资金，受益方主要是基建和市政工程企业。开发商能否从中获益，取决于其是否拥有城市更新项目的开发权，以及能否将财政补贴转化为项目层面的现金流改善。

\subsection{[R018] NOM：油价暴跌不是短期交易，而是亚洲资产重新定价的起点}
{\small\color{refgray}归类：地缘风险与汇率：和平红利下的货币重估与亚洲资产重新定价}

市场正在用交易员的思维理解美伊和平协议——油价跌了，能源进口国受益，这没错。但这份NOM研报真正值得关注的判断是：油价下行对亚洲的影响远不止于通胀和经常账户的数字改善，它正在改变三个更根本的结构——央行政策路径的重新校准、AI资本开支与能源成本下降的叠加效应、以及区域内国家之间“受益程度”的显著分化。

如果只看油价从100美元以上跌到83美元以下这个结果，你可能会认为这是一个短期交易机会。但NOM的分析框架告诉我们，当油价在80美元附近企稳，2026年全年均价将从基准假设的93.9美元降至85.4美元，下半年油价将比当前假设低约15\%。这意味着的，不是一个月的数据波动，而是一个季度的宏观假设重置。

更关键的是，这份报告在“大家都受益”的叙事下埋了一条暗线：受益的程度、方式、时间差，才是决定投资回报的核心变量。泰国、菲律宾、印度、韩国被列为最大赢家，但背后的逻辑完全不同——有的是进口账单下降，有的是通胀快速回落，有的是经常账户顺差扩大，还有一个可能享受“双重红利”。这些差异，才是资产定价的关键。

1. 油价每跌10\%，亚洲的宏观画面就换一次底色，但效果不是同步的

NOM用一组清晰的弹性系数告诉我们油价对亚洲经济的传导路径。每10\%的油价下跌，平均提升亚洲经常账户约0.3\%的GDP，压低通胀0.2-0.3个百分点，同时对增长和财政也有边际正向贡献。

但这些数字的平均值掩盖了最重要的信息：国别差异。泰国每10\%油价下跌对经常账户的改善是0.5个百分点，是所有亚洲经济体中最大的；印度是0.4个百分点，韩国是0.3个百分点。而在通胀端，印度和菲律宾的弹性最大，每10\%油价下跌可压低通胀0.5个百分点。

这意味着什么？当油价从100美元跌到83美元，降幅约17\%，泰国的经常账户改善幅度可能接近0.85个百分点的GDP，印度的通胀压力可能下降近0.85个百分点。这不是边际变化，而是足以改变宏观叙事的幅度。

但NOM同时提醒了一个容易被忽视的滞后效应：供应链正常化是渐进的。霍尔木兹海峡的重新开放需要清理水雷、疏散滞留船只、重启停产油井，保险成本在初期仍会高企，航运信心的恢复需要时间。因此，通胀改善的传导需要2-3个月才能完全体现。这是一个时间差交易机会——市场可能先定价通胀回落，但实际数据改善要等到三季度。

2. 四个最大赢家的受益逻辑完全不同，这决定了它们资产定价的驱动因素

NOM将泰国、菲律宾、印度和韩国列为油价下行的最大受益者，但细看每个国家的传导机制，你会发现它们对应的资产逻辑截然不同。

泰国是亚洲最大的净石油进口国（相对GDP），能源占CPI篮子高达11.8\%。油价的直接效果是进口账单下降和成本推动型通胀缓解。这对泰国的意义在于，经常账户从1.4\%的GDP占比开始改善，空间巨大。但NOM也明确指出，泰国的结构性约束——人口老龄化、生产率挑战、中国冲击——并未改变。因此，油价对泰国资产的影响更像是“减压”，而非“转型”。

菲律宾的情况不同。由于没有燃料补贴，国际油价直接传导至消费者价格，4月通胀已飙升至7.2\%。油价下跌在菲律宾的效果是“立竿见影”的通胀回落，这将减轻央行的加息压力。但这份报告也埋了一个伏笔：2026年是厄尔尼诺年，食品通胀的中期风险仍在。菲律宾的通胀叙事可能经历“先降后稳”的节奏。

印度是NOM认为可能享受“双重红利”的经济体。除了油价下降带来的进口成本降低、经常账户赤字收窄、通胀回落和财政补贴负担减轻之外，印度最近宣布的FCNR(B)措施预计将带来约550亿美元的流入，进一步提振国际收支盈余。这意味着印度可能同时获得“能源红利”和“资本流入红利”，两者叠加的效果可能远超单独计算。

韩国的逻辑则更多地与AI相关。油价下降叠加AI繁荣带来的半导体出口增长，韩国的经常账户顺差可能

[... middle omitted ...]

按兵不动”的阵营中，印尼央行、印度央行、澳大利亚央行、泰国央行和中国央行预计维持利率不变。但这里有一个值得深挖的细节：印尼央行维持利率的前提是印尼盾的稳定。如果油价下跌导致资本流入，印尼盾走强，那么印尼央行维持利率的空间会更大。但如果油价下跌引发全球风险偏好变化，资金流出新兴市场，印尼央行反而可能面临加息压力。

在“推迟加息”的阵营中，台湾央行是唯一一个被调整预测的。NOM将台湾央行的首次加息从6月推迟到12月，并预计12月和明年3月各加息12.5个基点，终端利率为2.25\%。这背后的逻辑是，油价下行缓解了通胀压力，让央行有更多空间维持宽松立场。但这也意味着，台湾的利率正常化进程将比市场预期更慢。

最值得关注的是“继续加息”的阵营。菲律宾央行、马来西亚央行、新加坡金管局、韩国央行和新西兰央行都预计将继续收紧政策。菲律宾央行预计累计加息75个基点（包括本周的25个基点），核心通胀仍然高企。韩国央行预计在7月、10月和明年1月各加息25个基点，终端利率达到3.25\%。

\subsection{[R019] Bernstein：液冷数据中心最被低估的环节不是冷板，而是冷却液分配单元}
{\small\color{refgray}归类：能源转型与电力基础设施：AI基建的真正瓶颈在电力供给}

Bernstein：液冷数据中心最被低估的环节不是冷板，而是冷却液分配单元

当英伟达的GPU集群单机柜功耗突破100千瓦，传统风冷已经彻底失效。液冷从边缘技术走向数据中心基础设施的核心，这已是行业共识。但多数讨论集中在冷板——那个贴在芯片上的金属块——而忽略了整个液冷系统中真正决定成败的组件：冷却液分配单元。

Bernstein这份最新的液冷专题报告，核心判断清晰且反直觉：CDU是一个被市场低估的好生意。它不是简单的泵阀组合，而是兼具技术壁垒、服务粘性和长期定价权的关键节点。更重要的是，在冷板可能走向标准化的过程中，CDU的不可替代性反而在增强。

这份报告的价值不在于告诉你液冷市场有多大——市场规模的争论仍在继续——而在于提供了一个评估CDU企业的分析框架，以及一个值得长期跟踪的技术拐点信号：两相直接接触冷却正在逼近商业化。

1. 为什么CDU比冷板更难被替代：技术复杂度和服务绑定的双重护城河

冷却液分配单元的核心功能是精确控制冷却液的流量、压力和温度，将其分配到多个机柜和机架，再回收升温后的冷却液。听起来像是一个工业泵站，但Bernstein指出，CDU的复杂程度和关键性远超表面认知。

从技术角度看，CDU需要同时管理多个物理量：流量比、压头、接近温度差、模块化扩展能力。以Bernstein对比的旗舰产品为例，Vertiv的Coolchip 2300可以处理2.3兆瓦的热负荷，流量高达每分钟3400升，压头超过50 psi；Trane Technologies的Giga-modular CDU则实现了14兆瓦的模块化扩展。这些参数背后是流体力学、热交换效率和系统可靠性的工程平衡。

但真正让CDU区别于冷板的是服务绑定。冷板一旦安装，几乎不需要维护。而CDU包含泵、阀、热交换器、传感器和控制逻辑，这些部件需要持续的监控、维护和更换。Bernstein强调，当单机柜的价值持续攀升——一个装满GPU的机柜可能价值数百万美元——CDU故障意味着多个机柜同时烧毁。这种风险让客户不会为了短期成本节省而选择最低价的服务商。

这意味着CDU市场存在一个天然的技术护城河：客户需要的是可靠的服务，而非便宜的硬件。这解释了为什么Bernstein认为CDU的标准化风险低于冷板——即使超大规模云厂商最终会指定CDU的设计规格，但服务环节仍然掌握在OEM手中。

2. 市场规模争论背后的真正变量：GW部署量、每千瓦成本和CDU使用寿命

Bernstein对CDU市场规模的估算相对克制：2026年约为低单位数十亿美元，到2030年增长至中高单位数十亿美元，复合年增长率约为15\%左右。但报告坦承，这个估算高度依赖于三个关键假设，而这三个假设正是当前市场争论的焦点。

第一个变量是数据中心新增电力容量。全球数据中心电力需求的预测从每年几吉瓦到几十吉瓦不等，差异主要来自AI推理需求的爆发速度。第二个变量是每千瓦冷却成本，这取决于液冷技术的演进路径——单相直接接触冷却的成本结构不同于两相，而浸没式冷却的成本结构又完全不同。第三个变量是CDU的使用寿命，这直接决定了替换周期和总需求。

Bernstein开发了一个专有的液冷模型来处理这些变量，但报告没有给出具体的模型细节。这实际上是一个明确的信号：CDU市场规模不是一个可以精确预测的数字，而是一个需要持续跟踪的变量。对于投资者而言，更重要的不是市场规模的具体数字，而是判断市场是否处于供不应求的状态——Bernstein明确表示，2028年之前CDU市场将处于供应受限状态，这意味着任何拥有CDU产品的企业都将获得大量订单。

Bernstein对CDU市场的竞争分析值得仔细拆解。报告没有简单列出所有参与者，而是根据技术实力和市场定位进行了分层。

第一梯队是英伟达液冷生态系统中的

[... middle omitted ...]

ne Technologies，报告用“实力超出预期”来形容。Trane的Giga-modular CDU在容量和模块化方面都处于领先地位，这与其在工业冷却领域的深厚积累有关。

第三梯队是Carrier和Johnson Controls。它们拥有CDU产品，但Bernstein发现，由于这两家公司更专注于冷水机组等传统冷却设备，其CDU的产品广度、规格参数和细节披露程度都不及前述企业。这暗示它们可能更多是跟随市场，而非定义技术路线。

值得关注的是CoolIT。报告指出，CoolIT更像是冷板公司，其CDU的接近温度差落后于竞争对手。这提醒我们，在液冷生态中，冷板能力和CDU能力并不自动等同。

Bernstein报告中最重要的技术判断，是关于两相直接接触冷却的演进。目前主流的单相直接接触冷却正在接近其理论散热极限。当机柜功耗超过某个阈值时，单相冷却液无法带走足够的热量，而两相冷却通过让冷却液蒸发来吸收更多热量，可以在不显著提高冷却液温度的情况下实现更高的散热效率。

\subsection{[R020] 摩根斯坦利：AI基础设施正在经历从“GPU短缺”到“系统级重构”的转折点}
{\small\color{refgray}归类：AI/算力与半导体：从GPU短缺到系统级重构，互联与存储定价权转移}

摩根斯坦利：AI基础设施正在经历从“GPU短缺”到“系统级重构”的转折点

市场对AI算力的讨论，过去两年几乎全部集中在GPU的产能与交付。但摩根斯坦利最新发布的这份大中华区科技硬件深度报告，传递了一个更值得关注的信号：AI基础设施的投资逻辑，正在从“谁拿到了芯片”切换到“谁能把芯片变成高效、可靠、可扩展的系统”。GPU的稀缺性正在被系统设计的复杂度取代。

这份报告覆盖了从服务器机架、电源架构、散热方案到PCB基板、数据网络互连的完整硬件链条，并明确指出了下一代平台（Vera Rubin、Kyber架构、LPX/Vera CPU机架）带来的设计升级机会。报告的核心判断可以提炼为一句话：AI服务器市场的价值分配，正在从芯片端向系统集成与关键组件端迁移，而这轮迁移的幅度，可能比市场当前定价所反映的要大得多。

为什么现在需要关注这个转折？因为市场对AI资本开支的预期已经相当充分，但对其结构变化的定价却远远不足。报告通过详尽的估值对比表显示，多数AI服务器ODM和组件供应商的2026年预期市盈率仍处于15-25倍区间，而它们的盈利增长轨迹——尤其是2027年的EPS增速——并未被充分计入。这不是一个关于“AI需求会不会持续”的问题，而是一个关于“需求结构变化下，哪些环节会获得超额分配”的问题。

报告明确指出，即将到来的Vera Rubin平台、Kyber架构以及LPX/Vera CPU机架，将带来“major design upgrades”。这不仅仅是代际更替，而是系统架构层面的重构。

从组件层面看，这些升级意味着：电源架构从传统方案向800VDC高压直流方案迁移，散热需求从风冷向液冷方案加速过渡，PCB与ABF基板的供应紧张程度进一步加剧，数据网络与互连技术向CPO（共封装光学）方向演进。每一个子系统的变化，都对应着一组供应商的重新洗牌。

市场往往低估了这种系统性升级的幅度。一个典型的例子是电源架构：从48V到800VDC的跃迁，不是简单的电压提升，而是整个配电拓扑、功率模块、连接器标准的重构。这意味着一套全新的供应链认证周期，以及伴随而来的供应商切换成本。对于已经卡位在800V DC方案的供应商，这是一个结构性壁垒，而非短期订单波动。

报告覆盖的Bizlink（3665.TW）就是一个值得关注的案例。该公司的2027年预期EPS增速接近90\%，但当前估值对应的2027年市盈率仅为19倍。这种增速与估值的错配，在报告覆盖的多家AI硬件供应商中并非孤例。

2. AI ASIC服务器的规模扩张，正在创造一条独立于GPU的供应链

报告特别强调了AI ASIC服务器——主要是TPU和Trainium平台——的升级与出货量扩张。这是一个容易被市场忽视的维度。当市场聚焦于NVIDIA GPU的迭代节奏时，定制化AI芯片（ASIC）的服务器部署正在形成一条独立的供应链逻辑。

与通用GPU服务器不同，ASIC服务器的系统集成度更高，对特定组件（如定制化散热模块、专用电源方案、高密度PCB）的需求更为集中。这意味着，ASIC服务器的放量，可能为某些细分组件供应商带来超越行业平均的订单弹性。

报告将TPU和Trainium平台单独列出，暗示这些平台的出货量增长可能超出当前市场预期。如果这一判断成立，那么相关供应链中具备定制化能力的供应商（如散热方案商AVC、连接器商Fositek），其盈利能见度将显著优于同行。

这里需要指出的是，报告并未详细展开ASIC服务器出货量的具体预测数字，但将其与GPU服务器并列讨论，本身就是一个信号：AI基础设施的供应链正在从“单极”（NVIDIA生态）向“双极”（NVIDIA + 定制ASIC）演变。这种演变对组件供应商的意义，可能比市场当前认知的要大。

3. 供应链的瓶颈正在从“

[... middle omitted ...]

“产能利用率”来评估供应链瓶颈。但进入 年，瓶颈正在向“系统级设计验证能力”迁移。

这种迁移改变了竞争壁垒的性质。在产能驱动的阶段，规模是最大的优势。但在工程能力驱动的阶段，技术积累、客户认证周期、跨学科整合能力才是真正的壁垒。报告覆盖的AI服务器ODM——如Wistron、Hon Hai、Wiwynn、Quanta——之所以被列为“key stock ideas”，不是因为它们的产能最大，而是因为它们在过去两年积累的系统级集成经验，正在转化为下一代平台的竞争护城河。

对于投资者而言，这意味着评估供应链公司时，传统的“产能扩张速度”指标的重要性正在下降，而“研发支出占比”、“客户认证数量”、“系统级解决方案的广度”等指标的重要性在上升。报告虽然没有直接给出这一框架，但从其选股逻辑中可以清晰地读出这一倾向。

报告在风险部分提到了两个需要警惕的因素：一是消费电子（智能手机/PC）需求疲软以及存储器涨价对利润率的影响；二是铜、镍等原材料价格上涨与供应紧张对成本的侵蚀。

\subsection{[R021] 摩根斯坦利：AI互联正在从“规模扩张”转向“架构重构”，市场低估了互联层定价权的转移}
{\small\color{refgray}归类：AI/算力与半导体：从GPU短缺到系统级重构，互联与存储定价权转移}

摩根斯坦利：AI互联正在从“规模扩张”转向“架构重构”，市场低估了互联层定价权的转移

这份摩根斯坦利半导体团队在2026年6月发布的巴士巡回路演纪要不是一份常规的企业拜访汇总。它传递了一个更值得关注的结构性信号：AI基础设施的投资重心，正在从算力芯片的比拼，转向互联架构的决定性博弈。三位CEO——Astera Labs、Marvell、Intel——不约而同地将叙事焦点放在互联层，而非单纯的算力性能上。这并非巧合，而是AI基础设施从粗放扩张进入精细化架构优化的必然产物。

研报的核心判断可以概括为：互联层正在从“必要组件”升级为“系统架构的定义者”。谁能在光学互联、PCIe生态、定制SerDes和先进封装上同时建立优势，谁就能在下一阶段的AI基础设施投资中占据定价权。而目前市场对这一转变的定价，仍然偏保守。

以下是基于该报告提炼的五个关键洞察，以及一个报告尚未完全回答的关键问题。

1. 互联层从“被动跟随”变为“主动定义”，Astera和Marvell正在争夺架构主导权

过去两年，AI基础设施的投资叙事几乎完全由GPU的算力密度和内存带宽主导。互联层（Scale-up与Scale-out网络）被视为“跟着GPU走”的附属品。但这份报告中的两个信号表明，情况正在逆转。

首先，Astera Labs管理层明确表示，其PCIe交换机正在被客户称为“开放生态的NVSwitch”。这不是一个营销话术，而是一个架构信号。NVSwitch是NVIDIA用于GPU间高速互联的专有技术。当客户开始将第三方PCIe交换机类比为NVSwitch时，意味着互联层正在从“标准化连接”走向“架构差异化”。Astera正在定义开放生态中的互联标准，而非仅仅提供兼容性产品。

其次，Marvell管理层将公司定位为“少数几家同时覆盖DCI、die-to-die IP、Scale-out交换、Scale-up交换、光学、SerDes、CXL、NIC和定制芯片的供应商”。这不仅是产品线广度的问题，更是系统集成能力的体现。当客户面对供应链约束和异构计算集成复杂性时，一个能提供“全栈互联”的供应商，其议价权远高于单一组件供应商。

这一转变的含义是：互联层的价值捕获正在从“按组件定价”转向“按架构价值定价”。市场对Astera和Marvell的估值模型，可能需要从“半导体公司”切换到“基础设施架构公司”。

2. PCIe的生命周期比市场预期的更长，中国市场的结构性机会被低估

市场普遍认为，随着AI工作负载向更高带宽的互联协议（如UAL、NVLink）迁移，PCIe将逐渐被边缘化。但这份报告提供了两个相反的论据。

第一，Astera管理层预计PCIe至少还会继续增长几年。原因在于，并非所有工作负载和协议都会迁移到新标准。在推理端，尤其是KV-cache offload驱动的CXL需求中，PCIe的兼容性和成熟度仍然是不可替代的。

第二，中国市场的特殊性被严重低估。由于中国许多AI系统基于Add-in-Card架构，这些系统天然是PCIe原生的。更关键的是，中国使用的GPU通常算力密度较低，需要更多的互联来弥补单卡性能的不足。这意味着，中国AI基础设施对PCIe的依赖程度和需求强度，可能高于美国市场。Astera管理层明确表示，虽然中国市场规模不会超过美国，但这是一个“有意义的增量机会”。

这一洞察对投资者的含义是：PCIe相关供应链（尤其是Astera的PCIe交换机）的需求韧性可能比市场预期的更强。市场对“PCIe被淘汰”的叙事可能过于线性，而低估了其在特定场景和地理市场中的结构性价值。

3. 光学互联的拐点正在到来，但路径并非单一的CPO，NPO可能成为中间态

光学互联（尤其是共封装光学CPO和近封装光学NPO）是AI基础设

[... middle omitted ...]

迟部分客户对CPO的采纳，而NVIDIA则可能直接跳到CPO。这意味着，光学互联的供应链格局不是线性的“CPO取代NPO”，而是两条并行路径，各自服务于不同的客户群体和部署场景。

这一洞察的关键在于：光学互联的TAM（总可寻址市场）被普遍看好，但市场对“谁将在哪个时间点捕获价值”的定价仍然模糊。Astera通过收购aiXscale获得了电IC、PIC和光学引擎的整合能力，这使其在光学互联的价值链中具有更完整的定位。但报告也承认，CPO对Astera而言更可能是2028年之后的机会。投资者需要区分“技术叙事”和“收入落地”之间的时间差。

4. Intel的转折点不在于CPU份额，而在于14A工艺能否成为“可信的备选”

Intel的叙事在这份报告中发生了有趣的转变。新任CEO Lip-Bu Tan的优先级已经从“文化重塑”转向“业务增长和产品路线图执行”。但报告的核心判断是：Intel的短期股价上涨更多来自先进制程短缺带来的CPU供应紧张，而非市场份额的结构性恢复。

\subsection{[R022] GS：霍尔木兹海峡若重开，最被低估的不是油价，而是油运定价权的结构性重置}
{\small\color{refgray}归类：交通运输：油运定价权结构性重置，VLCC日租金弹性飙升}

GS：霍尔木兹海峡若重开，最被低估的不是油价，而是油运定价权的结构性重置

市场正在关注霍尔木兹海峡可能重启的地缘信号，但多数讨论仍停留在“油价会跌多少”的层面。GS在6月14日发布的一份中国交通运输行业专题报告中，给出了一个更具穿透力的判断：如果海峡在2026年7月底前恢复正常通行，全市场最被低估的变量不是原油价格，而是VLCC（超大型油轮）日租金可能飙升至25万美元/天，甚至35万美元/天——这分别对应了航运公司盈利57\%和114\%的上行空间。

这份报告之所以值得认真读，不是因为它给出了一个乐观数字，而是因为它揭示了一个正在发生的结构性变化：过去两年中国和韩国船东的激进整合，已经从根本上改变了VLCC市场的定价机制。当供给侧的议价权从分散走向集中，同样的需求波动，会带来完全不同的价格弹性。

1. 市场真正没算清的账，是油运定价权已经从“船多货少”转向“船东说了算”

多数投资者对油运周期的理解，仍停留在“运价由供需缺口决定”的古典框架里。但GS的报告指出，这个框架已经过时了。

2021至2025年间，VLCC市场每1个百分点的净需求变化，只能带来约6000美元/天的运价波动。而到了2026年前两个月，这个弹性系数已经跃升至2万美元/天——翻了3倍以上。背后的原因，是前十大合规船东的市场份额从2025年的47\%飙升至68\%，集中度在短短一年内完成了质的飞跃。

这意味着什么？意味着即便需求恢复的速度没有超预期，船东的定价能力也已经今非昔比。过去，运价上涨会迅速触发新船订单和闲置运力回归，形成价格天花板。但在当前格局下，头部船东有能力通过控制运力投放节奏，把运价维持在高位更长时间。

这份报告的核心贡献，不是给出了一个情景假设，而是提醒市场：即使不考虑霍尔木兹重开，VLCC市场的定价逻辑也已经发生了不可逆的转变。海峡事件只是放大了这个转变的财务后果。

2. 补库存需求不是一次性脉冲，而是未来12个月的结构性支撑

霍尔木兹关闭期间，全球原油库存被持续消耗。GS大宗商品团队的数据显示，截至2026年6月底，全球原油库存较2月水平已下降约7.05亿桶。如果海湾出口在7月底恢复正常，补库需求将相当于全球海运原油贸易量的4.8\%。

这个数字本身并不惊人。但关键点在于：补库不是一个月内完成的脉冲，而是需要大约12个月去填补的持续需求。这意味着，即便地缘溢价消退，油运市场在未来一年仍将面临一个额外的、可量化的需求增量。

更值得注意的是，GS在报告中提出了一个“蓝色天空”情景——如果伊朗石油制裁同时解除，还会有约5个百分点的运输需求从影子船队转向合规船队。这部分需求虽然不直接增加总运量，但会显著压缩合规船东面临的“灰色运力”竞争，进一步强化头部船东的议价权。

这里有一个尚未被充分讨论的问题：影子船队的退出速度，取决于制裁执行力度和保险、港口合规要求的收紧程度，而这些变量并不完全由市场决定。GS的报告给出了一个方向性判断，但具体节奏仍需要持续跟踪。

3. 航空公司的受益逻辑比表面看起来更复杂，但弹性集中在三大航

油价每桶下降10美元，对三大航H股带来约60\%-70\%的股价上行空间——这个数字看起来很有吸引力。但仔细拆解会发现，不同航空公司的受益逻辑存在显著差异。

对于三大航，油价下跌的直接效果是降低燃油成本，同时通过降低燃油附加费间接刺激出行需求。GS估算，这一组合能推动三大航盈利改善16\%-26\%。但对于春秋航空这类低成本航司，受益幅度只有4\%，因为其燃油附加费占比更低，且票价敏感性不同。

这意味着，如果投资者想押注霍尔木兹重开带来的航空板块机会，需要区分不同商业模式的弹性差异。三大航的H股可能是更直接的受益标的，而春秋航空的逻辑更多依赖于自身运营效率，油价只是边际变量。

报告没有展开的一个问题是：如

[... middle omitted ...]

续性的信心。如果市场普遍认为高运价是“一次性补库”带来的短期现象，船东可能不会大规模追单。

报告提到了中国船厂产能扩张对远期供给的潜在压力，但没有深入讨论船东的订单行为是否会因为“担心错过当前窗口”而提前集中下单。这是一个值得在完整报告中细读的部分——GS对2027/2028年VLCC交付量的预测显示，2028年交付量占2025年底船队规模的10.7\%，这已经是一个不小的供给增量。

对于集装箱航运，报告给出了“适度下行风险”的判断，理由包括：低油价导致航速提升、以及红海复航可能释放约10\%的有效运力。这个判断与油运的乐观形成鲜明对比，也提示投资者：同一场地缘事件，对不同子行业的影响可能完全相反。

5. 这份报告尚未完全回答的关键问题：估值折扣是否足够补偿远期供给风险？

GS在情景分析中，对中远海能A/H股的目标市盈率分别打了25\%的折扣（10倍/8倍），理由是2028年之后中国船厂产能扩张可能带来更多新船交付。这个折扣幅度是否充分，取决于两个尚未被验证的变量：

\subsection{[R023] GS：美国供应链的“新常态”比市场预期的更早到来}
{\small\color{refgray}归类：宏观与利率：财政供给与通胀再平衡主导定价}

当所有人还在讨论关税冲击和地缘冲突是否会引爆新一轮供应链危机时，GS的最新周度追踪指数给出了一组反直觉的信号：截至2026年6月中旬，其供应链拥堵综合指数环比下降4\%，瓶颈评分连续多周维持在“2”的水平。这个数字意味着什么？在2021年峰值时期，该评分为“10”，而疫情前的基线水平约为“1.8”到“2.0”。换言之，美国供应链的拥堵程度，已经基本回到了疫情前的“正常”状态。

这不是一个关于“好转”的新闻，而是一个关于“结构性重置”的起点。市场目前对供应链风险的定价，依然停留在“高波动、易中断”的旧叙事中。但GS的数据暗示，一个更稳定、更可预测的供应链新常态可能已经提前到来。如果这一判断成立，那么它对全球贸易流、企业库存策略、航运运价以及资产定价的重新校准，将远比一次简单的数据改善更为深远。

这份报告最值得关注的判断，并非拥堵改善本身，而是改善的“质量”：它不是由需求崩溃驱动的被动去库存，而是在贸易流仍保持一定活跃度下的系统效率提升。这背后，是港口基建投资、铁路运营优化以及集装箱周转率的结构性改善。这意味着，即使未来关税政策或地缘局势出现扰动，供应链的韧性也可能比预期更强。

1. 拥堵指数回归“2”并非偶然，而是系统效率提升的累积结果

GS的周度瓶颈指数在2026年6月降至“2”，这是自2022年7月以来首次稳定在这一水平。从月度平均评分来看，自2024年初至今，该指数几乎一直围绕“1.8”到“2.0”的区间窄幅波动。相比之下，2021年8月至10月间，该指数曾连续三个月维持在“9”至“10”的极端水平。

这种长期、低波动率的改善，与2022年下半年的“假性正常化”有本质区别。当时拥堵缓解的主要驱动力是需求的急剧放缓——消费者从商品消费转向服务消费，零售商库存积压导致进口订单锐减。而当前的环境不同：GS追踪的西海岸铁路多式联运货运量同比增速在6月初加速至16\%，较前一周的10\%明显提升；东海岸集装箱船等待数量维持在5艘，西海岸等待数量仅为1艘。这说明贸易活动本身并未萎缩，拥堵的缓解是基础设施和运营效率提升的结果。

对于企业决策者而言，这意味着“库存安全垫”的逻辑需要重新审视。过去几年，许多企业为了对冲供应链中断风险，被迫维持高于正常水平的库存。如果供应链的流动性确实已回归疫情前水平，那么这些超额库存将不再具备风险溢价，反而可能成为资金占用的成本。那些率先调整库存策略的企业，将在资本效率和运营弹性上获得双重优势。

在供应链拥堵的所有指标中，集装箱底盘（chassis）的周转时间是最能反映“毛细血管”通畅程度的指标。底盘是连接码头、铁路和公路运输的核心设备，其周转时间直接决定了集装箱从船舶到最终目的地的效率。

GS的数据显示，20英尺集装箱底盘的街道停留时间在6月初为4.7天，与上周持平；40/45英尺底盘的街道停留时间为6.2天，较前一周的5.8天略有上升。但更值得关注的是终端停留时间：20英尺底盘从12.4天降至10.8天，40/45英尺底盘从5.5天降至5.2天。这些数字虽然仍有波动，但整体趋势是向下的，且远低于2021年峰值时期的15-20天。

底盘周转时间的改善，意味着港口拥堵的“最后一公里”瓶颈正在被系统性地解决。这背后是港口运营商的设备投资、底盘池的共享机制优化，以及铁路公司对多式联运节点的协调。报告没有详细拆解这些改善的具体来源，但一个合理的推论是，过去几年美国港口和铁路公司的资本开支开始产生回报。

这对航运公司和物流企业的含义是：运力的“硬约束”正在软化。过去，即使船舶运力充足，底盘和铁路的瓶颈也会人为制造拥堵，推高即期运费。如果这些瓶颈持续缓解，那么航运市场的供给弹性将显著增加，运价的高波动性可能随之下降。对于长期投资者而言，这意味着航运股的周期性溢价需要重新定价。

西海岸两

[... middle omitted ...]

升的空间已经有限。

GS的报告没有展开讨论这一点，但它引出了一个关键问题：如果货运量继续增长，铁路公司是否需要进一步增加资本开支来扩容？如果答案是肯定的，那么铁路公司的资本回报率和自由现金流将面临压力。反之，如果铁路公司能够通过数字化调度和资产共享来消化增量，那么其运营杠杆将非常可观。这个问题的答案，将直接影响投资者对铁路股的估值判断。

4. 报告尚未完全回答的关键问题：关税和地缘政治如何改变“新常态”的时间线

GS在报告的开头部分明确提出了这个悬而未决的问题：“关税和地缘冲突将对货运需求、时间节奏以及全球贸易正常化能力产生什么影响？”这是整份报告最诚实也最关键的伏笔。

报告指出，如果供应链压力继续缓解，指数有可能在2026年进入“1”的区间——这意味着比疫情前还要更流畅的流动性。但这一判断的前提是“没有重大外部冲击”。而现实是，美国大选后的贸易政策走向、中东局势对红海航线的持续影响，以及全球供应链“近岸化”和“去风险化”的结构性调整，都可能打破这一线性预测。

\subsection{[R024] JPM：新兴市场债券曲线正在发出被低估的信用分化信号}
{\small\color{refgray}归类：宏观与利率：财政供给与通胀再平衡主导定价}

这份JPM最新的EM USD 10s30s Spread Curve报告，表面上是一份每周例行的利差曲线追踪，但数据背后隐藏着一个被多数投资者忽视的结构性变化：新兴市场债券的期限结构正在从“风险偏好驱动”转向“信用质量驱动”。这意味着，过去那种用主权评级一揽子判断新兴市场久期策略的时代正在结束，投资者需要更精细地拆解单个发行主体的期限溢价。

报告显示，截至2026年6月12日，EM Aggregate 10s30s利差曲线为67个基点，较一周前收窄2个基点，较一年前收窄10个基点。乍看之下，曲线平坦化似乎是全球利率环境趋稳的普遍结果。但JPM的数据揭示，这一平坦化并非均匀分布——不同区域、不同信用等级、不同发行主体之间的曲线斜率分化正在加剧，而这种分化恰恰是市场尚未充分定价的关键变量。

真正值得关注的不是曲线本身的绝对水平，而是曲线斜率的相对变化所反映的信用分层逻辑。当EM Aggregate 10s30s在过去一年收窄10个基点时，EM IG（投资级）收窄了14个基点，而EM HY（高收益）仅收窄3个基点。这种分化不是偶然的噪音，而是市场正在对信用风险重新定价的信号。

JPM的数据清晰地展示了这一分化的程度。EM Aggregate 10s30s当前为67个基点，其内部结构已经出现明显裂痕：EM IG为61个基点，EM HY为79个基点。更值得注意的是，在过去一年中，IG曲线收窄了14个基点，而HY仅收窄3个基点。这意味着，投资级债券的久期溢价正在被系统性压缩，而高收益债券的期限溢价几乎纹丝不动。

这一现象与市场直觉相悖——如果曲线平坦化是全球利率环境趋同的结果，那么IG和HY应该同步收窄。但现实是，市场正在对两类信用质量施加不同的久期定价。对于IG发行人，投资者愿意接受更低的期限溢价来锁定收益率；而对于HY发行人，投资者要求更高的期限补偿来对冲尾部风险。

这种分化的含义是深远的：它表明，新兴市场债券的曲线定价已经从“宏观因子主导”转向“信用个体主义”。过去，投资者可以通过判断美联储政策路径或全球风险偏好来统一调整久期敞口；现在，久期策略必须嵌入信用筛选维度。

报告中最具洞察力的对比来自主权债券与准主权债券的曲线斜率分化。EMBIG Sovereign 10s30s当前为68个基点，而EMBIG Quasi为61个基点。虽然两者在过去一年均收窄了8-9个基点，但主权债券的绝对曲线水平始终高于准主权债券。

更关键的是历史对比。参照Figure 13，在2022年的市场动荡期，主权债券的10s30s曲线曾一度跌至-220个基点，而准主权曲线仅降至80个基点左右。这一极端分化说明，主权债券的曲线对系统性冲击的敏感度远高于准主权债券。当前，两者差距已从极端值回归，但主权曲线仍高出准主权7个基点。

这一差距隐含着一个尚未被充分讨论的风险：主权债券的曲线定价中是否包含了尚未定价的尾部风险？考虑到当前EM Aggregate 10s30s处于历史中位数附近（参照Figure 9，过去五年范围在-80至100个基点之间），主权曲线的相对陡峭可能反映了市场对部分主权发行体信用基本面恶化的预期，而非简单的流动性溢价。

对于投资者而言，这意味着主权债券的久期敞口需要比准主权债券更加谨慎。如果主权信用出现边际恶化，其曲线可能面临比准主权更剧烈的陡峭化，从而对长端多头头寸形成压力。

JPM报告的区域细分数据揭示了另一个重要趋势：不同区域的曲线斜率正在走向不同的均衡水平。EMBIG Asia 10s30s仅为36个基点，CEEMEA为76个基点，LatAm为65个基点。亚洲曲线在过去一年收窄了20个基点，而CEEMEA反而走阔了1个基点。

这种分化不是随机的。亚洲曲线的极端平坦化反映了该区域信用质量的系统性改善——

[... middle omitted ...]

洲，企业债的期限溢价远高于主权债；在拉美，情况则相反。

这一反转对久期策略有直接含义：在亚洲，投资者应该优先选择主权债来获取久期敞口，因为其期限溢价已经被压缩到极低水平；而在拉美，企业债可能提供更优的久期风险收益比。这种区域间的结构性差异，为久期策略的国别选择提供了一个全新的观察维度。

报告中的一周变化数据（Figure 4和Figure 7）提供了最动态的信号。在过去一周，EM Aggregate 10s30s收窄2个基点，而美国国债10s30s走阔3个基点。这意味着，EM曲线与美国国债曲线的走势出现了背离——EM曲线在独立地平坦化，而非跟随美债。

更细致的分析来自Figure 7的一周变化散点图。数据显示，1周内10年期利差的变化与10s30s曲线变化之间并非简单的线性关系。例如，CEMBI Asia的10年期利差走阔了2.5个基点，而10s30s曲线却走阔了3个基点；相反，EMBIG Quasi的10年期利差走阔了2.5个基点，但曲线却收窄了4.5个基点。

\subsection{[R025] JPM：亚洲正在从“被迫加息”转向“主动加息”，市场对此定价不足}
{\small\color{refgray}归类：亚洲央行与区域市场：从被动加息转向主动加息，市场定价不足}

JPM：亚洲正在从“被迫加息”转向“主动加息”，市场对此定价不足

亚洲央行的货币政策叙事正在经历一个被市场普遍低估的转折。JPM最新一期亚洲经济周报提出了一个核心判断：亚洲正从“弱势下的被动加息”阶段，进入“强势增长驱动的主动加息”阶段。这不是一个局部的、边缘的变化，而是一个结构性切换。印尼和菲律宾的加息属于前者——由汇率压力和通胀失控所迫；而韩国、台湾、新加坡正在酝酿的加息则属于后者——由科技繁荣、超预期的增长韧性和通胀传导风险所驱动。市场目前更多关注前者，对后者的定价仍不充分。

这份报告的关键信号在于：当增长本身成为加息的理由，而不是防御性工具时，资产定价的逻辑需要重新校准。对于关注亚洲市场的投资者而言，理解这两类加息的区别，比预测某一次会议的利率决议更重要。

1. 印尼和菲律宾的加息是“防御性”的，但市场可能低估了菲律宾的紧缩幅度

JPM明确指出，亚洲已经出现两类加息。第一类，也是市场相对熟悉的，是“弱势下的被动加息”。印尼央行在美元/印尼盾突破18200历史新高后，意外进行了25个基点的非周期加息，并辅以一系列监管措施应对汇率风险。菲律宾则是区域内通胀压力最大的经济体，尽管5月通胀数据略低于预期，但6.8\%的同比水平仍然令人不安。JPM预计菲律宾央行将在下周加息50个基点，作为锚定通胀预期的预防性措施。

这里的关键洞察不在于加息本身，而在于两点。第一，印尼的加息虽然暂时缓解了汇率压力，但JPM认为压力仍然足以让央行在下次例行会议上再加25个基点。这意味着印尼的紧缩周期尚未结束，市场可能低估了其持续性的风险。第二，菲律宾的加息决策并非板上钉钉。报告坦诚地指出，尽管基线预测是50个基点，但考虑到全球原油价格近期下跌，25个基点的加息也是可能的。这种不确定性本身就意味着，市场对菲律宾央行鹰派程度的定价可能存在偏差。

这两者合在一起的含义是：对于印尼和菲律宾的资产，市场可能已经price in了“一次加息”，但没有充分price in“加息周期的持续性和不确定性”。尤其是菲律宾，如果最终加息幅度低于预期，反而可能引发通胀预期脱锚的更大担忧。

2. 韩国、台湾、新加坡的“主动加息”才是真正的结构变化，市场对此定价不足

报告最具启发性的部分，是提出了“另一类央行正在浮现”。这类央行的加息动机不是防御，而是主动管理。推动因素是科技出口繁荣、中东危机的意外绝缘效果以及财政政策的持续支持，共同导致经济增长持续超预期。

数据本身很有说服力。韩国一季度GDP被上修至7.5\%的同比增速，JPM因此将全年增长预测上调至3.7\%。台湾在科技和非科技领域同步走强的推动下，二季度GDP追踪数据被大幅上修。报告将这种增长形容为“boomy”，这个词本身就传递了超出“稳健”的意味。

JPM的判断是，这种“繁荣式增长”加上通胀风险上升，将促使多家央行考虑从强势地位出发进行加息。韩国央行已经转向鹰派，JPM预计未来12个月将加息100个基点。台湾和新加坡预计将出现类似的转向。台湾下周的央行会议因此值得高度关注——5月CPI略超预期，PPI则飙升14.1\%，这带来了成本传导和利润率挤压的风险。

这里有一个容易被忽略的细节。报告提到，对于台湾下周的会议，“加息仍然是尾部风险事件”，但关键在于“央行是否发出信号，表明能源成本正在传导至核心通胀和通胀预期”。这意味着，即使下周不加息，央行释放的信号本身就会影响市场对年底前加息的概率定价。对于交易员而言，真正的变量不是利率决议本身，而是央行沟通的措辞。

3. 中国出口的“AI科技亮色”掩盖了生产动能的收窄，二季度GDP风险偏向下行

JPM对中国经济的分析同样提供了重要的结构性洞察。5月贸易和通胀数据指向AI科技的角色日益增强。出口超预期，背后是高技术出口的支撑，尤其是存储芯片/

[... middle omitted ...]

业增加值环比仅增长0.5\%，只能部分弥补4月1\%的环比下降。零售销售预计仅环比微增0.1\%，固定资产投资年初至今可能仍处于收缩区间。如果这些预测兑现，二季度增长风险将明显偏向下行。

这一判断对于理解中国宏观叙事非常重要。市场可能被出口的“表面繁荣”所迷惑，但真正的工业生产动能和国内需求仍然疲弱。AI科技带来的价格推升，并不能等同于广泛的工业活动复苏。

4. 二季度疲软反而为下半年提供了政策空间，但外部不确定性正在上升

报告对中国经济的分析并非一味悲观。一个重要的结构性判断是，二季度的部分疲软反映了财政政策支持的阶段性退坡，这反而意味着下半年的“回吐”压力比此前预期的要小。更重要的是，近期的经济放缓可能促使政府在未来几个月加快财政资金的投放节奏。

同时，外部环境存在对冲因素。下半年出口应该受益于依然坚韧的全球需求，以及科技之外资本支出的扩大。但报告也坦诚地指出了风险：美国关税不确定性（转向301/232条款）以及欧盟对华更强硬的立场，都可能部分抵消这些正面因素。

\subsection{[R026] 摩根斯坦利：市场低估的不是AI需求本身，而是需求在亚洲供应链中产生的结构性再定价}
{\small\color{refgray}归类：AI/算力与半导体：从GPU短缺到系统级重构，互联与存储定价权转移}

摩根斯坦利：市场低估的不是AI需求本身，而是需求在亚洲供应链中产生的结构性再定价

这份摩根斯坦利最新的“Three Actionable Ideas”报告，表面上给出了三个来自不同行业的买入建议，但其背后隐含着一个更深刻的判断：当前亚洲市场的定价逻辑，正在从“谁受益于AI需求”切换到“谁能在AI引发的供应链重构中，重新定义自己的议价权”。

报告选取的三家公司——臻鼎、华邦电、ONGC——分别来自PCB、存储芯片和能源上游三个看似不相关的领域。但将它们放在一起审视，一个清晰的图景浮现出来：摩根斯坦利认为，市场对AI主题的定价已经过于拥挤在少数几个环节，而真正的结构性机会，正在向那些被忽视的“瓶颈”和“再定价”环节扩散。这不是一个关于AI需求增速的讨论，而是一个关于供给结构如何被重塑、以及谁能在这一重塑中占据主动的问题。

1. AI互连的物理瓶颈正在从芯片转向传输材料，PCB行业的价值分配将迎来拐点

报告将臻鼎列为PCB领域的首选受益者，背后的逻辑并非简单的“AI服务器出货量增长”，而是一个更具体的产业现象：随着数据传输速率从400G向800G乃至1.6T演进，PCB作为信号传输介质的物理瓶颈正在凸显。摩根斯坦利在另一份相关报告中详细讨论了“Optics Drives the Board”这一主题，其核心判断是：高速互联对PCB的材料、层数、制造精度提出了质变级要求，这不仅是量的增加，更是供应链格局的重新洗牌。

臻鼎被看好的原因，在于其有能力在这一技术切换过程中“capture meaningful share”。这意味着，市场此前对PCB行业的认知——一个依赖产能扩张的周期性行业——正在被颠覆。新的竞争壁垒在于对高频高速材料的理解、与下游光模块和交换芯片厂商的协同设计能力，以及良率控制的技术深度。对于投资者而言，这释放了一个明确的信号：在AI硬件投资中，关注点应从已经高度定价的GPU和HBM，逐步向下游的互连和传输环节延伸，那里可能隐藏着市场尚未充分定价的结构性变化。

2. 存储芯片的“供应问题”正在制造买入窗口，华邦电的短期波动掩盖了中期基本面的走强

报告对华邦电的推荐，基于一个看似矛盾的判断：“memory optimization on servers is more of a supply issue”。这句话的意思是，当前服务器内存的优化调整，更多是由供给端而非需求端驱动的。市场可能将近期存储芯片价格的波动解读为需求疲软，但摩根斯坦利认为，这本质上是供给端的库存消化和产能调配问题。

“Recent share price volatility should provide a good entry point”——这句判断的隐含前提是，市场对供给问题的反应过度，而忽视了华邦电所处的基本面正在改善。报告提到“data points suggest a stronger outlook”，但没有展开具体数据。这里的悬念在于：华邦电的竞争优势到底来自何处？是其在特定利基市场的份额增长，还是成本结构的改善？报告没有完全回答这个问题，但一个合理的推断是，随着AI服务器对内存容量的需求持续增长，以及非HBM类型的存储芯片在边缘计算和终端设备中的应用扩大，华邦电可能正站在一个需求结构多元化的起点上。

3. 能源安全与AI用电需求正在重塑上游能源资产的估值逻辑，ONGC的“增长”属性被严重低估

ONGC的推荐是三个想法中最具反直觉色彩的。摩根斯坦利指出，“energy security and Powering AI increase long-term 'g' for upstream energy producers”。这里的“g”是增长率。传统上，上游能源公司被视为价值股，其增长逻辑依赖于油价

[... middle omitted ...]

 see its market cap in 20 years”。这是一个极具冲击力的表述：它意味着ONGC在过去二十年里向股东返还的资本，已经相当于其当前市值的全部。这暗示着，市场可能仍在用传统的能源周期股框架来给ONGC定价，而忽略了其作为“AI基础设施受益者”和“能源安全战略资产”的双重属性。如果这一重估逻辑成立，那么ONGC的估值框架将需要从“周期股”切换到“成长股+高分红股”的混合模式。

4. 报告尚未完全回答的关键问题：台股与韩国市场的相对表现差异，是否隐含了更深层的产业周期错位？

报告给出的三个想法分别来自台湾、台湾和印度，但如果我们把视野扩大到整个亚洲，一个未被回答的问题是：为什么摩根斯坦利在近期推荐中，台湾科技股（如臻鼎、华邦电、Aspeed、Accton）的密集度远高于韩国？从Exhibit 4和Exhibit 5的业绩记录来看，韩国市场的一些推荐（如三星SDI、Amorepacific）表现显著落后，而台湾的PCB和半导体设备公司则表现强劲。

\subsection{[R027] 摩根斯坦利：AI基建的真正瓶颈不在资本，而在电力与结构性供给约束}
{\small\color{refgray}归类：能源转型与电力基础设施：AI基建的真正瓶颈在电力供给}

摩根斯坦利：AI基建的真正瓶颈不在资本，而在电力与结构性供给约束

市场上关于AI基础设施投资的讨论，大多集中在资本规模、算力需求、以及科技巨头的资本开支计划上。一个被广泛接受的叙事是：只要资金到位，算力就能跟上。但摩根斯坦利最新的一份宏观策略报告提出了一个更冷静的判断——AI基建的真正瓶颈不在融资端，而在供给端，尤其是电力基础设施的结构性约束。这份报告的核心洞察是：电力不再是数据中心建设的“配套条件”，而是与算力建设并列的核心约束变量。理解这一点，才能理解未来几年AI产业链的竞争格局、定价权分布，以及投资逻辑的重塑方向。

报告指出，电力变压器交货周期已从此前的12-16周拉长至平均128周，部分型号甚至达到144周。伯克利实验室的数据显示，截至2025年初，美国电网互联积压项目规模已超过美国现有装机容量的两倍。这些数字背后是一个简单的物理现实：AI驱动的电力需求增长，正在与一个无法快速响应的供给体系发生碰撞。这不是一个边际问题，而是一个系统性问题。

1. 电力约束正在从“配套条件”升级为“独立瓶颈”，数据中心选址逻辑因此重构

报告的核心判断之一是，电力供给的约束已经不再是数据中心建设中的一个“需要协调的变量”，而是成为一个前置性的、独立的瓶颈。过去，数据中心开发商的逻辑是：先确定需求、选址、建设，再协调电力接入。但现在，这个顺序正在被逆转——开发者必须首先确认电力可用性，才能推进项目。

这一变化的影响是深远的。它意味着数据中心的选址将从“靠近用户”转向“靠近电力”。那些拥有富余电力容量、或者能够快速接入电网的区域，将获得显著的竞争优势。同时，这也意味着数据中心建设的节奏将不再由资本开支计划决定，而是由电网扩容和发电设施建设的节奏决定。后者通常以年为单位，甚至以五年为单位。

从这个角度看，市场对AI资本开支的线性外推可能过于乐观。报告虽然没有直接给出具体的建设时间表修正，但它暗示了一个逻辑：如果电力基础设施的瓶颈持续存在，那么算力供给的增长曲线将比需求曲线更平缓、更波动。对于投资者而言，这意味着需要重新评估那些依赖“算力快速扩张”假设的资产定价。

2. 融资模式正在发生结构性变化：AI与能源的资本池正在合并

报告的另一个重要洞察是，AI基建的融资需求和能源基础设施的融资需求正在趋同。传统的融资模式中，数据中心和电力设施是分开的：数据中心由科技公司或数据中心REITs融资，电力设施由公用事业公司融资。但现在，这种边界正在模糊。

报告提到，AI公司本身正在越来越多地直接收购、签约或融资电力资产，而不是将其视为独立的、由公用事业主导的投资。这意味着，AI基建的融资模式正在从“分别融资、分别建设”转向“一体化融资、协同建设”。这种变化的驱动因素很直接：当电力供给成为瓶颈，AI公司必须主动介入，才能确保自己的算力建设计划不被搁置。

这种趋势对资本市场的含义是双重的。一方面，它意味着AI相关融资的总规模可能被低估——因为电力资产的投资需要被纳入AI基建的资本开支框架。另一方面，它也意味着AI公司的资本结构和资产负债表将变得更加复杂，因为它们需要同时承担科技资产和能源资产的双重风险。

3. 结构性约束不止于电力：劳动力、水资源与政策阻力正在形成叠加效应

报告没有止步于电力分析，而是进一步指出，电力只是最显性的约束之一。在劳动力方面，美国未来十年预计将面临约30万电工的短缺，而现有电工队伍中超过五分之一年龄在55岁以上，即将退休。这意味着即使电力设备到位，安装和维护的人力也可能不足。

水资源是另一个正在浮现的约束。标普的分析显示，全球43\%的数据中心位于高水压力地区。数据中心的冷却系统对水资源有显著需求，在干旱频发的区域，这一约束可能成为限制扩张的硬性条件。

政策阻力也在加大。报告提到，纽约州正在审议

[... middle omitted ...]

之间可能出现结构性失衡，而稀缺性将成为市场的核心特征。在这种环境下，拥有规模化、可靠算力供给能力的玩家将获得显著的定价权——报告称之为“算力商人”（merchants of compute）。

更重要的是，报告指出，在当前周期阶段，需求的价格弹性相对较低，尤其是来自企业级应用的需求。这意味着，即使算力价格上涨，需求也不会显著放缓。如果有什么变化，那可能是价格上升会进一步推动算力向高价值应用集中，而低价值应用可能被挤出。

这个判断对投资者的含义是直接的：在算力供给受限的环境中，拥有稳定算力供给能力的公司，其议价能力和盈利能力可能被市场低估。相反，那些依赖“算力价格持续下降”假设的商业模式，可能需要重新审视。

尽管报告对结构性约束的分析非常透彻，但它并没有给出一个清晰的“临界点”判断——即这些约束何时会从“制约因素”变为“硬性天花板”。电力变压器交货周期拉长到128周，这是一个事实，但这是否意味着2026年、2027年的算力建设将显著低于预期？报告没有给出量化答案。

\subsection{[R028] NOM：人民币中间价模型正在发出一个被市场忽视的信号}
{\small\color{refgray}归类：地缘风险与汇率：和平红利下的货币重估与亚洲资产重新定价}

市场对人民币汇率的讨论，大多仍然停留在“中美关系”、“出口压力”或“央行干预”这些宏观叙事层面。但一份来自NOM亚洲外汇策略团队的最新研报，通过其定价模型给出了一个更具体、也更值得推敲的判断：人民币中间价的模型预测值，正在出现系统性偏离。

这份报告的核心不在于预测一个具体的点位，而在于揭示一个机制层面的变化——中间价的形成逻辑，可能正在从“被动跟随”转向“主动引导”。这个变化，对于任何持有人民币资产、或关注中国资本流动的决策者来说，都值得重新审视。

报告提供了一个关键数据点：模型对USD/CNY中间价的预测值为6.7575，较上一期大幅下调了513个基点。而如果计入逆周期因子，预测值则为6.7800，仍较前次中间价低288个基点。这不是一个随机波动，而是一个系统性的、由多个因素共同推动的修正。

下面，我们拆解这份报告的核心逻辑，并尝试回答一个更重要的问题：这个信号，到底意味着什么？

NOM这份报告最值得关注的，不是6.7575这个数字本身，而是模型预测值与实际中间价之间的“误差”。报告中的图2展示了模型误差的历史走势：从2025年初的-1800个基点（模型预测远低于实际中间价），逐步收窄，到2026年初接近零，甚至在今年4月转为正600个基点。

这个误差的演变，本质上反映的是“逆周期因子”的调节力度。当模型预测值显著低于实际中间价时，意味着央行在使用逆周期因子来“托底”汇率，防止人民币过快升值。而当误差收窄甚至转正，则意味着这种调节在减弱，或者方向在转变。

报告给出的最新模型预测（不含逆周期因子）为6.7575，比前次低513个基点。这本身就是一个强烈的信号：在当前的宏观环境下，市场力量本身就在推动人民币升值。而计入逆周期因子后的预测值6.7800，虽然仍低于前次中间价，但幅度已大幅收窄。这说明，央行没有完全放任市场力量，但也不再像过去那样强力干预。

*So what**：对于投资者而言，需要关注的不是人民币会升到多少，而是央行的“容忍区间”是否正在上移。误差收窄意味着政策干预的力度在下降，市场定价权在回归。这是一个结构性变化，而非短期波动。

2. 四个主要货币的贡献，暴露了人民币升值的“非对称”驱动力

NOM报告中的图1拆解了模型预测变化的四个最大贡献货币：欧元贡献了+13个基点，日元贡献了+5个基点，韩元贡献了-14个基点，卢布贡献了-18个基点。

这个分布非常有意思。欧元和日元对人民币中间价的贡献是正向的，意味着这些货币的走强在推动人民币升值。而韩元和卢布的负向贡献，则意味着这些货币的走弱在拖累人民币。

这揭示了一个关键事实：当前人民币的升值压力，并非来自中国自身的单一因素，而是全球货币格局分化的结果。欧元和日元因为各自的货币政策预期或地缘政治因素而走强，人民币被动跟随。而韩元和卢布的走弱，则反映了亚洲出口竞争和能源价格波动的影响。

*So what**：人民币汇率的驱动因素正在变得更加“非对称”。它不再仅仅取决于中美利差或贸易顺差，而是越来越多地受到全球主要货币之间相对强弱关系的影响。这意味着，单纯押注人民币单边升值的策略，风险正在上升。投资者需要同时关注欧元、日元、韩元等多个货币的走势，才能更好地判断人民币的短期方向。

报告最后一部分列出了2026年下半年的关键事件日历。这些事件本身并不新奇，但将它们放在一起，并与中间价模型的信号联系起来，就能看出一个清晰的逻辑链条。

7月底的政治局会议、11月的APEC峰会、12月的中央经济工作会议，以及年底可能的中美领导人会晤——这些事件都指向一个方向：政策层可能在为下半年的一系列重要外交和经济议程，创造一个更稳定、也更有利的汇率环境。

如果模型预测的升值趋势在下半年得到延续，那么中间价可能会在这些关键时间点前后出现更明显的调整。尤其

[... middle omitted ...]

用力度减弱的背景下，央行的“新常态”是什么？

报告展示了模型误差的收窄，但没有解释央行为何选择在这个时间点放松干预。是因为对人民币汇率的信心增强？还是因为干预成本上升？或者是为了应对国际压力？

此外，报告也没有讨论，如果模型预测的升值趋势在下半年加速，央行是否会重新加大逆周期因子的使用力度。从历史上看，央行在汇率快速升值时往往会有“恐高”心理，担心对出口造成冲击。

*So what**：对于深度参与者而言，真正的投资机会在于判断“政策容忍的边界”在哪里。NOM的模型给出了一个很好的起点，但最终的决定因素，仍然是央行对汇率波动的真实态度。这个态度，可能要到下半年的事件密集期才会真正明朗。

基于NOM这份报告，我们可以建立一个简单的观察框架，帮助读者形成自己的判断，而不是被动接受结论。

*第一，关注模型误差的持续方向。** 如果未来几周，模型预测值继续低于实际中间价，且误差没有明显收窄，说明市场升值压力仍在积累。如果误差开始扩大，则可能意味着央行在重新加码干预。

\subsection{[R029] Bernstein：上游资本开支的“超级周期”叙事，市场可能信错了}
{\small\color{refgray}归类：能源与大宗：供需错位结构性固化，上游资本开支超级周期存疑}

Bernstein：上游资本开支的“超级周期”叙事，市场可能信错了

市场正在为一个“上游资本开支超级周期”定价。地缘冲突推高布伦特油价至90美元/桶，霍尔木兹海峡的紧张局势让供给中断的担忧升温，投资者很自然地认为：石油公司会加大投资，一个新的开支上行周期已经启动。但Bernstein这份研报的核心判断恰好相反——这个叙事经不起推敲。该机构认为，2026年全球上游资本开支不仅不会进入超级周期，反而可能出现约400亿美元的收缩。这不是一个关于“会不会投”的问题，而是一个关于“谁在投、投在哪里、为什么投”的结构性判断。市场的定价逻辑，可能错配了供给侧的真正驱动力。

油价上涨确实带来了上游勘探生产公司现金流的显著改善。Bernstein的数据显示，美国大型E\&P公司的经营性现金流自疫情以来持续回升，部分得益于行业整合带来的规模效应。但关键问题在于：这些现金并没有流向资本开支。投资回报率、股票回购和自由现金流稳定性，依然是管理层和投资者共同关注的优先级。报告中的图表清晰表明，即便在2026年一季度地缘冲突推高油价后，那些将更高比例经营性现金流返还给股东的公司，股价表现反而弱于那些保留现金、减少分红的公司。这并非鼓励企业“少花钱”，而是说明市场当前对资本纪律的奖励机制仍然有效。Bernstein的结论是：现金流的改善，不等于投资意愿的恢复。对于上市公司而言，资本开支的决策逻辑已经从“有多少钱投多少”转变为“投多少才能让股东满意”。这个转变是结构性的，不会因为油价短期冲高而逆转。

2. 已发现资源中约三分之一尚未开发，但可转化为供应的比例正在下降

从资源禀赋的角度看，全球并不缺油。Bernstein梳理了过去一个世纪的数据，发现大约三分之一的已发现资源至今仍未开发。问题不在于资源总量，而在于这些资源的“可开发性”。深水和超深水是近年来发现的主要来源，但这些项目的开发周期长、成本高、技术难度大。报告明确指出，深水领域的“黄金时代”可能已经过去，将资源转化为产量的难度在上升。页岩油方面，美国页岩盆地的产量增长已经趋于平台化，生产率的进一步提升越来越困难。油砂项目则面临更明确的环境监管和成本压力。这意味着，即便企业有意愿增加开支，可供选择的高质量、低成本、短周期项目也变得越来越稀缺。资源禀赋的“质量”在下降，而不仅仅是“数量”的问题。对于投资者而言，这意味着未来新增供应的边际成本将系统性上升，但同时也意味着资本开支的效率可能持续走低。

3. 下一轮增长的边际桶，控制权正在从独立E\&P转向国家石油公司

这是报告中一个容易被忽略但极其重要的结构性判断。Bernstein指出，全球上游投资的“边际增长桶”——即那些决定产量增量方向的资源——正越来越多地掌握在国家石油公司和政府主导的企业手中，而不是美国页岩油独立生产商。这对于投资效率、回报率和决策速度都有深远影响。国家石油公司的投资决策往往受到非经济因素的驱动，包括能源安全、就业目标、地缘政治考量等。它们对资本回报率的敏感度低于上市公司，但同时也面临更复杂的政治和监管约束。Bernstein的分析暗示，即便全球资本开支总量有所上升，但由这些主体主导的投资，其转化为产量增长的速度和效率可能低于市场预期。对于关注E\&P股票的投资人而言，这意味着“谁在投”比“投了多少”更重要。那些仍然由独立E\&P主导的盆地和项目，可能拥有更好的资本纪律和股东回报潜力。

4. 最好的地质条件往往位于最差的司法管辖区，资本的地理错配难以解决

报告提出了一个堪称残酷的现实：全球最大的未开发资源，往往集中在财政征收率高、政治不稳定、受制裁或法治薄弱的国家。换句话说，资本最想去的地方（地质条件优越），恰恰是资本最难安全进入的地方。这是一个经典的“地理错配”问题。Bernstein的图表展示了这些未开发资

[... middle omitted ...]

受？如果资本无法流向最佳地质条件，那么全球供给增长的潜力就会受到实质性约束。

一个经常被用来论证“超级周期”的论据是：高油价会刺激钻井活动，从而带动油服公司的订单和收入增长。但Bernstein发现，这次的情况完全不同。报告指出，油服和设备公司尚未实质性上调2026年的业绩指引。钻井活动对油价的敏感度已经下降。更值得警惕的是，上游通胀正在消耗资本开支的实际购买力。企业即使维持名义资本开支不变，实际能够购买的钻井和完井服务也在减少。这意味着，要想保持产量不变，企业需要花费更多的钱。报告中的图表显示，通胀本身已经构成了一个“隐性成本”，它让资本开支的增长看起来比实际产出增长更“虚”。对于投资者而言，这意味着需要区分“名义资本开支”和“有效资本开支”。如果通胀吃掉了一部分增量资本，那么实际新增产能可能远低于市场预期。油服公司的业绩指引没有上调，恰恰说明终端需求并不像油价暗示的那么强劲。

报告尚未完全回答的关键问题：霍尔木兹海峡的溢价能持续多久，以及“规划价格”到底是多少

\subsection{[R030] Bernstein：福克斯收购Roku，战略逻辑成立但价格昂贵，市场低估了广告协同的潜力}
{\small\color{refgray}归类：未被主题摘要引用}

Bernstein：福克斯收购Roku，战略逻辑成立但价格昂贵，市场低估了广告协同的潜力

这笔交易标志着传统媒体公司向流媒体转型的一次关键尝试，但其真正的价值不在于硬件或用户规模，而在于广告库存的重新定价。

Bernstein最新发布的研报，对福克斯（Fox）拟以约220亿美元收购Roku的交易进行了深入剖析。这不是一份简单的交易点评，而是一份关于“在流媒体时代，什么资产才真正具有战略稀缺性”的思考框架。报告的核心判断是：Roku是一个独特的战略资产，但福克斯可能并不需要“拥有”它来实现类似的商业结果。然而，如果广告协同效应能够兑现，这笔交易对福克斯股价的重估意义将远超财务数字本身。

在当前媒体行业估值普遍承压的背景下，这笔交易为投资者提供了一个重新审视“流媒体平台价值”的窗口。市场可能低估了Roku作为广告平台与福克斯优质内容结合后，在广告填充率和定价权上的潜在提升空间。这不仅是两家公司的交易，更是对“内容+分发”商业模式未来走向的一次压力测试。

*以下是这份Bernstein研报的核心洞察与延伸思考。**

许多分析将这笔交易视为福克斯获取流媒体分发渠道、加速数字订阅增长的手段。Bernstein承认这一逻辑存在，但指出“商业合作可能同样有效”。真正让这笔交易变得“战略”而非仅仅是“财务”的，是广告。

福克斯在过去两年新增了超过500家广告客户，其广告库存需求旺盛，CPM持续走高。而Roku拥有美国最大的智能电视操作系统市场份额（以2025年出货量计，约800万台，领先于三星Tizen和亚马逊Fire TV），并积累了大量的第一方用户数据。将福克斯的广告销售能力与Roku的库存和数据结合，意味着福克斯可以更有效地填充Roku的广告位，并针对特定受众收取更高溢价。

*这意味着什么？** 这笔交易的核心并非“让福克斯的内容在Roku上卖得更好”，而是“让Roku的广告位卖出福克斯级别的价格”。如果这一协同效应得以实现，Roku的广告收入利润率将显著提升，而这正是当前市场对Roku独立估值时尚未完全定价的部分。

2. 220亿美元的代价：战略资产从不便宜，但稀释效应不容忽视

Bernstein指出，这笔交易的价格（约220亿美元）相当于Roku NTM EBITDA的约30倍，NTM EPS的约60倍。在当前的媒体和科技估值环境中，这无疑是高昂的。报告明确表示，“战略资产从不便宜”，而“战略”往往就是“昂贵”的同义词。

根据Bernstein的测算，在60\%现金、40\%股票的支付结构下，福克斯的净杠杆率在交易完成后约为2.9倍，依然健康。但每股收益（EPS）在2028年（假设协同效应完全实现前）将面临约10\%的稀释。这意味着一部分短期股东将承受EPS下降的压力。

*这意味着什么？** 这笔交易的财务回报并非立竿见影。Bernstein实际上在暗示，福克斯管理层愿意承受短期的EPS稀释和债务压力，来换取一个长期的结构性转型机会。对于投资者而言，需要关注的不是2028年的EPS数字，而是福克斯能否在 年证明其流媒体业务的增长能抵消传统线性业务的下滑，从而推动估值倍数（P/E）的重塑。

3. 监管风险可能被高估：流媒体市场的竞争格局远比想象中复杂

对于这笔交易可能面临的监管审查，Bernstein持相对乐观的态度。理由在于，无论是福克斯还是Roku，在当前的流媒体和操作系统（OS）市场中，都算不上是主导者。

从流媒体内容看，Netflix、Disney+、Amazon Prime Video等巨头占据主导。从操作系统看，Roku虽然在美国智能电视出货量中领先，但面临着来自Samsung Tizen、LG webOS、Amazon Fire TV、Google Android TV以及Walma

[... middle omitted ...]

交易失败的风险溢价。

Bernstein的研报对广告协同效应给予了高度关注，但并未给出具体的量化数字。报告提到“如果福克斯能证明在订阅和广告增长方面持续取得进展，其估值倍数可能重新设定”，但“现在或许还太早，不宜过于乐观”。

这里存在一个关键未解问题：广告协同效应的规模到底有多大？实现路径是否清晰？福克斯的广告销售团队能否有效渗透Roku的现有广告客户？Roku的第一方数据与福克斯的广告技术栈能否无缝整合？这些问题的答案将直接影响这笔交易最终是“战略胜利”还是“财务负担”。

*这意味着什么？** 市场目前对这笔交易的反应可能过于集中在“买贵了”或“EPS稀释”上，而对“广告协同”这一核心价值驱动因素的规模和可执行性缺乏深入理解。这正是专业投资者可以超越市场共识、形成独立判断的关键领域。

Bernstein在报告中提出了一个有趣的问题：还有谁可能加入这场竞购？报告指出，能真正从Roku的分发规模和广告货币化中受益的公司“名单很短，在160美元/股的价格下更短”。

\subsection{[R031] GS：市场低估了IRA谈判规则对皮下注射制剂的连锁影响}
{\small\color{refgray}归类：生物科技与医疗：并购结构性加速，IRA谈判规则影响被低估}

制药行业对《通胀削减法案》的定价影响已有充分预期，但多数讨论仍停留在IV制剂的谈判框架内。GS在最新一份研报中提出一个被广泛忽视的变量：CMS拟议规则正在将皮下注射制剂纳入2029年药品价格谈判的候选范围。这一政策动向的直接后果是，默沙东的Keytruda Qlex、百时美施贵宝的Opdivo Qvantig和强生的Darzalex Faspro可能面临额外定价压力。GS对此进行了量化分析，结论是即便在最保守的假设下，这些公司 年的EBIT也将受到低个位数百分比的冲击。

市场对这一消息的反应是迅速的——消息公布当日，默沙东下跌2.8\%，强生下跌2.2\%，百时美施贵宝下跌1.6\%，而同期标普医疗保健板块仅下跌0.6\%。但这轮下跌是否已经充分定价？GS的分析框架显示，影响远不止于表面数字。

1. 皮下制剂的纳入并非新问题，但此次规则修订的约束力显著升级

这不是CMS第一次将目光投向皮下制剂。GS在报告中回溯指出，在IPAY 2028谈判类别的草案指南中，CMS曾包含一个关于固定组合制剂的段落，但在最终指南中被删除。当时市场认为这只是一个试探性动作。然而，2026年的拟议规则不仅重新引入这一条款，而且措辞更为具体——明确提到含有两种或以上活性成分或疫苗抗原成分的制剂将被视为一个谈判单位，且特别点名了含有透明质酸酶的组合制剂。

这一变化的意义在于，皮下注射制剂通常是IV制剂的改良版本，本质上是“老药新剂型”。如果CMS成功将其纳入谈判范围，将打破制药公司通过剂型改良来延长专利生命周期、规避价格谈判的传统策略。GS对此的判断是，这不仅仅是针对三种药物的个案，而是对行业定价策略的一次结构性打击。

对于投资者而言，核心问题不是“哪些药会被纳入”，而是“这一先例是否会扩展到其他剂型改良药物”。如果答案是肯定的，那么整个肿瘤药领域的定价预期都需要重新校准。

2. 量化冲击：强生的Darzalex Faspro面临最大绝对风险，但相对影响可控

GS的敏感性分析是最值得仔细研读的部分。在假设净价格折扣36\%的基础上，基于2023年CMS药品支出数据，GS对三家公司的影响进行了测算。

强生的Darzalex Faspro承受的绝对影响最大。GS估算，在被纳入谈判后，该药每年将面临约19亿美元的销售冲击，占强生总销售的1.4\%，占调整后EBIT的4\%左右。考虑到Darzalex Faspro在2029年预计贡献强生Darzalex总销售的约56\%，这一比例并不令人意外。但强生对此的回应是，其立场没有改变——公司此前已公开表示，预计Darzalex Faspro在2034年之前不会受到IRA定价约束。这意味着强生可能认为CMS的拟议规则在最终版本中仍有可能被修改或推迟。

默沙东的Keytruda Qlex面临每年约9亿美元的销售冲击，占Keytruda总销售的约3\%，占默沙东调整后EBIT的约3\%。但这里有两点值得注意。第一，默沙东的回应策略与其他公司不同——公司明确表示，无论Keytruda Qlex是否被纳入谈判，其定价策略是一致的，且已考虑到IV Keytruda自2029年起可能受IRA定价约束以及2028年12月生物类似药进入市场的可能性。第二，GS在报告中提出了一个有趣的判断：如果Keytruda Qlex被纳入谈判，理论上可能从准入和覆盖角度创造顺风——因为定价被锁定后，支付方可能更倾向于将其纳入处方集。

百时美施贵宝的Opdivo Qvantig绝对影响最小，每年约2.7亿美元，占Opdivo总销售的约4\%，占调整后EBIT的约2.5\%。但百时美施贵宝当前正处于关键转型期，Cobenfy等新产品尚在爬坡，任何来自核心产品的额外压力都可能放大市场对其盈利能力的担忧。

3. 分析框架的局限性：三个关键假设

[... middle omitted ...]

的影响将显著放大。

第二，36\%的净价格折扣基于IPAY 2027谈判类别的平均折扣。但2029年谈判的药物类别和数量都与2027年不同，折扣率可能更高。CMS在首次谈判中可能采取较为温和的立场，但随着谈判经验的积累，后续年份的折扣力度可能加大。

第三，Medicare暴露比率基于2023年CMS药品支出数据。到2029年，这些药物的适应症扩展、市场份额变化以及人口老龄化等因素都将改变Medicare的暴露比例。特别是Keytruda和Opdivo正在向早期适应症扩展，而这些适应症的患者群体更年轻，商业保险覆盖比例更高，这可能降低Medicare暴露比率。但反过来，Darzalex的主要适应症是多发性骨髓瘤，患者年龄偏大，Medicare暴露比率可能进一步上升。

这些局限性的存在意味着，GS提供的数字更像是一个下限，而非中枢估计。如果上述任何一个假设被突破，实际影响都将超出报告给出的范围。

4. 报告尚未完全回答的关键问题：公司应对策略的差异将如何改变竞争格局

\subsection{[R032] 摩根斯坦利：生物科技并购不是脉冲，而是结构性的重新定价}
{\small\color{refgray}归类：生物科技与医疗：并购结构性加速，IRA谈判规则影响被低估}

这份报告来自摩根斯坦利SMID Cap生物科技团队的第65期周报。标题“Finger On The Pulse”暗示了团队对行业脉搏的持续追踪。在阅读完整份报告后，我们认为最值得提炼的判断不是某只股票的涨跌，而是：并购正在从偶发事件变为系统性力量，而市场对2026年生物科技资产定价的逻辑，可能仍低估了这一结构性变化的深度。

这份报告发布于2026年6月14日。此时，市场正处在一个微妙的窗口期：利率环境虽未大幅改善，但已趋于稳定；IPO市场在经历了2025年的极度低迷后，2026年年初至今的融资额已经超过了2025年全年；更重要的是，大型药企的并购动作正在加速。摩根斯坦利的团队没有简单地将这些信号归结为“市场回暖”，而是试图拆解这些现象背后真正驱动行业重新定价的力量。

以下是我们从这份报告中提炼出的四个关键层次，以及一个报告尚未完全回答的关键问题。

1. 2026年的并购不是在追赶潮流，而是在填补未来5年的产品线缺口

报告中最具冲击力的图表之一是“BioPharma M\&A Deals Announced in ”列表。短短一个多月内，GSK以106亿美元收购Nuvalent，Incyte以20亿美元收购Vega Therapeutics，Bayer以24.5亿美元收购Perfuse，再加上Eli Lilly在5月26日单日宣布的三笔收购。这些交易的密集程度，已经超出了简单的“行业整合”叙事。

摩根斯坦利团队在报告中给出了一个关键判断：并购是结构性的，而非偶发性的。这句话的分量在于，它指向了大药企正在面临的专利悬崖压力。以Incyte收购Vega Therapeutics为例，报告明确指出，这笔交易将在2029年后对Incyte的增长产生增值作用，而2029年正是其核心产品Jakafi面临专利到期的关键时间点。这不是一次战术性补强，而是战略性的管线重置。

更深一层的含义是，大型药企正在用资产负债表上的现金和低成本的融资，去购买那些经过初步临床验证、但尚未商业化的资产。这种行为模式的变化意味着，中小型生物科技公司的估值锚点正在从“能否独立上市”转向“能否成为大药企管线中的关键拼图”。对于投资者来说，这意味着评估一家SMID Cap生物科技公司的框架需要调整：不应只看其现金跑道和商业化能力，更要看其核心资产是否具备成为“大药企并购标的”的稀缺性和不可替代性。

报告中的IPO数据图表显示，2026年年初至今，生物科技IPO的总融资额已经达到35亿美元，超过了2025年全年的17亿美元。更值得注意的是，平均IPO融资额从2025年的1.9亿美元跃升至3.3亿美元。这是一个值得细读的信号。

摩根斯坦利团队在“Key Takeaway”中写道：“1Q26 has already seen IPO volumes surpass last year. Average IPO proceed is up meaningfully.” 这句话的潜台词是，市场正在重新向那些拥有更成熟数据、更大市场空间的资产敞开大门。2026年至今的10个IPO中，有7个处于临床二期、3个处于临床三期，没有一个是临床前或一期资产。这与2021年泡沫期大量临床前公司蜂拥上市形成了鲜明对比。

这种结构性变化意味着，投资者对生物科技的风险偏好正在修复，但修复的方向是“质量优先”。那些拥有清晰临床路径、明确监管路径和可观市场规模的资产，正在获得资本市场的溢价。而那些仍停留在概念验证阶段的资产，即便市场回暖，也未必能获得同等的融资机会。对于产业决策者而言，这个信号指向了一个更务实的策略：与其追求IPO的窗口期，不如先确保核心资产进入临床二期或三期，再考虑公开市场融资。

3. 免疫重置疗法的数据验证了机制，但持续性的验证才是定价的关键

报告

[... middle omitted ...]

alidating the mechanism in SLE and RA... but confirmation of the immune-reset potential and durability of response remains the crux of investor focus.” 这里的关键词是“immune-reset potential”（免疫重置潜力）。KOL（关键意见领袖）的评论进一步揭示了市场的担忧：大约3-6个月的持续性B细胞清除，是实现免疫重置效应的关键阈值。而当前的数据尚无法证明这一点。

这个案例揭示了当前生物科技定价中的一个核心矛盾：机制验证（Proof of Mechanism）已经不再是推动股价上涨的充分条件，投资者正在要求更持久、更可重复的临床终点数据。对于任何一家处于类似阶段的公司来说，单次治疗后的短期应答已经不够，市场需要看到的是“治疗性停药后的持久缓解”。这个门槛的提高，正在重塑整个自身免疫疾病治疗领域的估值逻辑。

\subsection{[R033] 摩根斯坦利：天然气轮市场的真正拐点不在2026，而在供给侧的隐性再定价}
{\small\color{refgray}归类：能源转型与电力基础设施：AI基建的真正瓶颈在电力供给}

摩根斯坦利：天然气轮市场的真正拐点不在2026，而在供给侧的隐性再定价

这份摩根斯坦利关于西门子能源的最新研报，表面上是在讨论天然气轮订单周期何时见顶，但真正值得产业决策者和投资者关注的，远不止于一个时间点的判断。

报告的核心信号是：市场对天然气轮需求侧的叙事已经高度拥挤，但供给侧的结构性变化——尤其是非传统发电解决方案的产能扩张——正在悄然重塑2030年的竞争格局。这种变化不是线性外推的，而是通过“时间成本”和“替代弹性”两个隐形变量在发挥作用。

摩根斯坦利分析师在密集的投资者交流后发现，即便是最专业的买方，对供给侧的认知也明显滞后于需求侧。超过50份索要供给模型的请求，本身就说明了一个问题：市场的注意力分配严重不均。而正是这种不均，可能孕育着未来最大的预期差。

以下是我们从这份研报中提炼出的五个关键洞察，以及它们对行业格局和投资框架的深层含义。

1. 市场共识正在从“订单还能涨多久”转向“利润弹性还有多大”，但后者的边际改善空间正在收窄

摩根斯坦利维持西门子能源的超配评级，但态度明显比 年更谨慎。核心原因在于：盈利上调的“速率”正在放缓。

从研报提供的图表可以清晰看到，2024年8月到2025年11月，每次业绩发布后30天内，市场对2028年EBITA的共识上调幅度都在7\%以上，最高达到16.2\%，伴随的是股价15\%-30\%的上涨。但进入2026年，上调幅度骤降至5\%左右，股价反应也变得平淡甚至为负。

这背后的逻辑很直接：当一家公司的盈利预期已经大幅上修后，继续超预期的难度呈指数级上升。摩根斯坦利目前对西门子能源2030年EBITA的预测仅比共识高5\%，而在 年，他们对2028年EBITA的领先幅度曾高达20\%。

这意味着，即使11月11日发布的新2030目标能触发新一轮共识上修，其幅度和股价弹性也将远低于此前几轮。对于持有者而言，真正的考验不是2026年订单是否见顶，而是当盈利上调的“加速度”消失后，市场愿意给什么估值。

2. 供给侧的真正风险不在传统燃气轮机巨头，而在“非传统解决方案”的产能涌入

摩根斯坦利的供给模型显示，到2030年，燃气轮机加电力解决方案的总供给能力可能达到135GW，其中传统燃气轮机约97GW。投资者对97GW这个数字基本没有异议，争议集中在剩余的38GW——包括发动机（如卡特彼勒、瓦锡兰、现代）、燃料电池（如Bloom Energy）等非传统方案。

反对者的逻辑是：这些“过渡方案”最终会被更高效的燃气轮机取代，因此不会对长期供需格局构成实质威胁。

但摩根斯坦利的反驳值得深思：这些非传统方案之所以能获得订单，恰恰是因为它们解决了数据中心“时间成本”这一核心矛盾——燃气轮机交付周期长，而发动机和燃料电池可以更快部署。问题在于，当大量这类产能同时涌入同一个“表后市场”时，竞争格局会发生质变。

研报明确警告：这将对定价和盈利能力产生“下游影响”。更关键的是，这种竞争压力首先会冲击的不是GE或西门子能源，而是瓦锡兰这类以发动机为主的公司——摩根斯坦利维持瓦锡兰的低配评级，正是基于这一判断。

对于投资者来说，这意味着：不能简单用“燃气轮机订单依然强劲”来否定供给侧风险。风险正在从传统OEM的产能规划，转移到那些不被市场视为直接竞争对手的替代方案上。

3. 数据中心的实际建设速度，正在成为比“算力需求”更硬的约束

投资者普遍认为，算力需求的爆发性增长会直接转化为燃气轮机订单的持续增长。但摩根斯坦利提出了一个更务实的质疑：订单可以下，但项目能建完吗？

报告指出，数据中心的建设速度最终取决于供应链中最薄弱的环节——目前这个环节已经不再是燃气轮机供应，而是工程总承包、劳动力和许可审批。大量的表后解决方案已经下单，但其中相当一部分可能无法在2030年前实际建成。



[... middle omitted ...]

”

当燃气轮机叙事出现裂痕时，很多投资者转向电网业务作为支撑估值的新支柱。但摩根斯坦利认为，电网业务的积极惊喜也在减弱。

数据很说明问题：市场对西门子能源电网技术业务2029年EBITA的共识预期已升至50亿欧元，较2025年初翻了约一倍。2030年的利润率共识为23.5\%，而摩根斯坦利自己的预测是25\%——差距已经很小。订单量从2021年的73亿欧元飙升至2026年的约250亿欧元，共识还假设 年继续以中个位数增长。

摩根斯坦利对此有两个关键质疑：第一，欧洲电网建设的大部分订单已经进入账本，后续更多是高位稳定而非持续加速；第二，由于产能约束，该业务的收入增速可能在 年出现放缓，而此前三年是约20\%的复合增长。

更深层的判断是：电网业务不应该被赋予与燃气轮机服务业务或施耐德、伊顿等电气化公司相同的估值倍数。原因在于其进入壁垒较低、利润率波动更大、且缺乏有粘性的后市场收入。

这意味着，如果投资者因为看好电网而给予西门子能源更高估值，这个逻辑本身可能就存在缺陷。

\subsection{[R034] GS：市场误读了美联储的利率路径，最不确定的不是加息，而是何时降息}
{\small\color{refgray}归类：宏观与利率：财政供给与通胀再平衡主导定价}

GS：市场误读了美联储的利率路径，最不确定的不是加息，而是何时降息

当通胀数据连续数月高于3\%的门槛，而劳动力市场又展现出意外的韧性时，市场对美联储下一步行动的定价几乎自动滑向了“加息”。这是过去一年反复出现的市场本能反应。但GS在6月FOMC会议前夕发布的最新研报给出了一个反直觉的判断：加息的可能性被市场系统性高估了，真正的政策博弈焦点不是“加不加”，而是“何时降”以及“降多少”。

这份研报的核心价值不在于预测点阵图，而在于它系统性地拆解了美联储不会因为供给冲击而加息的逻辑框架。GS认为，当前的通胀脉冲——主要由关税、油价和存储芯片价格构成——本质上是一次性的价格水平调整，而非需求驱动的持续性通胀。而美联储历史上对供给冲击的反应模式、当前劳动力市场的平衡状态，以及通胀预期并未脱锚的事实，共同指向了一个结论：加息的门槛远比市场想象的更高。

对于资产定价而言，这意味着市场可能正在为一种低概率情景（加息）支付过高的风险溢价。如果GS的框架成立，那么当前债券收益率曲线中包含的加息预期，将在未来几个月被逐步修正，从而为风险资产提供一个新的定价锚。以下是对这份GS研报关键逻辑的拆解。

GS的分析起点，是最近几个月就业增长的回升。数据图表显示，在经历了2025年下半年的极度疲软（甚至月度净负增长）之后，2026年初的就业数据明显反弹，1月至5月的月度新增非农就业从个位数回升至50-130万的水平。

市场很容易将这种反弹解读为经济过热或需求过强，从而推断通胀压力难以消退。但GS提出了一个更细致的观察框架：他们计算了一个“潜在就业趋势”指标，该指标并非简单看月度数据的波动，而是通过加权平均（0.75乘三个月均值加0.25乘九个月均值）来平滑噪音。这个趋势指标在2026年5月仅为3.5万/月，远低于名义新增数据所暗示的水平。

这意味着什么？意味着近期就业的“反弹”更多是统计上的均值回归，以及前期被压抑的招聘需求的集中释放，而非新一轮强劲扩张周期的开启。失业率从2025年中的4.1\%小幅上升至2026年初的4.5\%，也印证了劳动力市场正在从极度紧张状态向更平衡的状态过渡。

*这里的关键洞察是：一个趋于平衡的劳动力市场，不会产生工资-物价螺旋式上升的压力。** GS的数据显示，其编制的工资追踪器与劳动力市场闲置指标（Slack Tracker）正在向历史均值收敛。当工资增长不再加速，那么由成本推动的通胀传导链条就被切断了。这意味着，当前的通胀问题，本质上不是由劳动力市场过热引发的。

2. 这轮通胀脉冲是“加税”而非“发烧”，美联储不会用加息来应对

这是整份研报最具颠覆性的逻辑链条。GS明确指出，2026年核心PCE通胀之所以持续高于3\%，主要来自三个供给侧因素的叠加：关税、油价和存储芯片价格。这些因素在研报中被量化为对通胀的贡献度——关税贡献了约0.6个百分点，能源贡献约0.1个百分点，软件与配件贡献约0.1个百分点。

GS的核心论据是：美联储倾向于不对供给冲击引发的通胀进行加息。 他们展示了一组历史回归数据：当油价上涨是由于供给因素（而非需求）时，美联储做出强硬表态的概率仅为2.5\%，而欧洲央行则为36.5\%。美联储更关注的是需求侧的通胀压力，以及通胀预期是否失控。

这背后的逻辑非常清晰：供给冲击带来的是一次性的价格水平跃升。例如，关税提高导致进口商品价格一次性上涨，之后同比效应就会消退。加息无法降低关税，也无法压低石油供给短缺导致的价格。相反，加息只会通过抑制需求来制造经济衰退，从而以更大的代价去对抗一个本就会自行消退的通胀脉冲。

GS的预测模型显示，关税对通胀的贡献将在2026年下半年达到峰值（约0.84\%），随后在2027年迅速回落至0.34\%以下。综合来看，核心PCE通胀预计将在2026年全年维持在

[... middle omitted ...]

ommon Inflation Expectations Index）在2026年第二季度约为2.3\%，与美联储的长期目标基本吻合。尽管短期通胀数据高企，但市场、消费者和专业人士的长期通胀预期并未失控。这是美联储敢于按兵不动的最重要底气。

另一个更直观的指标是“高通胀扩散度”——即有多少细分品类的价格涨幅超过4\%、6\%或8\%。在2023年，有超过60\%的品类涨幅超过8\%，那是真正需要强力紧缩的时期。而到了2027年，GS的预测显示，涨幅超过8\%的品类占比将从峰值回落至约20\%，甚至低于 年的正常波动区间。这意味着通胀正在从“广泛上涨”回归到“局部上涨”，这是通胀压力趋于缓和的典型信号。

*这两个指标合在一起，构成了美联储的“安全网”。只要通胀预期不脱锚，且通胀压力不扩散为全面性的价格上涨，美联储就没有动机去打断当前的经济复苏。** GS据此判断，加息的概率仅为20\%，且这20\%的概率主要对应的是“更高通胀、更高增长、更高终端利率”的情景，而非传统意义上的紧缩周期。

\subsection{[R035] Citi：AI推理层正在从“成本竞赛”转向“稀缺性定价”，市场低估了中间层的结构性机会}
{\small\color{refgray}归类：AI/算力与半导体：从GPU短缺到系统级重构，互联与存储定价权转移}

Citi：AI推理层正在从“成本竞赛”转向“稀缺性定价”，市场低估了中间层的结构性机会

这份来自Citi的最新AI行业周报，揭示了一个正在发生的关键转变：AI基础设施的定价逻辑，正在从“算力成本下降驱动应用爆发”的线性叙事，转向一个更复杂、也更有利可图的结构——稀缺性正在被货币化，而货币化的速度超过了稀缺性被解决的速度。

这不是一个关于“大模型军备竞赛”的故事，而是一个关于“中间层如何捕获价值”的故事。Citi用六周的数据追踪，给出了一个清晰的信号：市场对AI的关注点，应该从“谁的模型最强”转向“谁在控制推理的分配权”。

Citi的数据显示，最前沿模型的推理价格在六周内几乎翻倍，而性能只提升了4个点。这并非技术进步放缓，而是一种主动的定价策略。报告明确指出，模型提供商已经开始将“配给”作为产品特征——新的顶级模型在6月22日后被置于使用积分之后，这意味着用户无法无限量按需调用。

与此同时，上一代A100 GPU的租赁价格在过去六周内上涨了11\%，而高端B200/B300的价格也在稳步攀升。这与市场普遍预期的“算力成本持续下降”形成了鲜明对比。

*这意味着什么？** 算力不再是纯粹的“commodity”，而是正在被分层定价。稀缺性（无论是物理GPU的稀缺，还是顶级模型推理能力的稀缺）正在成为定价的决定性因素。对于依赖推理成本假设来构建商业模型的投资者和企业决策者来说，这是一个需要重新校准的关键变量。

2. 最前沿的模型正在主动拉大与开源模型的差距，而非被动降价

一个反直觉的数据是：专有模型与开源模型之间的智能差距，在六周内从6个点扩大到了9个点。这意味着，前沿模型提供商（如报告提及的Google等）选择在“智能”这个维度上建立结构性壁垒，而不是在价格上与开源模型进行消耗战。

同时，中位数的推理速度在六周内从64 tok/s提升到了105 tok/s，而中位数的延迟则下降了7\%。速度正在向中端市场聚集，而智能则向上游集中。 这形成了一个清晰的“智能-速度-价格”三角分化：顶级模型卖的是智能，中端模型卖的是速度，低端模型卖的是价格。

*对于企业用户而言，这意味着“一刀切”的模型选择策略将失效。** 不同任务需要匹配不同的模型，而能够高效完成这种匹配的中间层——即报告所称的“推理编排层”——将获得巨大的价值。

Citi报告中最具洞察力的判断之一是：由模型分层定价创造的租金和节省，正在向“拥有推理编译和性能剖析数据”的中间层集中。这个中间层知道：哪个模型、在何种量化水平、在何种硬件上、对哪个子任务能产生可接受的输出。

这本质上是一个“路由”问题。报告指出，挑战在于如何在捕获这些数据并实现高效路由的同时，不泄露企业IP或PII。这是一项极高的技术壁垒和数据壁垒。

*这意味着，未来的AI竞争，将从“谁拥有最好的模型”转向“谁拥有最好的路由数据”。** 拥有跨模型、跨硬件的性能剖析数据，并能够基于这些数据做出实时路由决策的中间层，将拥有对上游模型提供商和下游企业的议价权。Databricks和AWS的峰会，正是观察这一趋势如何落地的关键窗口。

4. 数据中心的选址逻辑正在发生根本性变化：电力从运营成本变为资本成本

Citi报告提供了两个关于数据中心选址的关键图表。第一，数据中心集群集中在零售电价9-12美分/千瓦时的区域。第二，可再生能源占比超过25\%的州吸引了不成比例的数据中心容量。

这背后是一个更深层的变化：电力成本正在从“建成后的运营成本”转变为“建设前的资本成本”。 报告引用了一个案例：一家私有云服务商签约了4.9GW的电力，但其总规划管道超过40GW。12\%的签约/规划比例，意味着这些管道更像是开发商持有的“期权”，而非确定的建设计划。

*对于投资者而言，这意味着数据中心的真实供给弹

[... middle omitted ...]

IP和PII的保护要求，可能限制数据的跨客户流动，从而削弱网络效应。此外，模型本身也在快速迭代，今天积累的路由数据，是否适用于明天的模型，也是一个开放问题。

*这个未解问题，直接关系到“路由层”是否是一个可持续的投资主题，还是一个阶段性套利机会。** 这需要更深入的竞争格局分析和数据积累观察。

Citi报告提出了一个简洁而有力的分析框架：“智能-速度-价格”三角。 目前没有一家提供商能够同时占据三个顶点。前沿模型在智能上领先，中端模型在速度上优化，开源模型在价格上竞争。

对于产业决策者和投资者，这个框架提供了一个持续跟踪的坐标： 如果某个提供商（例如Google的Gemini 3.5 Pro，报告称其“仍在事件视界上”）能够同时提升智能和速度，同时不显著推高价格，那么它将打破现有的平衡，重塑竞争格局。 如果智能差距持续扩大，那么企业将更愿意为顶级模型支付溢价，推理编排层的价值将进一步提升。 如果开源模型在智能上取得突破，那么整个“智能溢价”的定价逻辑将受到挑战。

\subsection{[R036] JPM：先进核能真正的拐点不是技术，而是客户愿意为第一座电站付费}
{\small\color{refgray}归类：能源转型与电力基础设施：AI基建的真正瓶颈在电力供给}

JPM：先进核能真正的拐点不是技术，而是客户愿意为第一座电站付费

当Meta选择TerraPower的Natrium反应堆作为其数据中心供电方案时，这件事的意义远超一份商业合同。JPM近期与TerraPower管理层举行了一场炉边对话，这份研报揭示了一个正在被市场低估的结构性变化：先进核能的商业化路径，已经从“政府补贴驱动”转向“科技巨头需求驱动”。

这不是关于反应堆技术是否可行的问题。这个问题在学术界和工程界已经争论了二十年。真正的新信号是，一家年营收超过1600亿美元的科技公司，愿意在电站尚未建成、技术尚未大规模验证的阶段，就投入里程碑式的开发资金。Meta的选择不是对技术的赌注，而是对电力供给结构必须重构的判断。

JPM这份报告的价值不在于它对TerraPower的技术细节做了多少描述，而在于它清晰地展示了先进核能商业模式正在经历的三个根本性转变：客户结构从公用事业转向科技巨头，融资结构从政府拨款转向企业预付，竞争焦点从技术参数转向供应链整合能力。理解这些转变，才能理解为什么这次可能真的不一样。

以下是我们从这份研报中提炼出的核心洞察，以及它对产业竞争格局和投资逻辑的含义。

1. 科技巨头正在成为先进核能的第一推动力，但只有头部玩家能承受FOAK风险

Meta的参与方式值得仔细拆解。JPM报告中提到，Meta经过严格的RFP流程后选择了Natrium技术，并且不是简单地承诺未来购电，而是提供“里程碑式开发资本”——这笔资金将覆盖TerraPower在最终投资决策前的FEED、选址和监管工作，周期大约两年。

这个结构的意义是什么？它解决了先进核能商业化中最棘手的“先有鸡还是先有蛋”问题。第一座电站（FOAK）的成本和风险远高于后续电站，因为供应链尚未建立、监管路径尚未验证、施工经验尚未积累。传统上，这个风险要么由纳税人承担（通过政府补贴），要么由公用事业公司转嫁给消费者。Meta的模式提供了一种新的选择：由最终用户直接承担部分前期风险。

但这份报告也诚实地指出了边界条件。JPM提到，除了Meta这类超级科技公司，其他潜在客户如“neo-clouds”和芯片供应商虽然表现出兴趣，但“缺乏同样的资产负债表强度”。公用事业公司也有兴趣，但面临“成本超支风险和费率承受能力约束”。

这里的含义很清楚：先进核能的第一波商业化，将由极少数拥有超强资产负债表且电力需求急迫的企业推动。这不是一个普惠性的市场启动，而是一个高度选择性的过程。对于投资者而言，这意味着FOAK阶段的投资机会将高度集中在少数几个项目上，而不是整个行业普涨。

2. 供应链瓶颈的真正解法，在于把核能制造能力从西方转移到韩国

JPM报告中最被忽视的信息，可能是TerraPower与韩国合作伙伴的深度绑定。报告提到，KHNP、SK和Hyundai不仅持有TerraPower的股权，而且它们旗下的关联公司“在全球范围内建造核电站”。

这不是一个简单的供应商关系。韩国是全球少数几个仍然具备完整核电站建造能力的国家之一——KHNP是韩国核电的运营主体，Hyundai是EPC承包商，SK是综合能源集团。这三家公司的组合，意味着TerraPower的Natrium反应堆可以在韩国完成模块化制造，然后运往美国现场组装。

这个策略的价值在于，它绕过了美国本土核能供应链的瓶颈。美国在过去二十年几乎没有新建核电站，导致熟练工人、制造设施和供应链管理经验严重流失。韩国则一直保持核能建设能力，包括APR1400反应堆在阿联酋的出口成功。把制造环节放在韩国，相当于直接继承了世界上最成熟的核能工程管理体系。

但这里有一个尚未被充分讨论的问题：地缘政治风险。如果中美关系或朝鲜半岛局势发生变化，依赖韩国供应链的先进核能项目是否会面临中断风险？JPM报告没有展开

[... middle omitted ...]

RDP奖项，为两个反应堆的开发预留了32亿美元，但DOE迄今只分配了其中20亿美元。剩余的12亿美元将在50/50的成本分摊结构下逐步释放。这意味着TerraPower需要自己筹集至少12亿美元才能匹配政府资金，而这仅仅是第一座电站。

对于其他没有获得ARDP奖项的小型模块化反应堆开发商而言，融资挑战只会更大。这解释了为什么JPM在报告中反复强调“里程碑式开发资本”的重要性——在FOAK阶段，谁能找到愿意预付资金的客户，谁就能在竞争中占据决定性优势。

JPM报告在末尾提到一个被大多数市场参与者忽视的细节：TerraPower的医疗同位素业务，管理层将其定位为今年约10亿美元的收入机会。

这个数字如果准确，意味着TerraPower在核电站建成之前，就已经拥有一个具有规模意义的收入来源。Actinium-225是一种用于靶向癌症治疗的放射性同位素，自然界中不存在，只能通过人工生产。TerraPower从爱达荷国家实验室的遗留武器级废料中提取钍，然后生产这种同位素。

\subsection{[R037] 摩根斯坦利：市场正在严重低估HDD周期的长度与定价弹性}
{\small\color{refgray}归类：AI/算力与半导体：从GPU短缺到系统级重构，互联与存储定价权转移}

这份来自摩根斯坦利的研报，核心判断只有一句话：硬盘驱动器（HDD）的繁荣周期远未结束，恰恰相反，市场对它的长度和盈利弹性都严重低估了。

报告基于过去三周在亚洲的密集调研——包括台湾AI峰会、Computex展会以及西部数据亚洲路演——得出了一个清晰的结论：HDD的供需缺口正在扩大，定价权正在向供应商转移，而这一趋势至少将持续到2028年。

摩根斯坦利因此大幅上调了希捷科技（Seagate Technology）和西部数据（Western Digital）的盈利预测与目标价。希捷目标价从767美元上调至1035美元，西部数据从488美元上调至650美元。更重要的是，报告指出，在乐观情景下，这两家公司可能在2025至2028年间实现EPS的10倍增长。

这不是一个关于“周期高点”的故事。这是一个关于结构性供需失衡、供应商理性定价以及AI推理需求意外拉动传统存储的故事。对于关注科技硬件和AI基础设施的投资者来说，这可能是当前最具预期差的机会之一。

1. 供需缺口正在扩大，而不是收窄——短缺至少延续到2028年

市场对HDD周期的一个常见误判是认为这轮需求高峰已经过去。摩根斯坦利的最新调研完全推翻了这一假设。

报告明确指出，HDD需求正在以每年40-50\%的速度增长，而供给（以EB计）的增长仅为30-35\%。这意味着供需缺口不仅存在，而且在持续扩大。具体数据是：2026年近线HDD供给将比需求少300EB（约10-15\%的缺口），2027年和2028年缺口将进一步扩大至400EB。

为什么需求如此强劲？报告给出了两个关键驱动因素。第一，核心云服务的增长超出了预期。第二，AI推理和智能体（agents）的需求正在拉动通用服务器出货量，进而带动HDD需求。值得注意的是，调研显示超大规模云服务商目前部署HDD的比例接近100\%，而历史平均水平约为70\%。这意味着云服务商手中几乎没有库存缓冲。同时，ODM厂商的HDD库存仅有1-2周。整个供应链处于极度紧张的状态。

供给端则受到严格约束。HDD供应商没有新建产能的计划。过去几年行业经历了痛苦的整合与产能出清，供应商已经形成了高度理性的资本纪律。这意味着未来2-3年近线EB的增长将完全由供给决定，而供给增速被锁定在30-35\%的CAGR。需求增速远超供给增速，结果是短缺的持续加深。

这一判断的含义是：HDD的景气周期不是几个季度的事情，而是一个将持续至少三年的结构性机会。市场若仍以“周期股”的框架来估值HDD公司，很可能错过了最核心的定价逻辑变化。

2. 定价弹性被严重低估——每TB价格可能翻倍，而增量利润几乎全部落到底线

供需缺口只是第一层逻辑。真正让摩根斯坦利大幅上调盈利预测的，是定价端的巨大弹性。

报告提供了几个关键数据点。目前希捷和西部数据的混合近线HDD每TB价格在 美元。今年4月，调研显示供应商可能在 年实现每TB 20美元的价格。而最新的台湾调研显示，供应商的目标是每TB 25-30美元。更值得注意的是，现货市场（非长期协议客户）的HDD价格已经比此前高出约30\%，部分分销商的售价已达到每TB 30-35美元。

这些数字意味着什么？摩根斯坦利做了清晰的量化分析。如果混合近线每TB价格在2027年达到25美元、2028年达到30美元——与台湾调研一致——希捷在2027年的EPS可能超过70美元，2028年超过100美元；西部数据2027年EPS可能超过40美元，2028年超过70美元。作为参照，这两家公司2025年的EPS分别约为10.24美元和6.95美元。这意味着三年内EPS可能增长10倍。

需要强调的是，报告并没有把这种极端情景作为基准。摩根斯坦利的新基准预测是2027年每TB价格约19-19.5美元，2028年约22美元。即便如此，

[... middle omitted ...]

接受价格，而是主动管理定价权和商业条款。

报告揭示了一个非常有意思的细节：虽然超大规模云服务商希望签订长达2032年的长期协议（LTA），但HDD供应商并不愿意签署这么长的合同。原因有三。第一，HDD的应用场景（尤其是冷存储）正在拓宽，供应商对未来需求的可见性在增强，因此不愿意过早锁定不利的价格。第二，短缺在加剧，TCO（总拥有成本）相对于NAND的竞争优势在显著改善。第三，HDD的替代技术目前不存在。

这意味着供应商正在利用供需紧张的结构性优势，主动控制商业条款。他们愿意为客户提供5年的产能和技术路线图可见性，但不急于在12个月以上锁定商业条件。这种策略的核心目的是在未来几年内逐步实现更高的定价，而不是一次性锁定在可能偏低的水平。

这背后反映的是HDD行业竞争格局的根本性变化。经过多年的整合，希捷和西部数据已经形成了高度理性的双寡头格局。报告指出，HDD供应商在供给端极为自律，没有任何新建产能的意愿。这种供给纪律是过去十年行业痛苦的产物，也是当前定价权的基础。

\subsection{[R038] 摩根斯坦利：欧洲软件股正经历近二十年最深的估值回撤，但市场忽视了结构性分化}
{\small\color{refgray}归类：消费与科技硬件：AI眼镜渠道扩容与印度混动拐点}

摩根斯坦利：欧洲软件股正经历近二十年最深的估值回撤，但市场忽视了结构性分化

过去二十年，欧洲软件板块经历过两次深度回撤：一次是全球金融危机，一次是疫情后的利率与估值压缩。现在，第三次正在发生，而且持续时间之长、幅度之大，已经逼近前两次的水平。

摩根斯坦利欧洲软件团队在最新一期周报中，用一张跨越二十年的等权重股价走势图，揭示了一个被市场情绪掩盖的关键事实：自2021年高点以来，欧洲应用软件板块已经连续近1700天没有创出新高，累计跌幅超过35\%。这个回撤的深度和长度，只有2008年金融危机和2022年加息周期可以比拟。

但这份报告真正值得注意的判断，不是“软件股跌了很多”，而是“跌得如此之深，恰恰说明市场正在对软件资产的定价逻辑进行根本性重置”。而重置过程中，不同子板块的分化才刚刚开始。

1. 当前软件股的回撤深度与时长已与金融危机相当，但触发机制完全不同

摩根斯坦利选取了2005年1月之前上市的50余只中大型软件股，以等权重方式构建了一个长期价格指数。结果非常直观：这批股票在2021年底达到历史峰值后，至今未能突破新高。截至报告发布，距离上一次历史高点已经超过1700天，指数从峰值回撤超过35\%。

历史对比更具说服力。全球金融危机期间，软件板块同样经历了深度回撤，但彼时的触发因素是系统性金融风险蔓延至实体经济，软件公司的基本面遭受了广泛冲击。2022年的回撤则更多是利率上升带来的估值压缩，市场在重新定价未来现金流的折现率。

而本轮回撤的背景完全不同：AI技术浪潮正在重塑整个软件行业，企业软件支出并未出现断崖式下滑，多数头部公司的营收仍在增长。这意味着，当前的价格下跌并非基本面崩溃的结果，而是市场对“软件公司能否将AI转化为可持续利润”这一命题的深度怀疑。

这种怀疑有其合理性。过去两年，市场对AI概念的追逐曾将软件股推至不合理的高位，但随后发现，多数公司的AI变现路径并不清晰。摩根斯坦利的分析暗示，当前的回撤更像是一种“预期修正”，而非“盈利危机”。

更重要的是，这种修正并非均匀分布。报告注意到，在应用软件整体疲软的背景下，数据基础设施和网络安全子板块已经出现明显的相对强势。这引出了下一个核心问题：市场是否正在对不同软件资产的商业模式进行差异化定价？

2. 数据基础设施板块的强势反弹，揭示市场正在重新定义软件的价值锚点

摩根斯坦利在报告中特别指出，尽管应用软件整体承压，数据基础设施板块却出现了显著的资金流入。Snowflake在过去一个月上涨约58\%，MongoDB上涨约15\%，Datadog上涨约17\%。这些涨幅并非孤立事件，而是与Databricks新一轮融资传闻同步发生的——据The Information报道，Databricks正在洽谈以超过1650亿美元估值融资，较2月份的1340亿美元估值大幅跳升。

这一现象的深层含义在于：市场正在将软件公司的价值锚点从“用户数量”和“订阅收入”转向“数据资产的稀缺性与流动性”。数据基础设施公司本质上是在构建企业数据的“管道”和“仓库”，它们不直接面对终端用户，但控制着AI时代最关键的资源——结构化与非结构化数据的存储、处理与调用能力。

相比之下，应用软件公司面临的质疑更为尖锐：当AI能够自动生成代码、处理税务、管理客户关系时，传统应用软件的护城河在哪里？摩根斯坦利在报告中引用的伦敦科技周调研数据也印证了这一点——多数英国企业员工尝试过AI工具，但日常使用率、培训覆盖率和可量化的生产力提升仍然有限。这意味着，企业级AI应用的渗透率仍处于早期阶段，应用软件公司尚未找到让客户愿意为AI功能支付溢价的确定性路径。

因此，数据基础设施板块的反弹，本质上是市场在AI投资热潮中寻找“确定性”的结果。无论哪个应用层最终胜出，数据都需要被存储、处理和分析

[... middle omitted ...]

门已经在组织层面部署了生成式AI。这个数字远超一般企业软件的AI采用率，背后有两个不可逆的结构性驱动力：一是合格会计人才的持续短缺，二是全球税务合规复杂性的不断上升。

然而，摩根斯坦利的核心判断并非“AI将颠覆会计软件市场”，而是“AI将强化现有寡头的竞争壁垒”。报告明确指出，尽管AI原生初创公司正在攻击单个工作流程环节，但核心税务引擎市场仍将由Wolters Kluwer和Thomson Reuters等现有巨头主导。原因在于：税务合规的核心不是算法精度，而是对法规的理解、对合规系统的信任、以及与客户工作流程的深度嵌入。这些要素构成了数据、信任和习惯的三重壁垒，AI初创公司很难在短期内跨越。

这一点对于理解整个企业软件市场的竞争格局具有普遍意义。AI技术本身是开源且可复制的，但围绕AI构建的商业系统——包括数据积累、合规认证、客户转换成本——才是真正的护城河。摩根斯坦利的分析暗示，市场可能低估了现有软件巨头在AI时代的防御能力，同时高估了AI初创公司的颠覆速度。

\subsection{[R039] Bernstein：印度市场真正低估的不是电动化，而是混动供给侧的再定价}
{\small\color{refgray}归类：消费与科技硬件：AI眼镜渠道扩容与印度混动拐点}

Bernstein：印度市场真正低估的不是电动化，而是混动供给侧的再定价

印度汽车市场正在上演一场被大多数投资者忽略的结构性变化。2025年，在没有任何中央补贴、仅有8款在售车型、且大型混动车面临40\%高税率的环境下，强混动销量依然实现了超过85\%的增长，渗透率从1.4\%跃升至2.4\%。同期，印度市场有35至40款纯电动车型在售，享受5\%的低税率，但增速反而落后。

这不是一个需求问题。这是一个供给问题。Bernstein这份研报的核心判断是：印度强混动市场的规模几乎完全由车型供给决定，而供给侧的拐点即将到来。当车型目录在未来十八个月内翻倍，市场将进入一个全新的定价阶段。

为什么这个判断现在值得认真对待？因为在全球范围内，混动正在重新获得关注。美国、欧洲和日本市场在补贴退坡后，混动增速已反超纯电。印度作为全球第三大汽车市场，其混动渗透率目前仅2.3\%，但车型数量只有成熟市场的十分之一。如果供给是关键变量，那么当前的渗透率数字可能严重低估了潜在需求。

报告的核心主张是：印度强混动市场将从当前约2.3\%的渗透率，在2030年前提升至7\%至8\%，驱动因素不是政策变化，而是车型供给的集中释放。这一判断的置信度被刻意锚定在本十年的后半段，因为纯电成本下降和税收壁垒是真实存在的风险。但正是这种克制，让报告的结论更加可信。

1. 83\%的增长来自8款车型，这不是巧合，是“被压抑的需求”

印度市场实际上已经无意中做了一次压力测试。在8款强混动车型面对35至40款纯电动车型、混动车享受零补贴且大型混动税率高达40\%的条件下，混动销量在2025年增长了83\%。Bernstein的数据显示，2025年印度强混动零售量达到10.7万辆，而2024年仅为5.76万辆。

这个数字需要放在正确的框架中理解。有人可能会说，这只是一个低基数效应，或者是丰田一家公司的表现。确实，丰田在印度强混动市场占据了超过80\%的份额，Innova Hycross和Hyryder两款车型就占了全部混动销量的四分之三左右。

但这里的关键在于：丰田的高份额恰恰证明了供给限制的存在，而不是需求有限。直到最近，丰田几乎是唯一一家在印度规模化销售混动车型的车企。当玛鲁蒂铃木开始在2024年下半年推出基于丰田平台的混动车型后，其混动销量也出现了明显增长。玛鲁蒂在Victoris车型上做了一个关键决策：将全景天窗、360度摄像头、ADAS、通风座椅等高端配置全部集中在混动顶配版本上，而燃油版无法选配。消费者用行动回应了这种产品策略。

Bernstein用一个跨市场对比图清晰地说明了这一点。日本的混动渗透率约34\%，对应35至40款车型；欧洲约35\%，对应50至60款车型；美国约12\%至13\%，对应35至40款车型。印度以2.3\%的渗透率和8款车型，孤零零地落在图表的左下角。这不是一个不同类型的市场，这是一个门最少的市场。

这个跨市场对比意味着什么？它意味着如果印度混动车型数量向其他市场靠拢，渗透率的提升可能比大多数预测模型更陡峭。当然，印度的人均收入、油价、充电基础设施与发达国家不同，但车型供给与渗透率之间的强相关性表明，供给瓶颈是当前最关键的约束，而非需求不足。

如果混动市场的规模由车型目录决定，那么未来十八个月将是印度混动市场的决定性时刻。Bernstein梳理了确认和接近确认的车型计划：玛鲁蒂铃木将在2027年推出自研的串联混动系统，首先搭载于Fronx，随后扩展到Swift和Baleno；现代汽车计划到2030年推出8款混动车型，其中Creta混动版本已经确认；起亚Seltos混动、雷诺Duster E-Tech、马恒达XUV 3XO混动也在规划中。

这份名单的关键在于价格带。当前印度混动车型主要集中在25万卢比以上的高端市场，而即将推出的新车

[... middle omitted ...]

需求将出现数量级的变化。

从竞争格局看，玛鲁蒂的布局最有结构性优势。其目标是到2030至2031年，实现约35\%的CNG/CBG、25\%的燃油、25\%的混动和15\%的纯电的能源组合。这个组合是刻意设计的：无论哪种动力总成最终胜出，玛鲁蒂都有足够的产能和产品线来应对。它不需要押注哪个技术路线会赢，它只需要确保无论谁赢，自己都在牌桌上。

印度对混动车的税收政策是一个需要正视的约束。目前，长度超过4米的大型混动车与大型燃油车一样，面临40\%的商品服务税；而纯电动车仅需缴纳5\%。这种税收差距意味着，即使混动车的燃油效率显著优于燃油车，其初期购买成本仍然高于同级别燃油车，更不用说与纯电动车相比了。

但有几个因素正在改变这个等式。首先，即将推出的多款混动车型长度小于4米，符合18\%的小型车税率，这将显著缩小混动车与纯电动车之间的税收差距。其次，即使维持40\%的税率，混动车在全生命周期成本上仍然具有竞争力，因为其燃油消耗比同级别燃油车低40\%至50\%，且残值表现优于纯电动车。

\subsection{[R040] GS：AI眼镜渠道扩容的真正赢家不是Best Buy，而是能完成处方闭环的零售商}
{\small\color{refgray}归类：消费与科技硬件：AI眼镜渠道扩容与印度混动拐点}

GS：AI眼镜渠道扩容的真正赢家不是Best Buy，而是能完成处方闭环的零售商

当Best Buy宣布与Meta合作，在50多家门店内开设“Meta Lab”体验空间时，市场的第一反应往往是：电子产品零售商要抢眼镜店的生意了。这个直觉并不错，但它忽略了一个更关键的结构性事实——AI眼镜的购买决策，正在从“电子消费品”向“视力健康产品”迁移。而在这个迁移过程中，真正构成竞争壁垒的，不是门店数量或展示面积，而是能否在同一个场景里完成验光、处方、镜片定制和框架适配的闭环。

GS在6月发布的一份专题研报中，系统评估了Best Buy与Meta合作对美国主要眼镜零售商的影响。结论清晰且反直觉：对于National Vision（EYE）和Warby Parker（WRBY）而言，这个合作并非威胁；对于EssilorLuxottica（ESLX）而言，反而是利好。真正需要关注的，是当渠道扩容加速了消费者教育之后，谁有能力将“尝鲜流量”转化为“处方收入”。

这份报告的价值不在于它判断了股价涨跌，而在于它提供了一个分析框架：在AI眼镜从极客玩具走向大众消费品的过程中，竞争壁垒的构成要素正在被重新定义。

GS在报告中披露了一个关键数据：超过50\%的Best Buy消费者希望在购买AI眼镜前能够亲眼看到实物。这个数据本身并不令人意外，真正重要的是它揭示的消费者决策链条——人们走进门店，不只是为了摸一摸产品，而是为了确认“这副眼镜戴在我脸上是否合适、是否可以配上我的处方”。

这正是眼镜零售商与电子产品零售商之间最本质的差异。National Vision在美国拥有超过1000家门店，每家门店都配备了验光师和现场镜片加工能力。当消费者走进America‘s Best或Eyeglass World，他们可以在一次到店中完成：验光、选择Ray-Ban Meta框架、定制处方镜片、当场取货。而Best Buy的消费者如果需要处方镜片，则必须将眼镜寄送到Meta指定的在线验光平台，再等待镜片制作完成后寄回。

这个流程差异直接决定了转化率和客单价。GS报告引用National Vision管理层的话指出，Ray-Ban Meta已经成为其门店中周转最快的SKU，且购买Meta眼镜的消费者贡献了最高的平均交易价值。这背后的逻辑很清楚：当验光、镜片和框架可以在一个场景中解决时，消费者的决策成本大幅降低，同时零售商获得了镜片这一高毛利附加品的收入。

对于Best Buy而言，它的竞争优势在于展示和体验，但一旦消费者进入“我需要配一副能日常佩戴的眼镜”这个决策阶段，Best Buy的体验优势就会转化为体验断层。GS的判断是：National Vision提供的无缝体验是Best Buy无法复制的。

GS用Placer.ai和DataWorks的数据做了一个有趣的地理分析：将三家眼镜零售商的门店位置与Best Buy门店进行叠加，计算彼此在1英里、3英里、5英里半径内的重叠率。结果显示，Warby Parker有54\%的门店位于Best Buy周边3英里范围内，National Vision为53\%，EssilorLuxottica为49\%。从物理距离上看，三者与Best Buy的重叠度大致相当。

但物理距离的重叠并不意味着客流竞争的同质化。GS的分析揭示了一个更微妙的格局：当消费者在Best Buy体验了AI眼镜后，如果需要完成处方配置，他们会去哪里？答案是，他们会去最近的一家能提供验光和镜片服务的眼镜店。而在这个维度上，National Vision和Warby Parker的“服务密度”远高于Best Buy。

Warby Parker目前拥有约337家自有门店，覆盖44个州的主要城市。更重要的是，Warby Parker的核

[... middle omitted ...]

。但GS也谨慎地指出，通过Best Buy销售的AI眼镜，其处方镜片附加率可能低于通过EssilorLuxottica自有零售网络（如Sunglass Hut和LensCrafters）销售的产品。这意味着一部分增量收入是以牺牲利润率换来的。

3. 搜索趋势揭示了一个被忽视的信号：Warby Parker正在获得消费者心智

GS在报告中引用了Google Trends的数据，追踪了过去18个月消费者对“Best Buy AI glasses”、“America’s Best Meta”和“Warby Parker AI glasses”的搜索热度。结果值得玩味。

Best Buy的搜索热度在2025年6月之后迅速攀升，并在12月达到峰值，随后开始回落。而Warby Parker的搜索热度在2026年5月出现了明显的逆势上升，达到190的近期高点（以峰值100为基准）。这个趋势的变化时间点，恰好与Warby Parker宣布其AI眼镜即将在秋季发布的时间窗口吻合。

\subsection{[R041] GS：中国钛白粉出口韧性正在重塑全球定价权，而非简单的产能外溢}
{\small\color{refgray}归类：能源与大宗：供需错位结构性固化，上游资本开支超级周期存疑}

GS：中国钛白粉出口韧性正在重塑全球定价权，而非简单的产能外溢

中国钛白粉（TiO2）的出口数据，正在讲述一个与市场主流叙事不同的故事。

过去两年，市场普遍认为中国钛白粉的全球扩张主要依赖低成本优势，且极易受到上游原料短缺（如硫磺问题）或贸易政策（如印度反倾销税）的扰动。但GS最新发布的这份研报显示，一个更深刻的结构性变化正在发生：中国钛白粉的出口供给弹性远超预期，即便在面临硫磺供应紧张的情况下，2026年前四个月的净出口量仍同比增加了8.8万吨，增幅达14\%。这一数字不仅否定了“供给端会自动调节”的简单假设，更直接冲击了西方钛白粉生产商（如科慕、特诺、Kronos）在核心市场的定价能力。

这份报告的核心判断是：市场低估了中国钛白粉出口的“韧性”，而这将导致2025年下半年西方生产商预期的涨价幅度低于预期。这不是一次性的数据波动，而是全球钛白粉定价权从供给侧结构性转移的信号。理解这一点，是判断未来12个月该行业资产价格走势的前提。

GS分析师在报告中明确指出，他们长期将中国的净出口量视为衡量西方钛白粉生产商市场健康状况的关键指标。这一指标的逻辑在于：中国出口越多，全球市场供给越宽松，西方生产商的涨价空间就越小。

然而，2026年以来的数据打破了两个市场共识。第一个共识是“硫磺问题会自然限制中国出口”。GS在报告中坦言，4月份出口数据的高企令他们感到意外，因为此前市场普遍预期上游硫磺的供应紧张会抑制中国工厂的开工率。但现实是，中国生产商的供应链韧性远超预期，他们通过调整原料结构或利用库存，成功维持了高水平的出口。第二个共识是“印度反倾销税的不确定性会抑制出口”。印度反倾销税目前仍处于暂停状态，且司法裁决时间表未定。由于中国在印度没有保税仓库，一旦关税恢复，影响将是即时的。但至少在目前，中国企业正在充分利用这个窗口期。4月份中国对印度出口环比增加了1.4万吨，达到3.2万吨，前四个月累计同比增加3.2万吨，增幅高达28\%。

这意味着什么？市场不能再简单地将中国出口视为一种“顺周期”的产能溢出。它正在变成一种“逆周期”的结构性力量。即便上游成本上升、贸易政策悬而未决，中国生产商依然有能力维持甚至扩大出口。这种韧性正在从根本上改变全球钛白粉市场的供需平衡方程式。

GS的报告提供了两个关键市场的详细数据，这些数据揭示了中国出口结构性扩张的具体路径。

首先是欧盟市场。4月份中国对欧盟出口环比增加0.4万吨至2万吨，前四个月累计同比增加1.9万吨，增幅达32\%。欧盟是西方高端钛白粉生产商的核心利润来源地，中国产品在此处的渗透率提升，直接侵蚀了当地生产商的定价权和市场份额。值得注意的是，这一增长发生在欧盟自身需求并未显著改善的背景之下，这意味着中国产品更多是在进行“替代性”扩张，而非“增量性”扩张。

其次是印度市场。印度是全球钛白粉需求增长最快的市场之一，也是西方生产商和中国生产商争夺的焦点。如前所述，中国对印度的出口增长极为迅猛。GS特别指出，由于中国在印度没有保税仓库，一旦反倾销税恢复，影响会立即显现。但恰恰是因为这种不确定性，中国企业目前正在以高于正常水平的速度向印度发货，试图在政策窗口关闭前锁定更多客户关系。这种“抢跑”行为，本质上是在用短期的出口量，换取长期的市场份额和客户粘性。

这两个市场的共同点是：中国产品正在从“低端补充”转向“核心替代”。西方生产商如果仅仅依赖品牌溢价或技术优势来维持价格，可能会发现这些护城河正在被中国产品的性价比和供给稳定性所侵蚀。

报告中最具洞察力的部分，是GS将中国出口数据与终端需求数据进行了交叉验证。

GS观察到，油漆产量（钛白粉最主要的下游应用）在年初的表现弱于预期。这是一个关键信号。当终端需求放缓时，市场通常会预期供给端会随之调整，从而支撑价格。但中国出口的

[... middle omitted ...]

市场空间，很可能被中国产品迅速填补。这会导致一个尴尬的局面：西方生产商承担了减产的财务损失（GS在报告中指出，Kronos在2025年第四季度因高固定成本吸收而出现亏损），却无法享受到预期的价格回升红利。

对于投资者而言，这意味着需要重新审视对钛白粉行业下半年盈利的预测。那些依赖“供给端自律带来价格回升”的盈利模型，可能面临被修正的风险。

4. 报告尚未完全回答的关键问题：中国企业的成本曲线究竟有多陡？

尽管这份报告提供了极具价值的数据和判断，但它也留下了一个GS自身没有完全回答的关键问题：中国钛白粉生产商的成本曲线究竟在什么位置？或者说，中国出口的“韧性”能持续多久？

报告提到了硫磺问题，但没有深入分析中国企业在应对原料成本上涨时的具体手段。它们是依靠更高效的工艺消化了成本，还是通过牺牲利润率来维持出口量？如果是前者，那么中国出口的结构性优势将更加稳固；如果是后者，那么当前的高出口量可能只是暂时的“价格战”，一旦价格触及中国企业的盈亏平衡点，供给自然会收缩。

\section{Disclaimer}
{\scriptsize\color{refgray}
This community is a paid learning and information-sharing community created for general educational purposes only. The membership fee is charged solely for access to community discussions, educational content, organization, and operational costs. It is not an advisory fee, management fee, performance fee, brokerage fee, or compensation for personalized financial, investment, legal, tax, or professional advice.\par
I participate and share information solely as an individual and for educational discussion among community members. I am not acting as a financial adviser, investment adviser, broker, dealer, portfolio manager, legal adviser, tax adviser, accountant, fiduciary, or any other licensed professional adviser. Nothing shared in this community constitutes, and should not be interpreted as, financial advice, investment advice, legal advice, tax advice, a recommendation, solicitation, offer, endorsement, instruction, or guarantee to buy, sell, hold, short, trade, or invest in any security, cryptocurrency, fund, derivative, strategy, or other financial product.\par
All content shared in this community, including messages, discussions, links, files, charts, examples, market commentary, watchlists, AI-generated or AI-polished summaries, and third-party materials, is general, impersonal, and educational in nature. It does not take into account any individual's financial situation, investment objectives, risk tolerance, investment horizon, liquidity needs, tax status, legal restrictions, or suitability requirements. No content should be relied upon as suitable for any specific person, account, portfolio, or financial decision.\par
Information shared in this community may be incomplete, inaccurate, outdated, speculative, AI-generated, AI-edited, or based on third-party internet sources that have not been independently verified. I make no representation or warranty as to the accuracy, completeness, timeliness, reliability, or suitability of any information shared.\par
Financial markets involve substantial risk, including the possible loss of principal. Past performance, historical data, hypothetical examples, simulated results, backtests, personal opinions, or educational examples do not guarantee future results. Each member is solely responsible for conducting their own research, exercising independent judgment, and making their own decisions. Before making any financial, investment, legal, or tax decision, members should consult qualified and licensed professionals.\par
Participation in this community, payment of any membership fee, or communication with me or other members does not create any adviser-client, fiduciary, professional, contractual, agency, partnership, employment, or other advisory relationship. By joining or participating in this community, you acknowledge and agree that you are solely responsible for your own decisions, actions, transactions, profits, losses, risks, and consequences, and that I shall not be liable for any direct, indirect, incidental, consequential, financial, legal, tax, or other loss or damage arising from or related to any content, discussion, information, or material shared in this community.\par
}
\end{document}
